\documentclass[11pt,oneside,openany]{book}

% BEGIN SOURCE preamble.tex
\usepackage[T1]{fontenc}
\usepackage{lmodern,amsmath,amssymb,amsthm,mathtools}
\usepackage[a4paper,margin=25mm,headheight=15pt]{geometry}
\usepackage{microtype,enumitem,xcolor,longtable,booktabs,array,needspace}
\usepackage[colorlinks=true,linkcolor=blue!40!black,urlcolor=blue!40!black]{hyperref}
\usepackage[nameinlink,noabbrev]{cleveref}
\usepackage{fancyhdr}
\hypersetup{pdftitle={Point-Set Topology: Unified Notes},pdfauthor={MathGenome preparation series}}
\pagestyle{fancy}\fancyhf{}
\fancyhead[L]{\small MathGenome: Point-Set Topology}
\fancyhead[R]{\small\thepage}
\renewcommand{\headrulewidth}{0.3pt}
\setlength{\parindent}{0pt}\setlength{\parskip}{5pt}
\setlength{\emergencystretch}{3em}\allowdisplaybreaks[2]\raggedbottom
\setcounter{tocdepth}{1}
\newtheorem{theorem}{Theorem}[chapter]
\newtheorem{lemma}[theorem]{Lemma}
\newtheorem{proposition}[theorem]{Proposition}
\newtheorem{corollary}[theorem]{Corollary}
\theoremstyle{definition}
\newtheorem{definition}[theorem]{Definition}
\newtheorem{example}[theorem]{Example}
\theoremstyle{remark}
\newtheorem{remark}[theorem]{Remark}
\newcounter{foundation}
\renewcommand{\thefoundation}{F\arabic{foundation}}
\crefname{foundation}{foundation}{foundations}
\Crefname{foundation}{Foundation}{Foundations}
\newcommand{\R}{\mathbb R}\newcommand{\C}{\mathbb C}
\newcommand{\N}{\mathbb N}\newcommand{\Z}{\mathbb Z}\newcommand{\Q}{\mathbb Q}
\newcommand{\eps}{\varepsilon}
\newcommand{\abs}[1]{\lvert#1\rvert}
\newcommand{\norm}[1]{\lVert#1\rVert}
\newcommand{\ip}[2]{\langle#1,#2\rangle}
\newcommand{\powset}{\mathcal P}
\newcommand{\NN}{\mathcal N}\newcommand{\UU}{\mathcal U}
\newcommand{\BB}{\mathcal B}\newcommand{\FF}{\mathcal F}
\DeclareMathOperator{\id}{id}
\DeclareMathOperator{\cl}{cl}
\DeclareMathOperator{\Int}{Int}
\DeclareMathOperator{\supp}{supp}
\DeclareMathOperator{\diam}{diam}
\DeclareMathOperator{\dist}{dist}
\DeclareMathOperator{\St}{St}
\DeclareMathOperator{\mesh}{mesh}
\newcommand{\family}[1]{}
\newcommand{\source}[1]{\par\begingroup\small\sffamily\color{black!65}Source comparison: #1\par\endgroup}

% END SOURCE preamble.tex

\begin{document}
\frontmatter

% BEGIN SOURCE frontmatter.tex
\hypersetup{pageanchor=false}
\begin{titlepage}
\centering
\vspace*{22mm}
{\Large MathGenome Mathematical Preparation Series\par}
\vspace{18mm}
{\Huge\bfseries Point-Set Topology\par}
\vspace{8mm}
{\LARGE Unified Undergraduate and Selected Bridge Notes\par}
\vspace{16mm}
{\large Original exposition, source comparison, and explicit proofs\par}
\vfill
{\large Direction D10\par}
{\large 2 November 2026\par}
\vspace{10mm}
{\small Mathematical source: Munkres, Willard, Kelley, and Steen--Seebach.\par}
{\small This manuscript is not itself a Lean verification certificate.\par}
\end{titlepage}
\hypersetup{pageanchor=true}
\chapter*{Purpose and scope}
\addcontentsline{toc}{chapter}{Purpose and scope}
These notes organize a common body of point-set topology without concatenating
four textbooks or repeating each author's notation. Their core comprises
topological constructions and maps, general convergence, compactness,
connectedness, separation, countability, complete metric spaces and metrization.
Selected bridge material adds local compactness and compactifications, uniform
structures, function spaces and Ascoli, Baire category, paracompactness and
partitions of unity, Stone--Weierstrass approximation and finite-dimensional
Brouwer fixed points. Counterexamples test the actual hypotheses.

The chapter sequence is a mathematical exposition, not a schedule of Lean
workers and not a prescribed number of proof goals. Definitions and results have
stable labels. A later formalization must preserve the stated object, hypotheses
and entire conclusion; it may use a different valid proof. An intermediate step
of a particular argument is not automatically a separate deliverable.

The source comparison chooses substantial complementary material, including
selected exercises. It does not claim to adopt every exercise, every one of the
143 Steen--Seebach examples, or all advanced results of Kelley and Willard.
The source register records the boundaries. Fundamental groups, covering-space
classification, surface classification and general dimension theory belong to
D23; deeper set-theoretic topology belongs to D21/D24. They are not silently
assumed to prove a result here.

The source comparison lines identify topics actually compared. Every proof is
new exposition; references do not stand in for a missing argument.

\section*{How to read a proof}
Read the exact space and separation assumptions before an implication.
\emph{Compact} does not include Hausdorffness. \emph{Regular} and \emph{normal}
do not include the $T_1$ axiom. A sequence test is used only under its stated
countability or metric hypothesis. A completion concerns a specified metric or
uniformity, not merely the original topology. A quotient construction includes
its topology and map property, not just an underlying bijection.

% END SOURCE frontmatter.tex

\tableofcontents

% BEGIN SOURCE foundations.tex
\chapter*{Foundations and conventions}
\addcontentsline{toc}{chapter}{Foundations and conventions}
\begin{description}[style=nextline,leftmargin=12mm]
\item[F1: sets and choice]\refstepcounter{foundation}\label{base:sets}
Classical logic, sets, functions, equivalence relations and set quotients,
finite and countable families, elementary cardinal arithmetic, well-ordering
and Zorn's lemma are the background. Choice is allowed, not claimed to be
eliminated. Uses such as ultrafilter extension, arbitrary products and
well-ordering an open cover are made explicit.
\item[F2: scalar mathematics]\refstepcounter{foundation}\label{base:real}
The ordered complete real field, the complex field, elementary scalar
sequences and series, interval limits, scalar continuous functions, elementary
trigonometry and one-variable differentiation and integration are background.
The construction of these fields is not repeated. General-space compactness,
Urysohn/Tietze, Ascoli, Baire and the general Stone--Weierstrass theorem are
not background assumptions. The polynomial approximation needed in
\cref{ch:d10-16} is constructed there.
\item[F3: finite algebra and Euclidean calculations]\refstepcounter{foundation}\label{base:linear}
Finite-dimensional linear algebra, inner products and Cauchy--Schwarz,
determinants, finite sums and elementary affine geometry are background.
Metric and topological consequences used here are proved or reduced to earlier
notes results. The Sperner counting argument and arbitrarily fine
triangulations needed for Brouwer are proved in \cref{ch:d10-17}.
\item[F4: ordinals in selected examples]\refstepcounter{foundation}\label{base:ordinals}
For the ordinal counterexamples only, use the standard ordinal construction,
the first uncountable ordinal $\omega_1$, transfinite induction and the fact
that the supremum of countably many countable ordinals is countable. The last
fact follows because a countable union of countable sets is countable under
choice. No set-theoretic topological conjecture is assumed.
\end{description}

These are declared relative foundations, not an assertion that every branch
of mathematics has been developed inside this volume. MathGenome's earlier
linear algebra and analysis components supply relevant prior mathematics;
whether a particular Lean declaration matches an eventual specification still
requires the normal formalization checks.

\section*{Standing notation}
A topology is specified on its underlying set. A subset inherits the subspace
topology unless another topology is named. Empty spaces and empty indexing
families are permitted; the empty product is a singleton. Compactness means the
finite-subcover property, without an implicit separation axiom. $T_3$ means
regular and $T_1$; $T_4$ means normal and $T_1$; Tychonoff means completely regular
and Hausdorff. The manuscript proves each use of these properties in its
chosen convention rather than transferring a textbook's convention by name.

A proper filter does not contain the empty set. Its convergence to $x$ means
it contains the neighborhood filter $\NN(x)$. A Moore--Smith subnet is indexed
by an order-preserving cofinal map; bridges to other conventions are explicit.
An entourage uniformity on a nonempty space is proper. On the empty space its
unique degenerate entourage family is handled separately, not used as a
proper convergence filter. A proper map has compact preimages of compact sets;
equivalence with closed maps having compact fibres is asserted only with the
hypotheses in \cref{ch:d10-11}.

The usual Euclidean topology, its subspace topologies, the product topology and
the box topology are distinguished. Supremum metrics and compact-open
topologies are not interchanged just because they have the same underlying
function set. Every assertion involving a series or a formally infinite
partition of unity includes the convergence or local-finiteness argument
that gives it meaning.

% END SOURCE foundations.tex

\mainmatter

% BEGIN SOURCE chapters/01.tex
\chapter{Topologies, bases, and closure}\label{ch:d10-01}
\source{MUN sections 12--14,17; WIL sections 3--5; STE examples 8--11,18--20.}

Topology begins by specifying which tests of proximity count as local tests. The tests need not come from a distance and need not distinguish all points. Throughout these notes a neighborhood of a point contains an open set containing that point; the neighborhood itself need not be open. This agrees with WIL's convention and differs from MUN's use of ``neighborhood'' for an open neighborhood.

\begin{definition}[Topology and its order]\label{d10:01:topology}
A topology on a set $X$ is a family $\tau\subseteq\powset(X)$ containing $\varnothing,X$, closed under arbitrary unions and finite intersections. A member is open; a set with open complement is closed. The empty union is $\varnothing$ and the empty intersection is $X$. A topology $\sigma$ on the same set is finer than $\tau$ if $\tau\subseteq\sigma$. The discrete topology is $\powset(X)$ and the indiscrete topology is $\{\varnothing,X\}$.
\end{definition}

\begin{proposition}[Closed sets and generated topologies]\label{d10:01:generated}
Closed sets contain $\varnothing,X$, are closed under arbitrary intersections and finite unions, and any family with these properties is the family of closed sets of a unique topology. Every family $\mathcal A\subseteq\powset(X)$ is contained in a smallest topology. Intersections of topologies on a fixed set are topologies; unions need not be.
\end{proposition}
\begin{proof}
Complementation interchanges unions and intersections, including the empty cases, so the closed-set assertions follow from De Morgan's laws, in both directions. The intersection of any family of topologies has the topology axioms because each required operation belongs to every topology in that family. Apply this to all topologies containing $\mathcal A$; that is a nonempty family since it contains the discrete topology. For failure of unions let $X=\{a,b,c\}$, $\tau_1=\{\varnothing,\{a\},X\}$, and $\tau_2=\{\varnothing,\{b\},X\}$. Their union omits $\{a,b\}$ although it contains both of its singleton summands.
\end{proof}

\begin{definition}[Order topology]\label{d10:01:order-topology}
On a linearly ordered set $L$, the order topology is generated by the open rays $\{x:x<a\}$ and $\{x:x>a\}$, $a\in L$. This definition also covers empty and singleton ordered sets. For a set with more than one point a basis consists of open intervals, together with intervals $[m,b)$ when a minimum $m$ exists and $(a,M]$ when a maximum $M$ exists.
\end{definition}

\begin{proposition}[Order separation and convex subspaces]\label{d10:01:ordered-spaces}
Every linearly ordered space is Hausdorff. If $A\subseteq L$ is convex, meaning that it contains every point between any two of its points, its order topology for the restricted order equals its subspace topology from $L$, whose open sets are intersections of $A$ with open sets of $L$.
\end{proposition}
\begin{proof}
For $x<y$, if some $z$ lies strictly between them, the rays $(-\infty,z)$ and $(z,\infty)$ are disjoint neighborhoods of $x,y$. If no such $z$ exists, use $(-\infty,y)$ and $(x,\infty)$; their intersection would consist of points strictly between $x,y$ and is therefore empty. This proves Hausdorffness. An order ray in $A$ with endpoint $a\in A$ is the intersection of $A$ with the corresponding ray of $L$, so the order topology on $A$ is contained in its subspace topology. Conversely restrict any ambient subbasic ray with endpoint $b$. If $b\in A$, it restricts to a subbasic order ray. If $b\notin A$, convexity prevents $A$ from having a point on each side of $b$, so the restriction is either all of $A$ or empty. Every restricted ambient subbasic set is thus order-open in $A$, proving equality. The basis description in the definition follows by taking finite intersections of rays: the strongest finitely many lower and upper bounds determine an interval, with a one-sided interval at an endpoint; the whole space is covered by such intervals when it has at least two points.
\end{proof}
\begin{theorem}[Basis criterion and comparison]\label{d10:01:basis}
A family $\BB$ of subsets of $X$ is a basis for a topology precisely when it covers $X$ and, for $x\in B_1\cap B_2$ with $B_i\in\BB$, there is $B_3\in\BB$ with $x\in B_3\subseteq B_1\cap B_2$. The generated topology consists of all unions of basis members. If $\BB,\mathcal C$ are bases for $\tau,\sigma$, then $\tau\subseteq\sigma$ precisely when every $x\in B\in\BB$ admits $C\in\mathcal C$ with $x\in C\subseteq B$.
\end{theorem}
\begin{proof}
If $\BB$ is a basis for a topology, coverage follows from expressing $X$ as a union of basis members; the refinement condition follows from expressing $B_1\cap B_2$ in the same way. Conversely let $\tau$ consist of all unions of basis members. It contains $\varnothing$ and $X$ and is closed under arbitrary unions. If $U,V\in\tau$ and $x\in U\cap V$, choose basis members $B_1\subseteq U$, $B_2\subseteq V$ containing $x$, then refine their intersection. Thus $U\cap V$ is a union of basis members and belongs to $\tau$. Induction gives finite intersections. For comparison, the displayed refinement says exactly that each $B\in\BB$ is $\sigma$-open. This implies every union of such $B$ is $\sigma$-open. Conversely, if $B$ is $\sigma$-open, the basis property of $\mathcal C$ supplies the refinement at each of its points.
\end{proof}

\begin{proposition}[Subbases]\label{d10:01:subbasis}
For any family $\mathcal S\subseteq\powset(X)$, the finite intersections of members of $\mathcal S$, including the empty intersection $X$, form a basis. Its topology is the smallest topology containing $\mathcal S$. In particular, no covering assumption on a subbasis is necessary with this convention.
\end{proposition}
\begin{proof}
The basis contains $X$ and is closed under pairwise intersections, so \cref{d10:01:basis} applies. Every topology containing $\mathcal S$ contains its finite intersections and their unions, and the generated topology contains each member of $\mathcal S$.
\end{proof}

\begin{definition}[Neighborhoods and local bases]\label{d10:01:neighborhood}
A neighborhood of $x\in X$ is a set $N$ for which $x\in U\subseteq N$ for some open $U$. Write $\mathcal N(x)$ for all such $N$. A local base at $x$ is a subfamily $\mathcal B_x\subseteq\mathcal N(x)$ such that every $N\in\mathcal N(x)$ contains a member of $\mathcal B_x$.
\end{definition}

\begin{proposition}[Local tests]\label{d10:01:local-tests}
Every neighborhood contains $x$; supersets and finite intersections of neighborhoods are neighborhoods. A subset $G\subseteq X$ is open exactly when it is a neighborhood of every one of its points. If $\BB$ is a basis, its members containing $x$ form a local base at $x$.
\end{proposition}
\begin{proof}
Open witnesses for neighborhoods give the first assertions, with intersections of their witnesses for the intersection assertion. If $G$ is a neighborhood of each $x\in G$, take the union of all open subsets of $G$; it contains each such $x$ and therefore equals $G$. The reverse implication uses $G$ itself as a witness. Finally an open witness for $N\in\mathcal N(x)$ contains a basis member containing $x$.
\end{proof}

\begin{definition}[Closure, interior, boundary, and density]\label{d10:01:closure-def}
For $A\subseteq X$ put
\[
 \cl_X A=\bigcap\{F:A\subseteq F\subseteq X,\ F\text{ closed}\},\qquad
 \Int_X A=\bigcup\{U:U\subseteq A,\ U\text{ open}\}.
\]
The boundary is $\partial_X A=\cl_X A\cap\cl_X(X\setminus A)$; $A$ is dense in $X$ if $\cl_X A=X$. A point $x$ is an accumulation point of $A$ if every neighborhood of $x$ meets $A\setminus\{x\}$; write $A'$ for the set of these points.
\end{definition}

\begin{theorem}[Neighborhood tests]\label{d10:01:closure-test}
For $A\subseteq X$, $x\in\cl_X A$ if and only if every neighborhood of $x$ meets $A$, and $x\in\Int_X A$ if and only if $A$ is a neighborhood of $x$. Either test may be made using a local base. Moreover
\[
 \cl_X A=A\cup A',\qquad
 A\text{ dense}\ \Longleftrightarrow\ A\cap U\ne\varnothing
 \text{ for every nonempty open }U.
\]
\end{theorem}
\begin{proof}
If $x\notin\cl_X A$, the open complement of the closure is a neighborhood of $x$ missing $A$. Conversely an open neighborhood $U$ disjoint from $A$ has closed complement containing $A$, so $x\notin\cl_X A$. Any neighborhood contains such an open neighborhood; a local base refines any neighborhood. The assertion about interior follows directly from its definition as a union. A point of $\cl_X A$ outside $A$ has every neighborhood meeting $A\setminus\{x\}=A$, and points in $A$ belong to the closure. This proves the union formula and its reverse inclusion. The density assertion is the closure test applied at every point, together with the existence of a point in a nonempty open set.
\end{proof}

\begin{theorem}[Closure and interior calculus]\label{d10:01:closure-laws}
Closure is extensive, monotone, idempotent, preserves finite unions, and sends $\varnothing$ to $\varnothing$. A set is closed if and only if it equals its closure. Also
\[
 \Int_X A=X\setminus\cl_X(X\setminus A),\quad
 \partial_X A=\cl_X A\setminus\Int_X A,\quad
 \Int_X(A\cap B)=\Int_X A\cap\Int_X B.
\]
The boundary is closed; $A$ is both open and closed if and only if $\partial_X A=\varnothing$.
\end{theorem}
\begin{proof}
Closure is the smallest closed superset by its definition and the closed intersection axiom. This gives extensivity, monotonicity, idempotence, the empty case, and the fixed-set test. The closed set $\cl_X A\cup\cl_X B$ contains $A\cup B$, so it contains $\cl_X(A\cup B)$; monotonicity gives the opposite inclusion. A point is outside $\cl_X(X\setminus A)$ precisely when it has an open neighborhood contained in $A$, proving the complement formula and then the boundary formula. The intersection of the two interiors is an open subset of $A\cap B$; every open subset of $A\cap B$ lies in each interior. A boundary is an intersection of closed sets. If its value is empty, $\Int_X A=\cl_X A$ and both equal $A$, so $A$ is open and closed. The converse follows from the same formula.
\end{proof}

\begin{theorem}[Recovering a topology from closure]\label{d10:01:kuratowski}
Suppose $c:\powset(X)\to\powset(X)$ satisfies $A\subseteq c(A)$, $c(c(A))=c(A)$, $c(A\cup B)=c(A)\cup c(B)$, and $c(\varnothing)=\varnothing$. There is a unique topology whose closure operator is $c$.
\end{theorem}
\begin{proof}
The finite union law implies monotonicity: if $A\subseteq B$, then $c(B)=c(A)\cup c(B)$. Let the proposed closed sets be the fixed sets of $c$. They include $\varnothing,X$ and are closed under finite unions. If $F_i$ are fixed and $F=\bigcap_i F_i$, monotonicity gives $c(F)\subseteq c(F_i)=F_i$ for each $i$, hence $c(F)\subseteq F$; extensivity gives equality. They therefore define a topology by \cref{d10:01:generated}. For every $A$, $c(A)$ is a fixed superset of $A$, and every fixed superset $F$ satisfies $c(A)\subseteq c(F)=F$. Thus $c(A)$ is exactly the closure. Uniqueness follows because closure determines the closed sets.
\end{proof}

\begin{example}[Useful elementary topologies]\label{d10:01:examples}
On $\R$, intervals $(a,b)$ form a basis for the usual topology, while intervals $[a,b)$ form a basis for the lower-limit topology: at a point of an intersection choose a sufficiently small interval of the same form. The latter topology is strictly finer, because ordinary intervals are unions of lower-limit intervals, whereas $[0,1)$ is not ordinarily open. On a set $X$, the cofinite topology has as opens $\varnothing$ and complements of finite sets. The cocountable topology uses countable complements instead. Both satisfy the union axiom because the complement of a nonempty union is contained in any one of the finite, respectively countable, complements; finite intersections use finite unions of complements. On an infinite set a subset is cofinite-closed precisely when it is finite or all of $X$, so every infinite subset is dense. On an uncountable cocountable space every countable subset is closed, whereas every uncountable subset is dense. These conclusions follow either from the list of closed sets or the neighborhood test.
\end{example}

\begin{example}[Closure does not commute with every operation]\label{d10:01:closure-cautions}
In the usual real line, $A_n=\{1/n\}$ for $n\ge1$ have $\bigcup_n\cl A_n=\{1/n:n\ge1\}$ but $0\in\cl(\bigcup_n A_n)$. Also $A=(-1,0)$ and $B=(0,1)$ have empty intersection but $0\in\cl A\cap\cl B$. Thus the finite union law has no unrestricted union or intersection analogue. Nor can ``accumulation'' always be read as ``infinitely many points in each neighborhood'': in the indiscrete two-point space $\{a,b\}$, $a$ is an accumulation point of $\{b\}$. Under the $T_1$ condition that singletons are closed, the infinite-points reading is valid. Indeed, if a neighborhood met $A\setminus\{x\}$ in finitely many points, remove those closed singletons from an open neighborhood inside it to obtain a neighborhood missing $A\setminus\{x\}$. The reverse implication needs no separation assumption.
\end{example}

% END SOURCE chapters/01.tex


% BEGIN SOURCE chapters/02.tex
\chapter{Continuous maps, embeddings, and gluing}\label{ch:d10-02}
\source{MUN section 18, especially exercises 2--4,9,12; WIL section 7 and section 8, items 8.5--8.6.}

A map is continuous when every local test in the target pulls back to a local test in the domain. Open and closed maps test direct images instead. Distinguishing these directions prevents several common mistakes about inverses and quotient constructions.

\begin{definition}[Continuity and map types]\label{d10:02:maps}
A map $f:X\to Y$ is continuous if $f^{-1}(V)$ is open in $X$ for every open $V\subseteq Y$. It is continuous at $x$ if $f^{-1}(N)\in\mathcal N(x)$ for every $N\in\mathcal N(f(x))$. It is open, respectively closed, if it sends every open, respectively closed, subset of $X$ to an open, respectively closed, subset of $Y$. A homeomorphism is a bijection continuous in both directions. An embedding is a map inducing a homeomorphism $X\to f(X)$, where the image has the subspace topology $\{f(X)\cap V:V\text{ open in }Y\}$.
\end{definition}

\begin{theorem}[Equivalent continuity tests]\label{d10:02:continuity-tests}
For $f:X\to Y$ the following are equivalent: continuity; continuity at every point; inverse images of closed sets are closed; and
\[
 f(\cl_X A)\subseteq\cl_Y f(A)\quad(A\subseteq X).
\]
To test continuity it suffices to test inverse images of a subbasis, and hence of a basis, of $Y$.
\end{theorem}
\begin{proof}
If inverse images of opens are open and $N$ is a neighborhood of $f(x)$, pull back an open witness inside $N$; this proves continuity at $x$. Conversely, for open $V$, pointwise continuity says that $f^{-1}(V)$ is a neighborhood of each of its points, hence is open by \cref{d10:01:local-tests}. Complementation proves equivalence with the closed-set test. If $f$ is continuous, the closed set $f^{-1}(\cl_Y f(A))$ contains $A$, hence contains $\cl_X A$; this proves the displayed inclusion. Conversely, apply the inclusion to $A=f^{-1}(F)$ for a closed $F\subseteq Y$. It gives $f(\cl_X A)\subseteq\cl_Y f(A)\subseteq F$, so $\cl_X A\subseteq A$ and $A$ is closed. Finally inverse images commute with arbitrary unions and finite intersections. The topology generated by a subbasis consists of unions of its finite intersections, so the subbasis test is sufficient.
\end{proof}

\begin{proposition}[Composition, restrictions, and range]\label{d10:02:restriction}
Identity maps, constant maps, and composites of continuous maps are continuous. If $A\subseteq X$ and $f:X\to Y$ is continuous, $f|_A$ is continuous. If $f(X)\subseteq B\subseteq Y$, then $f:X\to B$ is continuous if and only if $f:X\to Y$ is continuous. If $A\subseteq B\subseteq X$, the subspace topology induced on $A$ through $B$ is the one induced directly through $X$.
\end{proposition}
\begin{proof}
The inverse image under an identity is the set itself, and that under a constant is either the whole domain or empty. For composition, $(g\circ f)^{-1}(U)=f^{-1}(g^{-1}(U))$. For restriction, $(f|_A)^{-1}(U)=A\cap f^{-1}(U)$, which is subspace-open. A subset of $B$ open in $B$ has the form $B\cap U$ with $U$ open in $Y$, whose inverse image is $f^{-1}(U)$ because the range lies in $B$. This proves one direction of the range assertion; the other follows because $f^{-1}(U)=f^{-1}(B\cap U)$. Finally $(U\cap B)\cap A=U\cap A$ proves the transitivity assertion.
\end{proof}

\begin{theorem}[Homeomorphism and embedding tests]\label{d10:02:embedding}
A continuous bijection is a homeomorphism if and only if it is open, equivalently if and only if it is closed. A map $e:X\to Y$ is an embedding if and only if it is injective, continuous, and every open subset of $X$ has the form $e^{-1}(V)$ for an open $V\subseteq Y$. Equivalently, the topology on $X$ is exactly $\{e^{-1}(V):V\text{ open in }Y\}$ and $e$ is injective. Composites of embeddings are embeddings.
\end{theorem}
\begin{proof}
For a bijection $f$, the inverse image of $U\subseteq X$ under $f^{-1}$ is $f(U)$; this proves the open criterion. The closed criterion follows by complements because $f(X\setminus A)=Y\setminus f(A)$ for a bijection. For an injective $e$, saying that its inverse from $e(X)$ is continuous says precisely that each open $U\subseteq X$ is $e^{-1}(e(X)\cap V)=e^{-1}(V)$ for some open $V\subseteq Y$. Continuity supplies the reverse containment between the two topologies. If $e:X\to Y$ and $j:Y\to Z$ satisfy this test, write an open $U\subseteq X$ as $e^{-1}(V)$ and then $V=j^{-1}(W)$. Thus $U=(j\circ e)^{-1}(W)$. Composition is injective and continuous, so the test proves the last assertion.
\end{proof}

\begin{proposition}[Topology comparison and characteristic maps]\label{d10:02:characteristic}
For topologies $\tau,\sigma$ on $X$, the identity $(X,\sigma)\to(X,\tau)$ is continuous precisely when $\tau\subseteq\sigma$. Let $S=\{0,1\}$ have topology $\{\varnothing,\{1\},S\}$. The characteristic function $\chi_A:X\to S$ is continuous precisely when $A$ is open. With discrete topology on $\{0,1\}$ it is continuous precisely when $A$ is both open and closed.
\end{proposition}
\begin{proof}
The inverse image under the identity of an open set is that same set, giving the first assertion. The only nonautomatic subbasic inverse image for $S$ is $\chi_A^{-1}(\{1\})=A$. For the discrete target one must also check $\chi_A^{-1}(\{0\})=X\setminus A$.
\end{proof}

\begin{theorem}[Pasting over an open cover or a finite closed cover]\label{d10:02:pasting}
Suppose $X=\bigcup_{i\in I}A_i$ and continuous maps $f_i:A_i\to Y$ agree on overlaps. They determine a unique map $f:X\to Y$. This map is continuous if all $A_i$ are open, or if $I$ is finite and all $A_i$ are closed.
\end{theorem}
\begin{proof}
Agreement makes $f(x)=f_i(x)$ for $x\in A_i$ well-defined and forces uniqueness. If the pieces are open, for open $V\subseteq Y$ each $f_i^{-1}(V)$ is open in the open subspace $A_i$ and therefore open in $X$. Their union is $f^{-1}(V)$. For the closed case, if $F$ is closed in $Y$, each $f_i^{-1}(F)$ is closed in the closed subspace $A_i$, hence closed in $X$. The finite union $f^{-1}(F)$ is closed. The claim about an open or closed subset of an open or closed subspace follows respectively from intersecting two open or two closed sets, using the complement description for subspace-closed sets.
\end{proof}

\begin{lemma}[Locally finite unions of closed sets]\label{d10:02:locally-finite-closed}
A family $(A_i)_{i\in I}$ is locally finite if each point has an open neighborhood meeting only finitely many indexed members. For a locally finite family of closed sets, $\bigcup_i A_i$ is closed. The same holds for any family $(F_i)$ with each $F_i$ closed and $F_i\subseteq A_i$.
\end{lemma}
\begin{proof}
The family $(F_i)$ is locally finite whenever $(A_i)$ is, so it suffices to prove the first assertion. If $x\notin\bigcup_i A_i$, choose an open $U$ containing $x$ meeting only $A_{i_1},\ldots,A_{i_n}$. Then $U\setminus(A_{i_1}\cup\cdots\cup A_{i_n})$ is open, contains $x$, and misses every $A_i$. The complement of the union is therefore open.
\end{proof}

\begin{theorem}[Locally finite closed pasting]\label{d10:02:locally-finite-pasting}
The pasting conclusion of \cref{d10:02:pasting} holds for a locally finite closed cover, with any index set.
\end{theorem}
\begin{proof}
For a closed $F\subseteq Y$, the sets $f_i^{-1}(F)$ are closed in $X$, contained in $A_i$, and their union is $f^{-1}(F)$. The preceding lemma makes that union closed. Apply the closed-set continuity test.
\end{proof}

\begin{example}[Limits of elementary map tests]\label{d10:02:cautions}
The identity from a two-point discrete space onto the same set with indiscrete topology is a continuous bijection whose inverse is not continuous. Thus the word ``continuous'' alone cannot replace ``homeomorphism.'' A constant map sends every accumulation point of a nonempty set to the unique image point, which is not an accumulation point of its singleton image; continuity preserves closure, not necessarily the derived-set operation.

A countable closed cover already shows why arbitrary closed pasting fails. Take $X=\{0\}\cup\{1/n:n\ge1\}$ with its usual subspace topology and singleton pieces. Set $f(0)=0$, $f(1/n)=1$. Every piece restriction is continuous, yet $f^{-1}((-1/2,1/2))=\{0\}$ is not open. The cover fails local finiteness at $0$.

Even separate continuity on a product does not imply continuity. Put $h(x,y)=xy/(x^2+y^2)$ away from $(0,0)$ and $h(0,0)=0$. Fixing either variable gives a continuous real function (the zero function when that variable is zero). Along $(t,t)$, $t\ne0$, its value is $1/2$, so every neighborhood of $(0,0)$ contains points sent outside $(-1/4,1/4)$. Hence it is not continuous at the origin.
\end{example}

\begin{example}[Open and closed are independent]\label{d10:02:open-closed-examples}
Projection $p:\R^2\to\R$ onto the first coordinate is open: the image of a basic rectangle $U\times V$ is $U$ if $V\ne\varnothing$, and is empty otherwise; take unions. It is not closed, because $\{(x,y):xy=1\}$ is closed by ordinary real continuity, while its image is $\R\setminus\{0\}$. The inclusion $[0,1]\hookrightarrow\R$ is closed since subspace-closed subsets of the closed set $[0,1]$ are closed in $\R$; it is not open since its image of the whole domain is not open. Neither openness nor closedness by itself implies continuity: any map from an indiscrete two-point space bijectively to a discrete two-point space is both open and closed (its only domain-open or domain-closed sets are empty and whole) but is not continuous.
\end{example}

% END SOURCE chapters/02.tex


% BEGIN SOURCE chapters/03.tex
\chapter{Subspaces, products, sums, and quotients}\label{ch:d10-03}
\source{MUN sections 15--16,19,22, including section 22 examples 1--2,7 and exercises 2--3; WIL sections 6,8--9, including 8.12 and 9.13--9.15.}

The topology of a constructed space is best specified together with the maps that characterize it. Initial constructions control maps into the new space; final constructions control maps out of it. Neither description should be confused with an assertion that every topological property is preserved.

\begin{theorem}[Subspace formulas and universal property]\label{d10:03:subspace}
For $A\subseteq X$, the sets $A\cap U$, $U$ open in $X$, form a topology on $A$. Its closed sets are $A\cap F$ with $F$ closed in $X$; for $B\subseteq A$,
\[
 \cl_A B=A\cap\cl_X B.
\]
The inclusion $i:A\hookrightarrow X$ is an embedding, and $g:Z\to A$ is continuous if and only if $i\circ g$ is continuous. If $\BB$ is a basis of $X$, $\{A\cap B:B\in\BB\}$ is a basis of $A$.
\end{theorem}
\begin{proof}
Intersection with $A$ commutes with unions and finite intersections and sends $X$ to $A$, proving the topology axioms. Complementation within $A$ gives the closed-set description. Neighborhoods of a point in $A$ can be refined by sets $A\cap U$ with $U$ an ambient open neighborhood. Thus a point of $A$ lies in $\cl_A B$ precisely when every ambient neighborhood meets $B$, proving the closure formula. The inclusion and mapping assertion follow by pulling back $A\cap U$: $g^{-1}(A\cap U)=(i\circ g)^{-1}(U)$. The topology is exactly the one pulled back by the injective inclusion, so \cref{d10:02:embedding} proves the embedding assertion. Finally intersect an expression of each ambient open $U$ as a union of basis members with $A$.
\end{proof}

\begin{remark}[Interior depends on the ambient space]\label{d10:03:interior-caution}
There is always $A\cap\Int_X B\subseteq\Int_A B$ for $B\subseteq A$, but equality can fail. For $A=B=\{0\}\subseteq\R$, the left side is empty and the right side is $A$. Similarly $\partial_A B$ need not equal $A\cap\partial_X B$: in this example they are empty and $A$, respectively. These computations use that $\{0\}$ is closed but not open in $\R$.
\end{remark}

\begin{theorem}[Initial topology]\label{d10:03:initial}
Given maps $f_i:X\to Y_i$ from a set to topological spaces, there is a coarsest topology on $X$ making them continuous. It is generated by the subbasis $f_i^{-1}(V)$, $V$ open in $Y_i$. In this topology $g:Z\to X$ is continuous if and only if every $f_i\circ g$ is continuous.
\end{theorem}
\begin{proof}
Any suitable topology must contain the displayed inverse images. Their generated topology makes all $f_i$ continuous, so it is coarsest. For the mapping property, necessity follows by composition. Sufficiency follows because $g^{-1}(f_i^{-1}(V))=(f_i\circ g)^{-1}(V)$ is open for every subbasic set, and the subbasis continuity test applies. When the family is empty the topology is indiscrete, and every map into it is continuous, as the assertion requires.
\end{proof}

\begin{definition}[Product and box topologies]\label{d10:03:product}
For a family $(X_i)_{i\in I}$ the product set $P=\prod_{i\in I}X_i$ consists of functions $x$ with $x(i)\in X_i$; write $x_i=x(i)$ and $\pi_i(x)=x_i$. The product topology is the initial topology of all $\pi_i$. Its basis consists of finite intersections $\bigcap_{i\in F}\pi_i^{-1}(U_i)$, where $F\subseteq I$ is finite and each $U_i$ is open. The box topology instead has all rectangles $\prod_i U_i$ as a basis. The empty product is a singleton. If any factor is empty, a product with that factor is empty.
\end{definition}

\begin{theorem}[Product universal property and projections]\label{d10:03:product-universal}
A map $g:Z\to\prod_i X_i$ is continuous if and only if every coordinate $\pi_i\circ g$ is continuous. The projections are continuous and open. If the product is nonempty, every projection is surjective. Products over a finite index set have the box topology as their product topology.
\end{theorem}
\begin{proof}
The first assertion is \cref{d10:03:initial}. A basic product-open set has open projection: if it is empty this is immediate; otherwise, fix a point in it. Its $j$th projection is its prescribed open $U_j$ if $j$ is restricted, and $X_j$ otherwise, since the fixed point's other coordinates may be retained while its $j$th coordinate varies in that set. Arbitrary opens are unions of basic opens, proving openness. A fixed point in a nonempty product also lets the $j$th coordinate be replaced by any element of $X_j$, proving surjectivity without another choice. For finitely many coordinates every box is a finite intersection of coordinate inverse images; the two bases agree.
\end{proof}

\begin{theorem}[Products and subspaces]\label{d10:03:product-subspaces}
If $A_i\subseteq X_i$, the product topology of $\prod_i A_i$ equals its subspace topology in $\prod_i X_i$. If $e_i:A_i\to X_i$ are embeddings, their product is an embedding. Coordinate permutations and regroupings of products are homeomorphisms.
\end{theorem}
\begin{proof}
An intersection of $\prod_i A_i$ with an ambient basic open is exactly a basic open formed by the finitely many relative opens $A_i\cap U_i$. Conversely each relative open in a finite coordinate test has such a representation, so the bases agree. For the second assertion, first identify each $A_i$ homeomorphically with $e_i(A_i)$. The product of these homeomorphisms is bijective and continuous in both directions by \cref{d10:03:product-universal}, and the first assertion identifies its image topology. For a coordinate permutation or regrouping, both the stated bijection and its inverse have coordinate maps that are projections or tuples of projections. They are continuous by the same universal property.
\end{proof}

\begin{proposition}[Closure of rectangles]\label{d10:03:product-closure}
In classical set theory with choice, for $A_i\subseteq X_i$,
\[
 \cl_{\prod_iX_i}\!\left(\prod_i A_i\right)=\prod_i\cl_{X_i}A_i.
\]
In particular a product of dense subsets is dense. The product of closed subsets is closed, without needing choice for that assertion.
\end{proposition}
\begin{proof}
The right side is closed, since it is the intersection of the closed sets $\pi_i^{-1}(\cl A_i)$, and contains the rectangle. This gives one inclusion. If the right side is empty there is nothing more to prove. Otherwise take a point $x$ there and a basic neighborhood restricting a finite set $F$ of coordinates. For $i\in F$ choose $a_i\in A_i\cap U_i$ using the closure test. For the other coordinates choose $a_i\in A_i$; these sets are nonempty because $x_i\in\cl A_i$. This last simultaneous selection uses choice. The resulting point is in the rectangle and in the neighborhood, proving the other inclusion. If all $A_i$ are closed the rectangle equals $\bigcap_i\pi_i^{-1}(A_i)$, which directly proves the last assertion.
\end{proof}

\begin{example}[What the product property does not say]\label{d10:03:product-cautions}
The diagonal map $d:\R\to\R^{\N}$, $d(t)_n=t$ for $n\ge1$, is continuous for the product topology by its coordinate functions. It is not continuous for the box topology, since the inverse image of $\prod_{n\ge1}(-1/n,1/n)$ is $\{0\}$. Projection need not be closed, by the hyperbola in \cref{d10:02:open-closed-examples}. A product of infinitely many nontrivial discrete spaces need not be discrete: in $\{0,1\}^{\N}$ every basic neighborhood leaves some coordinate unrestricted, hence contains at least two points. The universal property concerns maps into products; separate continuity does not settle maps out of them, as \cref{d10:02:cautions} shows.
\end{example}

\begin{theorem}[Disjoint sums]\label{d10:03:sum}
Let $S=\coprod_{i\in I}X_i=\{(i,x):i\in I,\ x\in X_i\}$. Declare $U\subseteq S$ open if $\{x:(i,x)\in U\}$ is open in every $X_i$. This is a topology; each summand inclusion $j_i$ is an embedding with open and closed image. A map $f:S\to Y$ is continuous if and only if every $f\circ j_i$ is continuous. The sum over an empty family is empty.
\end{theorem}
\begin{proof}
The topology axioms hold separately on each summand. The induced topology on a summand is its original topology because any open subset there, empty on all other summands, is open in $S$. A whole summand and its complement are open, proving the image assertion. Finally the intersection of $f^{-1}(V)$ with summand $i$ corresponds exactly to $(f\circ j_i)^{-1}(V)$, so the defining open-set condition proves both directions of the continuity assertion.
\end{proof}

\begin{theorem}[Final topology]\label{d10:03:final}
For maps $q_i:X_i\to Y$ to a set $Y$, the family
\[
 \{V\subseteq Y:q_i^{-1}(V)\text{ is open in }X_i\text{ for every }i\}
\]
is the finest topology making every $q_i$ continuous. For this topology $h:Y\to Z$ is continuous if and only if every $h\circ q_i$ is continuous.
\end{theorem}
\begin{proof}
Inverse images commute with unions and finite intersections, so the displayed family is a topology. Every topology making the maps continuous is contained in it. For an open $W\subseteq Z$, the set $h^{-1}(W)$ is open in this final topology exactly when each $q_i^{-1}(h^{-1}(W))$ is open, which proves the mapping criterion. If no maps are given the topology is discrete; the condition and the mapping assertion still hold.
\end{proof}

\begin{definition}[Quotients and saturation]\label{d10:03:quotient}
A quotient map is a surjection $q:X\to Y$ such that $V\subseteq Y$ is open exactly when $q^{-1}(V)$ is open in $X$. A subset $A\subseteq X$ is saturated if $A=q^{-1}(q(A))$, equivalently if it is a union of fibers. The quotient topology of a surjection is the final topology for that map. For an equivalence relation $\sim$ on $X$, the quotient $X/{\sim}$ is the set of equivalence classes with this topology for the class map.
\end{definition}

\begin{theorem}[Descent and the quotient universal property]\label{d10:03:quotient-universal}
If $q:X\to Y$ is quotient, $h:Y\to Z$ is continuous if and only if $h\circ q$ is continuous. In particular a continuous $f:X\to Z$ constant on fibers of $q$ descends uniquely to a continuous $\bar f:Y\to Z$. Conversely a continuous surjection with the stated mapping test for every target space is quotient. Any continuous open surjection or continuous closed surjection is quotient. Composites of quotient maps are quotient. A continuous surjection admitting a continuous section is quotient.
\end{theorem}
\begin{proof}
The mapping property follows from \cref{d10:03:final}, and fiber constancy gives the unique set map $\bar f$ with $\bar f q=f$. For the converse let $V\subseteq Y$ have open inverse image under $q$ and use its characteristic function $h:Y\to S$ to the Sierpi\'nski space of \cref{d10:02:characteristic}. The composite is continuous, hence $h$ is continuous and $V$ is open. If $q$ is open, an open $q^{-1}(V)$ has open image $V$; for a closed $q$, use closed inverse images and their images instead. For quotient $q:X\to Y$ and $r:Y\to Z$, a set $V\subseteq Z$ is open exactly when $r^{-1}(V)$ is open, exactly when $q^{-1}(r^{-1}(V))$ is open. Finally if $s:Y\to X$ is continuous and $qs=\id_Y$, then $h q$ continuous implies $h=(h q)s$ continuous; apply the converse mapping test.
\end{proof}

\begin{proposition}[Quotient restrictions and open products]\label{d10:03:quotient-restriction}
If $q:X\to Y$ is quotient and $A=q^{-1}(B)$ is an open or a closed saturated subspace, then $q|_A:A\to B$ is quotient. If $q$ is open (respectively closed), its restriction over any $B\subseteq Y$ is open (respectively closed), and hence quotient. A product of continuous open surjections is a continuous open surjection in the presence of choice, so is quotient.
\end{proposition}
\begin{proof}
Suppose $A$ is open. If $V\subseteq B$ has $(q|_A)^{-1}(V)$ open in $A$, it is open in $X$ and equals $q^{-1}(V)$; hence $V$ is open in $Y$ and therefore in $B$. The restriction is continuous, giving the other direction. For closed $A$ use closed sets and the analogous closed-set criterion. If $q$ is open, a relative open $A\cap U$ has image $B\cap q(U)$ because $A$ is the full inverse image of $B$; this is relatively open in $B$. Replace ``open'' by ``closed'' to prove the closed version. For a product of open surjections, a basic product rectangle has image the rectangle of its coordinate images with only finitely many restrictions. Equality of these images uses a simultaneous choice of preimages in the other coordinates, unless the relevant image is empty. It follows that the product map is open; choice also gives surjectivity. Its continuity follows from the product universal property.
\end{proof}

\begin{example}[A restriction can lose the quotient topology]\label{d10:03:restriction-failure}
Let $X=[0,1]\coprod[2,3]$ and $q:X\to[0,2]$ take $t$ to $t$ on the first interval and to $t-1$ on the second. This is continuous. Images of closed sets are unions of two sets closed in the closed intervals $[0,1]$ and $[1,2]$, hence closed in $[0,2]$. Thus $q$ is quotient. Restrict to $A=[0,1)\coprod[2,3]$, still mapping onto $[0,2]$. The inverse image of $[1,2]$ is the open summand $[2,3]$, whereas $[1,2]$ is not open in $[0,2]$. This restriction is not quotient; $A$ is not saturated since it keeps one preimage of $1$ and discards another.
\end{example}

\begin{example}[A product of quotient maps can fail]\label{d10:03:quotient-product-failure}
Let $q:\R\to Y$ collapse $\{1,2,\ldots\}$ to a point $b$ and leave all other points distinct, with the quotient topology on $Y$. The identity of $\Q$ is quotient, but $q\times\id_{\Q}$ is not. Here is a full open-set verification. Put $a_n=\sqrt2/n$ and
\[
 U_n=\{(x,r)\in\R\times\Q:|x-n|<1/4,\quad |r-a_n|>|x-n|\},\qquad U=\bigcup_{n\ge1}U_n.
\]
Each $U_n$ is open by real continuity, and $U$ contains $\{n\}\times\Q$ for every positive integer $n$, because $a_n$ is irrational. Hence $U$ is saturated for $q\times\id$; its only non-singleton fibers consist of all $(n,r)$ at a fixed rational $r$. Its image $V$ contains $(b,0)$. If $V$ were open, it would contain $W\times((-\delta,\delta)\cap\Q)$ for an open neighborhood $W$ of $b$ and $\delta>0$. Choose $n$ with $0<a_n<\delta$. The open $q^{-1}(W)$ contains an interval $(n-\eta,n+\eta)$. Choose $0<t<\min(\eta,1/4)$ and then a rational $r$ with $|r|<\delta$ and $|r-a_n|<t$. Such an $r$ exists by density of $\Q$ in $\R$. The point $(n+t,r)$ belongs to the inverse image of that rectangle but to no $U_m$: it fails the inequality for $m=n$, and its first coordinate is outside the other strips. This is a contradiction. Thus $U$ is open while $V$ is not, as required.
\end{example}

% END SOURCE chapters/03.tex


% BEGIN SOURCE chapters/04.tex
\chapter{Nets, proper filters, and general convergence}\label{ch:d10-04}
\source{WIL sections 10--12, especially 10A--10C,11B--11C,12D; MUN section 17 for closure and convergence comparisons.}

Sequences test only countably many stages. General neighborhoods need not admit a countable cofinal collection, so topology requires a broader language of convergence. Nets retain ordered stages; filters retain sets on which a process is eventually confined. We use nonempty directed index sets and proper filters, so no empty process converges vacuously to every point.

\begin{definition}[Nets and Moore--Smith subnets]\label{d10:04:nets}
A directed set is a nonempty set $D$ with a reflexive transitive relation $\le$ such that any two elements have a common upper bound. Antisymmetry is not required. A net in $X$ is a function $d\mapsto x_d$ on a directed set. It is eventually in $A$ if $x_d\in A$ for all $d\ge d_0$ for some $d_0$, and frequently in $A$ if for every $d_0$ some $d\ge d_0$ has $x_d\in A$. It converges to $x$ if eventually in every $N\in\mathcal N(x)$, and clusters at $x$ if frequently in every such $N$. A Moore--Smith subnet is $x_{\phi(e)}$ indexed by a directed set $E$, where $\phi:E\to D$ is order-preserving and has cofinal range: for every $d\in D$ some $e\in E$ has $d\le\phi(e)$.
\end{definition}

\begin{lemma}[Subnet calculus]\label{d10:04:subnet-calculus}
An order-preserving cofinal $\phi$ is eventually above each $d\in D$. A subnet of a convergent net converges to the same point, and a subnet of a subnet is a subnet. If a subnet clusters at $x$, the original net clusters at $x$.
\end{lemma}
\begin{proof}
Choose $e_0$ with $\phi(e_0)\ge d$; for $e\ge e_0$, monotonicity gives $\phi(e)\ge d$. This proves preservation of eventual membership and hence of convergence. Composition of order-preserving maps is order-preserving; if $\psi:F\to E$ is cofinal, choose $e$ with $\phi(e)\ge d$ and then $f$ with $\psi(f)\ge e$, which proves cofinality of $\phi\psi$. For the last assertion, given a neighborhood $N$ and $d_0$, take $e_0$ with $\phi(e_0)\ge d_0$. Clustering of the subnet supplies $e\ge e_0$ with $x_{\phi(e)}\in N$, and $\phi(e)\ge d_0$.
\end{proof}

\begin{theorem}[Cluster points and subnets]\label{d10:04:cluster-subnet}
A net clusters at $x$ if and only if it has a Moore--Smith subnet converging to $x$.
\end{theorem}
\begin{proof}
Suppose $(x_d)$ clusters at $x$. Put
\[
 E=\{(d,U):d\in D,\ U\text{ an open neighborhood of }x,\ x_d\in U\},
\]
with $(d,U)\le(d',U')$ when $d\le d'$ and $U'\subseteq U$. This set is nonempty. Given two elements, choose $d_0$ above their indices and use clustering to choose $d'\ge d_0$ with $x_{d'}\in U_1\cap U_2$; then $(d',U_1\cap U_2)$ is an upper bound. The projection $\phi(d,U)=d$ is order-preserving. It is cofinal since $U=X$ permits any $d$. For an open neighborhood $V$ of $x$, clustering supplies an index $d$ with $x_d\in V$. Every index above $(d,V)$ has its term in $V$. Thus the subnet converges. Conversely a convergent subnet clusters (every sufficiently late term lies in a given neighborhood, and directedness permits also passing any specified stage); \cref{d10:04:subnet-calculus} transfers clustering to the original net.
\end{proof}

\begin{proposition}[An alternative subnet convention]\label{d10:04:subnet-interface}
Some authors call $x_{\phi(e)}$ a subnet whenever $\phi:E\to D$ is eventually above each $d\in D$, without requiring monotonicity. Such a net has a common Moore--Smith refinement with the original net. Consequently its convergence to $x$ still implies that the original net clusters at $x$.
\end{proposition}
\begin{proof}
Let $y_e=x_{\phi(e)}$ and put $H=\{(d,e):d\le\phi(e)\}$. Order it by $(d,e)\preceq(d',e')$ if $d\le d'$, $e\le e'$, and $\phi(e)\le\phi(e')$. To find an upper bound of two elements, choose $a\in D$ above $d_1,d_2,\phi(e_1),\phi(e_2)$ and choose $e'$ above $e_1,e_2$ sufficiently late that $\phi(e')\ge a$. Then $(a,e')$ works. The maps $H\to E$, $(d,e)\mapsto e$, and $H\to D$, $(d,e)\mapsto\phi(e)$, are order-preserving and cofinal. Cofinality in $E$ uses $(\phi(e),e)$; cofinality in $D$ uses eventual cofinality of $\phi$. The same net $x_{\phi(e)}$ indexed by $H$ is therefore a Moore--Smith subnet of both $y$ and $x$. If $y$ converges, this refinement converges, and \cref{d10:04:cluster-subnet} applies. We will use the monotone convention throughout; this proposition supplies the interface when the other convention occurs.
\end{proof}

\begin{theorem}[Closure by nets]\label{d10:04:closure-net}
For $A\subseteq X$, $x\in\cl_X A$ if and only if a net with all its terms in $A$ converges to $x$. A subset is closed if and only if it contains every limit in $X$ of every net with terms in that subset.
\end{theorem}
\begin{proof}
If $x\in\cl A$, index by pairs $(U,a)$ with $U$ an open neighborhood of $x$ and $a\in U\cap A$, ordered by reverse inclusion of $U$, ignoring $a$. This is a nonempty directed preorder: intersections of two neighborhoods meet $A$. The term at $(U,a)$ is $a$. For any open neighborhood $V$, an index with $U=V$ exists and its entire subsequent tail lies in $V$, so this net converges to $x$. This construction makes no simultaneous selection of one point for every neighborhood. Conversely each neighborhood of a limit contains a tail, hence at least one term in $A$, so the closure test applies. The closed-set assertion follows from $A=\cl A$.
\end{proof}

\begin{theorem}[Continuity by nets]\label{d10:04:continuity-net}
A map $f:X\to Y$ is continuous at $x$ if and only if $x_d\to x$ implies $f(x_d)\to f(x)$ for every net in $X$. Thus continuity is equivalent to preservation of all net limits.
\end{theorem}
\begin{proof}
Under continuity at $x$, inverse images of neighborhoods of $f(x)$ are neighborhoods of $x$, and eventual membership pulls forward. If continuity fails, choose a neighborhood $V$ of $f(x)$ whose inverse image is not a neighborhood of $x$. Every open neighborhood $U$ of $x$ then contains a point $a$ with $f(a)\notin V$. The pairs $(U,a)$ with this property, ordered by reverse inclusion of $U$, form a directed set as in \cref{d10:04:closure-net}. The net of their points tends to $x$, while all image terms are outside $V$, so the image does not tend to $f(x)$.
\end{proof}

\begin{definition}[Proper filters]\label{d10:04:filters}
A proper filter $\FF$ on $X$ is a nonempty family of subsets, closed under supersets and finite intersections, with $\varnothing\notin\FF$. A filter base is a nonempty family $\mathcal B$ of nonempty subsets such that for $B_1,B_2\in\mathcal B$ some $B_3\in\mathcal B$ is contained in $B_1\cap B_2$. Its generated filter consists of the supersets of its members. A filter $\mathcal G$ is finer than $\FF$ when $\FF\subseteq\mathcal G$. It converges to $x$ when $\mathcal N(x)\subseteq\FF$, and clusters at $x$ when every $F\in\FF$ meets every neighborhood of $x$.
\end{definition}

\begin{proposition}[Basic filter facts]\label{d10:04:filter-facts}
A filter base generates a proper filter. A family of subsets generates a proper filter by taking finite intersections and supersets exactly when all its finite intersections, including the empty intersection $X$, are nonempty. In particular no proper filter exists on the empty set. Neighborhoods form a proper filter at each point. For a map $f:X\to Y$ its image filter is
\[
 f_*\FF=\{B\subseteq Y:f^{-1}(B)\in\FF\};
\]
it is the proper filter generated by the sets $f(F)$, $F\in\FF$.
\end{proposition}
\begin{proof}
The directed-refinement condition for a base gives closure under intersections for its supersets; no base member is empty, so its generated filter is proper. For a general family, finite intersections constitute such a base precisely under the stated nonemptiness condition. Necessity follows because a proper filter contains each finite intersection and $X$. This also proves the empty-set assertion. The neighborhood assertion is \cref{d10:01:local-tests}, together with the fact that every neighborhood contains its point. Inverse images show the formula for $f_*\FF$ is a proper filter. If $f^{-1}(B)\in\FF$, take $F=f^{-1}(B)$ to get $f(F)\subseteq B$; conversely $f(F)\subseteq B$ gives $F\subseteq f^{-1}(B)$, hence membership. This proves the base description.
\end{proof}

\begin{theorem}[Clustering and finer filters]\label{d10:04:filter-cluster}
A proper filter $\FF$ clusters at $x$ exactly when $x\in\bigcap_{F\in\FF}\cl F$, and exactly when some finer proper filter converges to $x$. Also $x\in\cl A$ exactly when $A$ belongs to some proper filter converging to $x$.
\end{theorem}
\begin{proof}
The first equivalence is the neighborhood closure test for each $F$. If $\FF$ clusters, the sets $F\cap N$, $F\in\FF$, $N\in\mathcal N(x)$, are nonempty and form a filter base. Their filter contains both $\FF$ and $\mathcal N(x)$ and thus converges. Conversely if $\mathcal G\supseteq\FF$ converges, $F,N\in\mathcal G$ imply $F\cap N\ne\varnothing$. For the last assertion, if $x\in\cl A$, the nonempty sets $A\cap N$ form a base whose filter contains $A$ and converges. If $A\in\FF\to x$, every neighborhood meets $A$ because both belong to the same proper filter.
\end{proof}

\begin{theorem}[Continuity by filters]\label{d10:04:continuity-filter}
A map is continuous at $x$ if and only if $f_*\FF\to f(x)$ for every proper filter $\FF\to x$. It suffices to test $\FF=\mathcal N(x)$.
\end{theorem}
\begin{proof}
Continuity at $x$ means that for each neighborhood $V$ of $f(x)$, $f^{-1}(V)\in\mathcal N(x)$. If $\FF\supseteq\mathcal N(x)$, this says precisely $V\in f_*\FF$, so that image filter converges. Conversely apply the proposed property to the neighborhood filter itself. The image filter's convergence says $f^{-1}(V)\in\mathcal N(x)$ for each $V$, exactly pointwise continuity.
\end{proof}

\begin{theorem}[The net--filter bridge]\label{d10:04:filter-net}
For a net $(x_d)$, the tail sets $T_d=\{x_e:e\ge d\}$ generate a proper filter $\FF_x$. It has exactly the same limits and cluster points as the net. Conversely, for a proper filter $\FF$, the pairs $(F,a)$ with $a\in F\in\FF$, ordered by reverse inclusion of $F$, define a net with value $a$ whose tail filter is exactly $\FF$. A Moore--Smith subnet has a tail filter finer than the original tail filter.
\end{theorem}
\begin{proof}
Each tail is nonempty; a common upper bound of two indices gives a tail contained in both. A neighborhood belongs to the tail filter exactly when it contains a tail, proving the limits assertion. It meets every tail exactly when the net is frequently in it, proving the cluster assertion. The canonical pair-index set for $\FF$ is nonempty and directed because $F_1\cap F_2$ is nonempty and belongs to $\FF$. The tail after $(F,a)$ has as its set of values exactly $F$: all subsequent values lie in a member contained in $F$, and each $b\in F$ occurs at $(F,b)\ge(F,a)$. Thus its tail filter equals $\FF$. Finally, for a subnet, every original tail contains a sufficiently late subnet tail by eventual cofinality; hence every member of the original tail filter belongs to the new one.
\end{proof}

\begin{lemma}[Refining a tail filter by a subnet]\label{d10:04:finer-subnet}
If a proper filter $\mathcal G$ contains the tail filter of $(x_d)$, there is a Moore--Smith subnet whose tail filter contains $\mathcal G$. In particular, if $\mathcal G\to x$, that subnet converges to $x$.
\end{lemma}
\begin{proof}
Use indices $(d,G,e)$ where $G\in\mathcal G$, $e\ge d$, and $x_e\in G$. Order them by increasing $d,e$ and decreasing $G$. Such indices exist: every $G$ meets every tail because both belong to $\mathcal G$. To dominate two indices, first take $a\in D$ above their $d$ and $e$ coordinates. The set $G_1\cap G_2$ meets the tail after $a$, so choose $e'\ge a$ with $x_{e'}\in G_1\cap G_2$; then $(a,G_1\cap G_2,e')$ is an upper bound. Projection onto $e$ is order-preserving and cofinal, and once the $G$ coordinate lies in a fixed $G_0\in\mathcal G$, every later term lies in $G_0$. Thus the subnet's tail filter contains $\mathcal G$.
\end{proof}

\begin{theorem}[Products and unique limits]\label{d10:04:product-convergence}
A net in a product converges to $x$ if and only if every coordinate net converges to $x_i$. A proper filter $\FF$ on that product converges to $x$ if and only if $(\pi_i)_*\FF\to x_i$ for every $i$. A space is Hausdorff if and only if each net has at most one limit, equivalently if and only if each proper filter has at most one limit. Here Hausdorff means that distinct points have disjoint open neighborhoods.
\end{theorem}
\begin{proof}
Continuity of projections proves necessity of both coordinate assertions. A basic neighborhood restricts finitely many coordinates. For nets choose a common upper bound of the finitely many stages after which these restrictions hold; for filters intersect their finitely many inverse images, each of which belongs to $\FF$. This proves sufficiency. In a Hausdorff space, two distinct limits would force eventual membership in disjoint neighborhoods, impossible for a net on a nonempty directed set; for a filter their empty intersection would belong to the filter. If distinct $x,y$ have no disjoint neighborhoods, the sets $U\cap V$ for $U\in\mathcal N(x)$, $V\in\mathcal N(y)$ are a nonempty filter base and their filter converges to both. Its canonical net from \cref{d10:04:filter-net} likewise converges to both, proving both converses.
\end{proof}

\begin{theorem}[When sequences suffice]\label{d10:04:first-countable-sequences}
If $X$ is first countable, meaning that every point has a countable local base, then $x\in\cl A$ if and only if some sequence in $A$ converges to $x$. In such a domain, preservation of sequence limits characterizes continuity into any target space. No countability hypothesis on the target is needed.
\end{theorem}
\begin{proof}
At $x$ replace a countable local base $(B_n)_{n\ge1}$ by the decreasing base $C_n=\bigcap_{k\le n}B_k$. If $x\in\cl A$, choose $a_n\in A\cap C_n$. This uses countable choice and gives $a_n\to x$. The converse follows from the net closure theorem. If $f$ fails continuity at $x$, choose a neighborhood $V$ of $f(x)$ whose inverse image is not a neighborhood of $x$, and choose $a_n\in C_n\setminus f^{-1}(V)$. The sequence tends to $x$ but its image never enters $V$. Continuous maps preserve sequence limits by \cref{d10:04:continuity-net}.
\end{proof}

\begin{example}[Sequence tests can miss the topology]\label{d10:04:sequence-failure}
Give an uncountable set $X$ the cocountable topology. A convergent sequence $x_n\to x$ is eventually equal to $x$: the countable set $\{x_n:x_n\ne x\}$ is closed and misses $x$, so its complement is a neighborhood that excludes every term unequal to $x$. Conversely eventual equality implies convergence in any topology. Consequently the identity from this space to the discrete topology preserves every convergent sequence, although it is not continuous because singletons are not cocountable-open. Every uncountable $A\subsetneq X$ is dense, yet no sequence in $A$ converges to a point outside $A$. Nets recover those closure points by \cref{d10:04:closure-net}.

Coordinate cluster points also require care. The sequence alternating $(0,1)$ and $(1,0)$ in $\R^2$ has $0$ as a cluster point in each coordinate, while $(0,0)$ is not a cluster point of the sequence: $(-1/2,1/2)^2$ contains no term. Coordinate convergence, unlike coordinate clustering, includes simultaneous eventual constraints.
\end{example}

% END SOURCE chapters/04.tex


% BEGIN SOURCE chapters/05.tex
\chapter{Compactness, ultrafilters, and product compactness}\label{ch:d10-05}
\source{MUN sections 26--28,37, including section 37 exercise 1; WIL sections 12,17, especially 17.3--17.9 and exercises 17F--17G; STE examples 4--5,18--20.}

Compactness converts infinitely many local tests into finitely many sufficient tests. It does not include a separation axiom. Empty spaces are compact. The finite intersection formulation, followed by the convergence formulation, will let us prove product compactness without supposing that sequences detect an arbitrary topology.

\begin{definition}[Compactness and the finite intersection property]\label{d10:05:compact}
A space $X$ is compact if every cover of $X$ by open sets has a finite subcover; an empty finite subcover is permitted. A family $\mathcal A$ of subsets has the finite intersection property (FIP) if every finite subfamily has nonempty intersection, including the empty subfamily whose intersection is $X$. A subset is compact when it is compact in its subspace topology.
\end{definition}

\begin{theorem}[Closed-set compactness criterion]\label{d10:05:fip}
A space is compact if and only if every family of closed subsets with the FIP has nonempty total intersection. Equivalently, for every family $\mathcal A$ with the FIP,
\[
 \bigcap_{A\in\mathcal A}\cl_X A\ne\varnothing.
\]
\end{theorem}
\begin{proof}
For closed $F_i$, an empty total intersection says that the open complements cover $X$. A finite subcover is equivalent, by De Morgan's law, to a finite subfamily with empty intersection. Thus compactness is exactly the first assertion. If $\mathcal A$ has the FIP then so do its closures: every finite intersection of closures contains the corresponding nonempty intersection of original sets. The first assertion therefore implies the second. Conversely apply the second to a family already closed. These arguments include $X=\varnothing$: no family has the FIP there under our empty-intersection convention, and every open cover has the empty subcover.
\end{proof}

\begin{theorem}[Subspaces and finite unions]\label{d10:05:compact-heredity}
A subset $K\subseteq X$ is compact if and only if every cover of it by ambient open sets has a finite subcover. Closed subsets of compact spaces are compact. A finite union of compact subsets is compact.
\end{theorem}
\begin{proof}
Ambient open covers restrict to subspace-open covers. Conversely, for a subspace-open cover, use all ambient open sets whose intersections with $K$ equal a member of that cover. They cover $K$, so finitely many suffice and their intersections give a finite subcover from the original family. This formulation avoids choosing ambient representatives simultaneously. If $F$ is closed in compact $X$, append $X\setminus F$ to an ambient open cover of $F$. A finite subcover of $X$ then covers $F$ after the appended member is discarded. Finally a cover of $K_1\cup\cdots\cup K_n$ restricts to a cover of each compact $K_i$; the union of their finite subcovers is finite and covers the union. The case $n=0$ is the empty compact space.
\end{proof}

\begin{theorem}[Continuous images]\label{d10:05:compact-image}
If $f:X\to Y$ is continuous and $X$ is compact, its image $f(X)$ is compact. In particular a continuous real-valued function on a nonempty compact space attains both its maximum and its minimum.
\end{theorem}
\begin{proof}
An ambient open cover $(V_i)$ of $f(X)$ pulls back to an open cover $(f^{-1}(V_i))$ of $X$. A finite subcover pulls forward to the required finite cover of the image. For the real-valued assertion, its compact image $K\subseteq\R$ is bounded because the sets $(-n,n)$, $n\ge1$, cover it and a finite subcover lies in one such interval. Compact subsets of the Hausdorff real line are closed, as proved in \cref{d10:05:compact-hausdorff} below. Let $s=\sup K$ by order completeness. Every interval about $s$ meets $K$: otherwise its left half would produce a smaller upper bound. Thus $s\in\cl K=K$. The infimum is handled in the same way. This last part depends only on the subsequent elementary separation lemma, not on product compactness.
\end{proof}

\begin{theorem}[Compact subsets of Hausdorff spaces]\label{d10:05:compact-hausdorff}
If $K$ is compact in a Hausdorff space $X$ and $x\notin K$, there are disjoint open sets containing $x$ and $K$. Hence $K$ is closed. Two disjoint compact subsets of a Hausdorff space have disjoint open neighborhoods.
\end{theorem}
\begin{proof}
For each $y\in K$, Hausdorffness supplies disjoint open neighborhoods $U$ of $x$ and $V$ of $y$. Use the family of all such $V$ to cover $K$ and choose a finite subcover $V_1,\ldots,V_n$ and their corresponding $U_i$. Then $\bigcap_iU_i$ and $\bigcup_iV_i$ are disjoint open neighborhoods of $x$ and $K$. For $K=\varnothing$, take $X$ and $\varnothing$. Thus every point outside $K$ has a neighborhood missing it, proving closedness. For disjoint compact $A,B$, first separate each $a\in A$ from $B$ by open $U_a,V_a$ using the first assertion. A finite collection of $U_a$ covers $A$; the union of those $U_a$ and the intersection of the corresponding $V_a$ separate $A$ and $B$. Empty cases are immediate by using an empty neighborhood on the empty side.
\end{proof}

\begin{corollary}[The compact-to-Hausdorff test]\label{d10:05:compact-bijection}
A continuous map from a compact space to a Hausdorff space is closed. If it is bijective, it is a homeomorphism; if it is injective, it is an embedding; if it is surjective, it is quotient.
\end{corollary}
\begin{proof}
A closed subset of the domain is compact by \cref{d10:05:compact-heredity}; its image is compact by \cref{d10:05:compact-image}, hence closed by \cref{d10:05:compact-hausdorff}. The bijective assertion follows from the closed-map criterion in \cref{d10:02:embedding}. For an injection apply that assertion with the Hausdorff subspace image as target: subspaces are Hausdorff by intersecting separating neighborhoods with the subspace. The surjective assertion follows from \cref{d10:03:quotient-universal}.
\end{proof}

\begin{lemma}[Tube lemma and closed projections]\label{d10:05:tube}
Let $K$ be compact, $x\in X$, and $W\subseteq X\times K$ open with $\{x\}\times K\subseteq W$. There is an open neighborhood $U$ of $x$ such that $U\times K\subseteq W$. Consequently the projection $X\times K\to X$ is a closed map for every space $X$.
\end{lemma}
\begin{proof}
For each $k\in K$, choose a basic rectangle $U_k\times V_k$ containing $(x,k)$ and contained in $W$. By compactness finitely many $V_k$ cover $K$. The intersection $U$ of their $U_k$ is a neighborhood of $x$, and for $(u,k)\in U\times K$ choose one of those $V_k$ containing its second coordinate; the corresponding rectangle contains $(u,k)$. If $K$ is empty use $U=X$. For a closed $C\subseteq X\times K$, if $x$ is outside its projection, $W=(X\times K)\setminus C$ contains the entire fiber above $x$. The tube just found gives an open $U$ disjoint from the projection. Thus that projection is closed.
\end{proof}

\begin{theorem}[Finite products]\label{d10:05:finite-product}
A finite product of compact spaces is compact.
\end{theorem}
\begin{proof}
The empty product is a singleton and compact, since any cover contains a member covering its point. It suffices by induction to treat $X\times K$. If either factor is empty the assertion is immediate. Given an open cover, a slice $\{x\}\times K$ is homeomorphic to $K$ by the coordinate map and so is covered by finitely many cover members. Their open union $W_x$ contains the slice. The tube lemma supplies $U_x\times K\subseteq W_x$ with $U_x$ open about $x$. Compactness of $X$ gives finitely many $U_x$ covering $X$, and all the finitely many original cover members used for their $W_x$ cover the product. One can use the family of all possible finite choices and their tubes, so the argument needs only finite selections rather than choice over all $x$.
\end{proof}

\begin{definition}[Ultrafilters]\label{d10:05:ultrafilter}
An ultrafilter is a proper filter maximal among proper filters on its underlying set. The principal ultrafilter at $x$ is $\{A\subseteq X:x\in A\}$. A filter is free if the intersection of all its members is empty. An ultranet is a net eventually in $A$ or eventually in $X\setminus A$ for every $A\subseteq X$.
\end{definition}

\begin{theorem}[Ultrafilter alternative]\label{d10:05:ultrafilter-alternative}
A proper filter $\UU$ is an ultrafilter if and only if, for each $A\subseteq X$, exactly one of $A$ and $X\setminus A$ belongs to $\UU$. An ultrafilter converges to each of its cluster points. For any map $f$, $f_*\UU$ is an ultrafilter.
\end{theorem}
\begin{proof}
If $A\notin\UU$ and every $U\in\UU$ met $A$, the sets $U\cap A$ would form a proper filter base generating a strictly finer filter containing $A$, contradicting maximality. Thus some $U\in\UU$ misses $A$, giving $X\setminus A\in\UU$. Both cannot belong because their intersection is empty. Conversely if the alternative holds and $\mathcal G$ is a proper filter properly containing $\UU$, choose $A\in\mathcal G\setminus\UU$. Then $X\setminus A\in\UU\subseteq\mathcal G$, a contradiction. If $\UU$ clusters at $x$, a neighborhood $N$ absent from it would have its complement in $\UU$, contrary to clustering; hence every neighborhood belongs and $\UU\to x$. Finally the inverse images of $B$ and its complement are complementary. Apply the alternative to these sets and the definition of $f_*\UU$ to obtain the alternative for the image filter.
\end{proof}

\begin{theorem}[Ultrafilter extension; choice is used]\label{d10:05:ultrafilter-extension}
Every proper filter is contained in an ultrafilter. More generally any family with the FIP is contained in an ultrafilter. This proof uses Zorn's lemma.
\end{theorem}
\begin{proof}
For a proper filter $\FF$, order by inclusion all proper filters containing $\FF$. The collection is a set, is nonempty, and every nonempty chain has its union as an upper bound. Indeed two members of this union lie together in one filter of the chain, so their intersection belongs to the union; upward closure also holds and the empty set occurs in no chain member. The empty chain has $\FF$ as an upper bound. Zorn's lemma gives a maximal member, which is maximal among all proper filters because any larger proper filter would still contain $\FF$. A family with the FIP generates a proper filter by \cref{d10:04:filter-facts}; extend that filter. No proper filter exists on the empty underlying set, so no exceptional empty case is hidden in the statement.
\end{proof}

\begin{proposition}[Free ultrafilters and ultranets]\label{d10:05:ultranet}
On an infinite set the complements of finite subsets generate a proper filter, whose ultrafilter extensions are free. The tail filter of a net is an ultrafilter exactly when the net is an ultranet. Every net has a Moore--Smith subnet that is an ultranet, using ultrafilter extension.
\end{proposition}
\begin{proof}
Finite unions of finite subsets cannot exhaust an infinite set, so the complements have the FIP. Every point is excluded by one such complement, proving freeness of any extension. A set belongs to a tail filter precisely when the net is eventually in it, so the ultrafilter alternative is exactly the ultranet property. Extend the tail filter to an ultrafilter $\UU$ and apply \cref{d10:04:finer-subnet}; the resulting subnet has a proper tail filter containing $\UU$, hence equal to $\UU$ by maximality. It is therefore an ultranet.
\end{proof}

\begin{theorem}[Compactness and convergence]\label{d10:05:compact-convergence}
For a space $X$ the following are equivalent, working with choice:
\begin{enumerate}[label=(\roman*)]
\item $X$ is compact;
\item every proper filter on $X$ has a cluster point;
\item every net in $X$ has a cluster point;
\item every net in $X$ has a Moore--Smith convergent subnet;
\item every ultrafilter on $X$ converges to at least one point;
\item every ultranet in $X$ converges to at least one point.
\end{enumerate}
The common limit need not be unique. The ultrafilter-extension implication below is the choice-dependent step.
\end{theorem}
\begin{proof}
If $X$ is compact and $\FF$ proper, the family $\{\cl F:F\in\FF\}$ has the FIP because each finite intersection contains the corresponding nonempty filter intersection. The closed-set criterion gives a point in all these closures, a cluster point by \cref{d10:04:filter-cluster}. Conversely a closed family with the FIP generates a proper filter. A cluster point of that filter lies in the closure of each member of the family and thus in each member itself. This proves (i)$\Leftrightarrow$(ii). The net--filter bridge proves (ii)$\Leftrightarrow$(iii), and the cluster-subnet theorem proves (iii)$\Leftrightarrow$(iv). Statement (ii) implies (v) because an ultrafilter converges to each of its cluster points. If (v) holds, extend any proper $\FF$ to an ultrafilter $\UU$ using \cref{d10:05:ultrafilter-extension}. A limit of $\UU$ is a cluster point of $\FF$, since each neighborhood and each $F\in\FF$ belong to $\UU$ and therefore meet. This gives (ii). Finally ultranets have ultrafilter tail filters, and the canonical net of any ultrafilter has exactly that tail filter. The bridge and \cref{d10:05:ultranet} prove (v)$\Leftrightarrow$(vi). If $X$ is empty there are no proper filters and no nets with our index convention, so all the convergence assertions hold vacuously, in agreement with compactness.
\end{proof}

\begin{theorem}[Tychonoff product theorem]\label{d10:05:tychonoff}
An arbitrary product of compact spaces is compact in the product topology. Conversely, if a product is nonempty and compact, every factor is compact. These statements do not require Hausdorffness. The forward proof uses choice.
\end{theorem}
\begin{proof}
If the product $P$ is empty, it is compact. Otherwise take any ultrafilter $\UU$ on $P$. For every $i$, its image $(\pi_i)_*\UU$ is an ultrafilter and converges to some point of compact $X_i$ by \cref{d10:05:compact-convergence}. Choose one such limit $x_i$ for every $i$; this simultaneous selection is an explicit use of choice, since limits may be nonunique. Then $x=(x_i)\in P$, and \cref{d10:04:product-convergence} says $\UU\to x$. Every ultrafilter converges, so the compactness criterion proves $P$ compact. That criterion itself used ultrafilter extension, another explicit choice principle. If the index set is empty, the singleton product is compact directly, also consistent with this proof. For the converse, nonemptiness makes each projection surjective by \cref{d10:03:product-universal}. Compactness of the factor follows from the continuous-image theorem. Without nonemptiness the converse fails, for $\varnothing\times\R$ is compact although $\R$ is not.
\end{proof}

\begin{proposition}[Why full arbitrary-space product compactness has choice strength]\label{d10:05:choice-strength}
In set theory without assuming choice, the assertion that every product of compact spaces is compact implies the axiom of choice. Thus full arbitrary-space Tychonoff is stronger than simply ignoring the nonuniqueness of coordinate limits.
\end{proposition}
\begin{proof}
Let $(A_i)_{i\in I}$ be a family of nonempty sets. Adjoin a distinguished new point $\infty_i$ to each, using tagged copies if necessary, and give $X_i=A_i\cup\{\infty_i\}$ the topology $\{\varnothing,\{\infty_i\},X_i\}$. Each $X_i$ is compact: since $A_i$ is nonempty, a cover must contain $X_i$ to cover even one point of $A_i$. The product is nonempty without choice, since the specified function $i\mapsto\infty_i$ belongs to it. Assume product compactness. Each $A_i$ is closed in $X_i$, so $C_i=\pi_i^{-1}(A_i)$ is closed in the product. The family $(C_i)$ has the FIP: for finitely many specified coordinates one can make finitely many choices in the corresponding nonempty $A_i$ and use $\infty_i$ elsewhere. Finite choice is proved by induction in ordinary set theory and does not require the axiom of choice. Compactness gives a point in $\bigcap_i C_i$, precisely a choice function for $(A_i)$. For an empty $I$, the empty function already makes the conclusion true.
\end{proof}

\begin{theorem}[Compact intervals]\label{d10:05:interval}
Every closed bounded real interval is compact in its usual topology.
\end{theorem}
\begin{proof}
An empty interval is compact. For $a\le b$ fix an open cover of $[a,b]$, and let $E$ be the set of $t\in[a,b]$ such that $[a,t]$ is covered by finitely many cover members. A member containing $a$ shows $a\in E$. Let $c=\sup E$. Choose a cover member $V$ containing $c$ and $\delta>0$ such that $(c-\delta,c+\delta)\cap[a,b]\subseteq V$. If $c>a$, the supremum property supplies $t\in E$ with $t>c-\delta$; if $c=a$, take $t=a$. A finite cover of $[a,t]$ together with $V$ covers $[a,c]$, so $c\in E$. If $c<b$, choose $s$ with $c<s<\min(b,c+\delta)$. The same finite cover together with $V$ covers $[a,s]$, contradicting that $c$ is an upper bound of $E$. Hence $c=b$ and $b\in E$, proving a finite subcover. For $a=b$, the member containing $a$ already suffices.
\end{proof}

\begin{definition}[Countable and sequential compactness]\label{d10:05:compact-variants}
A space is countably compact if every countable open cover has a finite subcover; it is sequentially compact if every sequence has a convergent subsequence indexed by a strictly increasing sequence of natural numbers. It is limit point compact if every infinite subset has an accumulation point in the sense of \cref{d10:01:closure-def}.
\end{definition}

\begin{theorem}[The countable convergence criterion]\label{d10:05:countable-compactness}
Countable compactness is equivalent to the countable closed-family FIP criterion and to the assertion that every sequence has a cluster point. Sequential compactness implies countable compactness. In a first-countable space these two properties coincide. Compactness implies countable compactness and limit point compactness; for $T_1$ spaces, countable compactness and limit point compactness coincide.
\end{theorem}
\begin{proof}
The countable closed-family criterion is obtained by taking complements of a countable open cover, exactly as in \cref{d10:05:fip}. For a sequence $(x_n)_{n\ge1}$ the closed tail sets $F_m=\cl\{x_n:n\ge m\}$ are nonempty, decreasing, and have the FIP. Countable compactness gives a point in their intersection, precisely a cluster point of the sequence. Conversely, if a countable open cover $U_1,U_2,\ldots$ has no finite subcover, choose $x_n\notin U_1\cup\cdots\cup U_n$. A cluster point $x$ would belong to some $U_k$ while all terms with $n\ge k$ are outside $U_k$, a contradiction. This choice of terms uses countable choice. A finite cover presents no issue, since it already is a finite subcover.

A convergent subsequence makes its limit a cluster point of the original sequence, so sequential compactness implies countable compactness. If the space is first countable and $x$ is a cluster point of $(x_n)$, take a decreasing countable local base $B_k$ at $x$. Frequency supplies successively $n_k>n_{k-1}$ with $x_{n_k}\in B_k$; that subsequence converges to $x$. This proves the first-countable assertion without assuming separation.

Compactness implies countable compactness by the definition of a cover. If an infinite subset $A$ has no accumulation point, then $\cl A=A$ by \cref{d10:01:closure-test}. For every $a\in A$ there is an open neighborhood meeting $A$ only at $a$. Together with $X\setminus A$ all such neighborhoods cover $X$, but finitely many meet $A$ in only finitely many points. This contradicts compactness and proves limit point compactness. More generally, countable compactness implies limit point compactness in the presence of choice: take distinct points $a_n$ in any infinite $A$ and take a cluster point $x$ of this sequence. Every neighborhood of $x$ meets infinitely many of the distinct $a_n$, so meets $A\setminus\{x\}$. Conversely assume $T_1$ and limit point compactness. For any sequence with finite range some value occurs infinitely often and is a cluster point. For infinite range, choose an accumulation point of that range. By the $T_1$ observation in \cref{d10:01:closure-cautions}, each of its neighborhoods contains infinitely many range points, hence terms at arbitrarily large indices. It is a sequence cluster point. The sequence criterion yields countable compactness.
\end{proof}

\begin{example}[Two strict distinctions]\label{d10:05:variant-counterexamples}
Compact need not imply sequentially compact. The space $\{0,1\}^{\powset(\N)}$ is compact by Tychonoff, with each factor discrete and finite. Define a sequence by $x_n(A)=1$ if $n\in A$, and $0$ otherwise. For any subsequence indexed by $n_1<n_2<\cdots$, let $A=\{n_2,n_4,\ldots\}$. Its $A$-coordinate alternates $0,1$, hence does not converge in the discrete two-point factor. Coordinate convergence shows the subsequence cannot converge. The choice of $A$ is made after the proposed subsequence is specified.

Without $T_1$, limit point compactness need not imply countable compactness. On $\N\times\{0,1\}$ take the product of the discrete topology on $\N$ and the indiscrete two-point topology. Every nonempty subset has an accumulation point: if it contains $(n,i)$, the other point $(n,1-i)$ has every neighborhood containing $(n,i)$. Yet the countable cover by the open pairs $\{n\}\times\{0,1\}$ has no finite subcover.
\end{example}

\begin{corollary}[Heine--Borel in finite-dimensional Euclidean space]\label{d10:05:heine-borel}
A subset of $\R^n$, with its usual topology, is compact if and only if it is closed and bounded. This includes $n=0$ and the empty subset.
\end{corollary}
\begin{proof}
For $n\ge1$, a compact subset is closed because $\R^n$ is Hausdorff: distinct points differ in a coordinate, and disjoint real intervals in that coordinate pull back to separating neighborhoods. The open cubes $(-m,m)^n$, $m\ge1$, cover $\R^n$, so a compact subset lies in one such cube and is bounded. Conversely a bounded subset lies in $[-m,m]^n$ for some $m$. That cube is compact by \cref{d10:05:interval,d10:05:finite-product}. A closed subset of it is compact. The product topology here is the Euclidean topology: an open Euclidean ball about $x$ contains a cube about $x$ with side radius $\eps/\sqrt n$, and a sufficiently small Euclidean ball lies in any basic rectangle containing $x$. For $n=0$, the space is a singleton and each of its subsets is closed, bounded, and compact.
\end{proof}

\begin{example}[Compactness without hidden separation or sequence assumptions]\label{d10:05:cautions}
Every finite space is compact: choose one cover member for each of its finitely many points. Every indiscrete space is compact, since any cover of a nonempty one must contain the whole space; an infinite discrete space is not, since its singleton cover has no finite subcover. Every cofinite space and every one of its subspaces is compact: one nonempty cover member leaves only finitely many points uncovered, which can be covered by finitely many additional members. Consequently on an infinite cofinite space an infinite proper subset is compact and not closed. This shows precisely why Hausdorffness occurred in \cref{d10:05:compact-hausdorff}. In a two-point indiscrete space every proper filter converges to both points; the limit in the compactness criterion cannot be declared unique. Compactness supplies convergent subnets of arbitrary nets; it does not in this generality assert convergent subsequences. Metric versions and their additional countability mechanism are treated later.
\end{example}

% END SOURCE chapters/05.tex


% BEGIN SOURCE chapters/06.tex
\chapter{Connectedness, components, and local connectedness}\label{ch:d10-06}
\source{MUN sections 23--25; WIL sections 26--29; STE examples 116--118.}

Connectedness concerns the impossibility of splitting a space into two open pieces. It requires no separation axiom. Its local version concerns the availability of small connected neighborhoods and is a different condition.

\section{Separations and their consequences}
\begin{definition}\label{d10:06:connected}
A separation of a space $X$ is a decomposition $X=U\cup V$ into disjoint nonempty open sets. The space is connected if it has no separation. A subset is connected when its subspace is connected. Thus the empty space and every one-point space are connected. Two subsets $A,B\subseteq X$ are separated when $A\cap\cl B=B\cap\cl A=\varnothing$.
\end{definition}

\begin{proposition}\label{d10:06:separation-tests}
The following are equivalent: $X$ is connected; its only clopen subsets are $\varnothing,X$; every continuous map from $X$ to the discrete space $\{0,1\}$ is constant (where the map with empty domain counts as constant). A decomposition $Y=A\cup B$ is a separation of the subspace $Y$ precisely when $A,B$ are nonempty separated subsets of $X$.
\end{proposition}
\begin{proof}
In a separation each piece is the complement of the other, so it is clopen; conversely a proper nonempty clopen set gives a separation. A map to $\{0,1\}$ is continuous precisely when its two fibers are open, and a nonconstant such map has two nonempty fibers. Finally, for disjoint $A,B$ with union $Y$, the equalities $Y\cap\cl_X A=A$ and $Y\cap\cl_X B=B$ say exactly that both pieces are closed in $Y$, equivalently open in $Y$.
\end{proof}

\begin{theorem}\label{d10:06:permanence}
Continuous images of connected spaces are connected. If $A$ is connected and $A\subseteq B\subseteq\cl A$, then $B$ is connected. If $C$ is nonempty and connected and each connected set $A_i$ meets $C$, then $C\cup\bigcup_iA_i$ is connected. In particular, a union of connected sets having a common point is connected.
\end{theorem}
\begin{proof}
A separation of an image pulls back to a separation of the domain. For the second assertion suppose $B=U\cup V$ is separated. Connectedness puts $A$ in one piece, say $U$; if $A$ is empty, so is $B$. A nonempty relatively open piece $V$ would meet $A$, since $A$ is dense in $B$, a contradiction. For the union assertion, in any separation of $C\cup\bigcup_iA_i$ the connected subset $C$ lies on one side. Each $A_i$, meeting it, must lie on that same side, leaving the other side empty. For a nonempty family with common point $c$, take $C=\{c\}$, already contained in the union; an empty union is connected by definition.
\end{proof}

\begin{theorem}\label{d10:06:product}
An arbitrary product of connected spaces is connected. A nonempty product is connected if and only if every factor is connected.
\end{theorem}
\begin{proof}
An empty product is a singleton, and a product with an empty factor is empty. Otherwise choose a point $a=(a_i)$ in the product. For two nonempty connected spaces $X,Y$, the sets $(X\times\{b\})\cup(\{x\}\times Y)$ are connected by the union result and their union, as $x$ varies, is $X\times Y$. Induction proves the finite case. In the arbitrary product let $S_F$ consist of points whose coordinates outside the finite set $F$ agree with $a$. It is homeomorphic to a finite product and is connected. The union $S=\bigcup_F S_F$ is connected because all $S_F$ contain $a$. Every nonempty basic open set meets $S$, so $S$ is dense; the preceding theorem completes the proof. Conversely, every coordinate projection of a nonempty product is surjective, and continuous images preserve connectedness. Choice is used to select $a$ if the factors are merely given as nonempty.
\end{proof}

\begin{theorem}\label{d10:06:intervals}
A subset $A\subseteq\R$ is connected if and only if it is order convex: whenever $a,b\in A$ and $a<t<b$, one has $t\in A$. Empty sets and singletons are included. Consequently a continuous real-valued function on a connected space takes every intermediate value between any two values it assumes.
\end{theorem}
\begin{proof}
If a point $t$ between $a,b\in A$ is missing, intersection with $(-\infty,t)$ and $(t,\infty)$ separates $A$. Conversely suppose $A$ is convex and $A=U\cup V$ is a separation. Choose $a\in U,b\in V$, relabeling the pieces if necessary so $a<b$. The interval $[a,b]$ lies in $A$. Set $c=\sup(U\cap[a,b])$. If $c\in V$, relative openness gives a left-hand interval just below $c$ missing $U$; this contradicts the supremum property, since $c>a$. If $c\in U$ and $c<b$, relative openness gives a point of $U\cap[a,b]$ above $c$, also a contradiction. The remaining possibility $c=b\in U$ is excluded by $b\in V$. Thus $A$ is connected. The last assertion follows by applying the image theorem to the function.
\end{proof}

\begin{proposition}\label{d10:06:linear-continuum}
Let $L$ be a linearly ordered set with at least two points, dense order, and the least-upper-bound property for every nonempty subset bounded above. Then every convex subspace of $L$, with the order topology inherited from $L$, is connected.
\end{proposition}
\begin{proof}
If a convex subspace had a separation, select points $a<b$ on opposite sides. Its restriction would separate $[a,b]$ into relatively open pieces $U,V$, with $a\in U,b\in V$. The set $U$ has a supremum $c\in[a,b]$. If $c\in V$, then $c>a$ and openness in the order topology supplies $d<c$ with $(d,c]\cap[a,b]\subseteq V$; choosing $d\geq a$ if necessary contradicts the defining property of the supremum. If $c\in U$, then $c<b$ and openness together with order density supplies a point of $U$ strictly between $c$ and $b$, again a contradiction. Empty and singleton convex subsets are already connected. The real-line theorem is the special case $L=\R$.
\end{proof}

\section{Components, paths, and local structure}
\begin{definition}\label{d10:06:components}
For $x\in X$, the component $C(x)$ is the union of all connected subsets containing $x$. A path from $x$ to $y$ is a continuous map $p:[0,1]\to X$ with endpoints $x,y$. A space is path connected if each pair of its points can be joined by a path, and $P(x)$ denotes the set of points reachable from $x$. Both connectedness conventions include the empty space.
\end{definition}

\begin{theorem}\label{d10:06:component-partition}
Components are the maximal nonempty connected subsets; they partition $X$ and are closed. Path components also partition $X$ and are maximal nonempty path connected subsets. Each path component is contained in a component. Path components need not be closed.
\end{theorem}
\begin{proof}
The union defining $C(x)$ is connected. If a connected set meets it, their union is connected, and hence lies in $C(x)$ by maximality. This proves that two intersecting components agree and that every nonempty connected subset lies in one component. Since $\cl C(x)$ is connected, maximality also proves closedness. Constant paths give reflexivity, reversing time gives symmetry, and concatenating paths on $[0,1/2]$ and $[1/2,1]$ gives transitivity; continuity of concatenation follows from finite closed pasting. Every two points in one class can therefore be joined within the class, since the image of each witnessing path remains in the class. A path image is connected by \cref{d10:06:intervals,d10:06:permanence}, proving containment in a component. The next example supplies failure of closedness.
\end{proof}

\begin{example}\label{d10:06:sine-curve}
Let $G=\{(t,\sin(1/t)):0<t\leq1\}$ and $S=G\cup(\{0\}\times[-1,1])$. Then $S=\cl_{\R^2}G$ is connected but has exactly two path components, $G$ and the vertical segment. In particular $G$ is a path component which is not closed.

Indeed $G$ is a continuous image of $(0,1]$ and hence connected and path connected. Its stated closure follows by choosing sequences of arguments tending to infinity with prescribed sine values. The vertical segment is path connected. If a path joined it to $G$, write its coordinates as $(u(t),v(t))$. A component $(a,b)$, or a corresponding interval ending at $1$, of the open set $\{t:u(t)>0\}$ can be chosen with $u(a)=0$ and $u(t)>0$ immediately to its right: take the component containing a time when the path is in $G$, after a time when it is on the segment. On $[a,t_0]$ the function $u$ takes every sufficiently small positive value. For any sequence $t_k\downarrow a$, its positivity at $t_k$ and the intermediate value theorem provide times $s_k,r_k\in(a,t_k)$ at which $u$ is respectively $1/(\pi/2+2\pi n_k)$ and $1/(3\pi/2+2\pi m_k)$, with sufficiently large integers. Hence $v(s_k)=1$ and $v(r_k)=-1$, contradicting continuity at $a$. This proves the assertion without using any unproved theorem about images of locally connected spaces.
\end{example}

\begin{definition}\label{d10:06:local-definitions}
A space is locally connected at $x$ if every neighborhood of $x$ contains a connected neighborhood of $x$; it is locally path connected at $x$ if ``connected'' is replaced by ``path connected.'' A neighborhood need not itself be open. The space has the respective local property if it has it at every point.
\end{definition}

\begin{theorem}\label{d10:06:local-characterization}
A space is locally connected if and only if the components of every open subset are open in $X$, equivalently if it has a basis of connected open subsets. The same assertions hold with path connectedness and path components. In a locally path connected space, components and path components agree. Both local properties pass to open subspaces.
\end{theorem}
\begin{proof}
Let $U$ be open and $C$ one of its components. For $x\in C$ choose a connected neighborhood $V$ with $x\in V\subseteq U$. It lies in $C$; because $V$ is a neighborhood, $x$ is an interior point of $C$. Thus $C$ is open. Conversely, the component of an open neighborhood containing $x$ is a connected open neighborhood, which gives a basis. The same argument works for paths. In the path version each path component of $X$ is open. Within a connected component $C$, a path component and the union of the others would therefore separate $C$ unless there is only one. Finally the basis characterization restricts to open subspaces.
\end{proof}

\begin{proposition}\label{d10:06:local-quotients}
A quotient image of a locally connected space is locally connected. A continuous image of a path connected space is path connected. Arbitrary products of path connected spaces are path connected.
\end{proposition}
\begin{proof}
Let $q:X\to Y$ be a quotient surjection and $U\subseteq Y$ open. For a component $C$ of $U$, every component of $q^{-1}(U)$ meeting $q^{-1}(C)$ maps into $C$. Consequently $q^{-1}(C)$ is a union of components of $q^{-1}(U)$, all open by the preceding theorem. The quotient test makes $C$ open. The image assertion follows by composing a path with the map. For the product assertion choose coordinate paths between two given product points; their product is continuous by the product universal property. Empty products and products with an empty factor satisfy the stated convention.
\end{proof}

\begin{proposition}\label{d10:06:local-products}
A nonempty product is locally connected if and only if each factor is locally connected and all but finitely many factors are connected. The same statement holds with path connectedness in both places. A compact locally connected space has only finitely many components.
\end{proposition}
\begin{proof}
Suppose the product is locally connected and fix a point $a$. Given a neighborhood $U_i$ of $a_i$ in a factor, take a connected neighborhood $C$ of $a$ inside $\pi_i^{-1}(U_i)$. Its projection is a connected neighborhood of $a_i$, because $C$ contains a basic product neighborhood whose projection is a neighborhood. Varying $a$ proves local connectedness of every factor. Taking just one connected neighborhood $C$ of $a$, choose a basic neighborhood $V\subseteq C$ with finite support $E$. For $i\notin E$, the projection $\pi_i(C)$ equals the whole factor, so that factor is connected. Conversely, enlarge the support of any prescribed basic neighborhood by the finitely many possibly disconnected factors. In these finitely many coordinates choose connected open neighborhoods inside the prescribed ones. Their product with the remaining, connected, whole factors is a connected open neighborhood by \cref{d10:06:product}. For the path version repeat the argument using continuous images and products of path connected spaces as proved above. Nonemptiness permits all coordinate values to occur in these projections. Finally the components of a locally connected space are open, and cover it. Compactness yields a finite subcover, which, since components are disjoint, lists all components. This includes the empty space with zero components.
\end{proof}

\begin{remark}\label{d10:06:local-limitations}
Local connectedness alone does not imply connectedness: a discrete space with two points is a counterexample. Connectedness alone does not imply local connectedness: in the space $S$ above, any connected subset meeting the vertical segment and also a point of $G$ has first-coordinate image containing an interval $[0,c]$. It therefore contains points of height $1$ and $-1$ with positive first coordinate arbitrarily near zero. A sufficiently small neighborhood of $(0,0)$ cannot contain such a connected neighborhood. Arbitrary subspaces need not inherit local connectedness, since $\Q\subseteq\R$ has singleton components and is not discrete. In fact every connected subset of $\Q$ has at most one point and is a singleton when nonempty, by cutting between two rationals at an irrational.
\end{remark}

\section{Compact connected refinements}
\begin{definition}\label{d10:06:quasicomponent}
The quasi-component $Q(x)$ is the intersection of all clopen subsets of $X$ containing $x$. A space is totally disconnected if every component is a singleton. In these notes ``zero-dimensional'' means that clopen sets form a basis; it does not include a separation axiom and does not refer to covering dimension.
\end{definition}

\begin{theorem}\label{d10:06:quasicomponent-compact}
In a compact Hausdorff space, $Q(x)=C(x)$ for every $x$. Moreover, whenever $O$ is open and $C(x)\subseteq O$, some clopen $H$ satisfies $C(x)\subseteq H\subseteq O$.
\end{theorem}
\begin{proof}
A connected set containing $x$ cannot cross a clopen set, so $C(x)\subseteq Q(x)$. The set $Q=Q(x)$ is closed and compact. If it were disconnected, write $Q=A\cup B$ with disjoint nonempty relatively clopen pieces. They are compact and therefore closed in $X$. Compact Hausdorff normality, proved independently in \cref{d10:07:compact-normal}, gives disjoint open sets $U,V$ containing $A,B$. For each $y\in X\setminus(U\cup V)$ some clopen neighborhood $H_y$ of $x$ omits $y$, by the definition of $Q$. Compactness gives finitely many whose intersection $H$ lies in $U\cup V$. If the complementary set is empty take $H=X$. Since $Q\subseteq H$, both $H\cap U$ and $H\cap V$ are nonempty. They are relatively clopen in the clopen set $H$, hence clopen in $X$. The one containing $x$ would have to contain all of $Q$, a contradiction. Thus $Q$ is connected and equals $C(x)$. Apply the same finite-intersection argument to the compact set $X\setminus O$: for each of its points choose a clopen neighborhood of $x$ omitting that point. Their finite intersection is the required $H$.
\end{proof}

\begin{corollary}\label{d10:06:zero-dimensional}
A compact Hausdorff space is totally disconnected if and only if it has a clopen basis. More generally a zero-dimensional $T_1$ space is Hausdorff and totally disconnected. Without $T_1$, a clopen basis need not imply total disconnectedness.
\end{corollary}
\begin{proof}
For the forward compact assertion use the last theorem with $C(x)=\{x\}$. If a space has a clopen basis and is $T_1$, then for $x\ne y$ a clopen set $H$ satisfies $x\in H\subseteq X\setminus\{y\}$. The disjoint open sets $H,X\setminus H$ separate the two points and prevent any connected subset from containing both. The two-point indiscrete space has the clopen basis $\{X\}$ and is connected, showing the necessity of a separation qualification.
\end{proof}

\begin{theorem}\label{d10:06:filtered-continua}
Let $K$ be compact Hausdorff and let $(K_i)_{i\in I}$ be a nonempty family of nonempty closed connected subsets, filtered downwards: for $i,j$ some $k$ satisfies $K_k\subseteq K_i\cap K_j$. Then $\bigcap_iK_i$ is nonempty and connected.
\end{theorem}
\begin{proof}
Filteredness gives the finite intersection property, so compactness gives a nonempty intersection $L$. If $L=A\cup B$ were separated, its relatively closed pieces would be disjoint closed subsets of $K$. Normality gives disjoint open sets $U,V$ containing them. Some $K_i$ must lie in $U\cup V$: otherwise the closed sets $K_i\setminus(U\cup V)$ would have the finite intersection property and yield a point of $L$ outside $U\cup V$. But that $K_i$ meets both $A$ and $B$, and hence both $U,V$, which separate it. This contradiction proves connectedness.
\end{proof}

\begin{corollary}\label{d10:06:irreducible-continuum}
Call a nonempty compact connected Hausdorff space a continuum. If $K$ is a continuum and $\varnothing\ne A\subseteq K$, there is a subcontinuum $L$ containing $A$ minimal under inclusion among such subcontinua. In particular any two points of a continuum belong to a subcontinuum irreducible between them.
\end{corollary}
\begin{proof}
Order the subcontinua containing $A$ by reverse inclusion. This collection is nonempty because it contains $K$. Every nonempty chain has as an upper bound its intersection, which is a continuum by the preceding theorem and still contains $A$. The empty chain has upper bound $K$. Zorn's lemma gives a maximal element for reverse inclusion. Nonemptiness of $A$ makes the intersection's nonemptiness automatic as well; no assertion about a unique minimal continuum is intended.
\end{proof}

\begin{remark}\label{d10:06:continuum-scope}
The compact refinements proved here are the preceding three results. The Cantor binary/ternary representation, classification of nonempty perfect totally disconnected compact metrizable spaces, and Cantor surjections onto nonempty compact metric spaces are proved separately in \cref{d10:18:cantor-model,d10:18:cantor-characterization,d10:18:cantor-surjection}. The square-filling existence theorem is then obtained in \cref{d10:18:peano}. WIL section 28's cut-point characterizations of arcs and circles, section 29's general inverse-limit representation, and section 31's full Hahn--Mazurkiewicz characterization are outside the adopted scope and are not used as foundations. The word ``continuum'' has competing metric and nonmetric conventions; our definition explicitly allows nonmetrizable compact Hausdorff spaces and excludes the empty space.
\end{remark}

% END SOURCE chapters/06.tex


% BEGIN SOURCE chapters/07.tex
\chapter{Separation and extension by continuous functions}\label{ch:d10-07}
\source{MUN sections 31--35; WIL sections 13--15, especially 15.6--15.8 and exercises 15B--15C.}

\section{Conventions and elementary separation}
\begin{definition}\label{d10:07:axioms}
A space is $T_0$ if distinct points are distinguished by an open set, $T_1$ if every singleton is closed, and Hausdorff (or $T_2$) if distinct points have disjoint open neighborhoods. It is regular if a point and a closed set not containing it have disjoint open neighborhoods. It is normal if two disjoint closed sets have disjoint open neighborhoods. Neither regular nor normal includes $T_1$. We write $T_3$ for regular plus $T_1$, and $T_4$ for normal plus $T_1$.

A space is completely regular if for every closed $F$ and $x\notin F$ there is continuous $f:X\to[0,1]$ with $f(x)=1$ and $f|_F=0$. It is Tychonoff if it is completely regular and Hausdorff. Empty closed sets are permitted in all these definitions. All the conditions hold for the empty space.
\end{definition}

\begin{proposition}\label{d10:07:shrinking}
A space is regular if and only if, for $x\in U$ with $U$ open, some open $V$ satisfies $x\in V\subseteq\cl V\subseteq U$. A space is normal if and only if, for closed $F\subseteq U$ with $U$ open, some open $V$ satisfies $F\subseteq V\subseteq\cl V\subseteq U$.
\end{proposition}
\begin{proof}
In the regular case separate $x$ from $X\setminus U$ by open $V,W$. Since $V\cap W=\varnothing$, one has $\cl V\subseteq X\setminus W\subseteq U$. Conversely $V$ and $X\setminus\cl V$ separate $x$ and $X\setminus U$. The same argument with $F$ instead of $x$ proves the normal version. It also works when either closed set is empty.
\end{proof}

\begin{proposition}\label{d10:07:hierarchy}
Hausdorff implies $T_1$, which implies $T_0$. Complete regularity implies regularity. A regular $T_0$ space is $T_3$, and every $T_3$ space is Hausdorff. Normality together with $T_1$ implies regularity, and hence $T_4$ implies $T_3$. Normality alone does not imply regularity.
\end{proposition}
\begin{proof}
In a Hausdorff space every point outside $\{x\}$ has a neighborhood omitting $x$, proving $T_1$. Complements of closed singletons distinguish points. A separating function for $x,F$ gives disjoint neighborhoods $\{f>2/3\}$ and $\{f<1/3\}$. If regularity and $T_0$ hold, for distinct points an open $U$ contains one, say $x$, and omits $y$. Choose $x\in V\subseteq\cl V\subseteq U$; then $X\setminus\cl V$ contains $y$ and omits $x$. Thus each point has a neighborhood omitting the other, giving $T_1$. In a regular $T_1$ space separate $x$ from the closed singleton $\{y\}$ to obtain Hausdorffness. In a normal $T_1$ space apply normality to $\{x\},F$. For the last assertion give $\R$ the topology consisting of upper rays $(a,\infty)$, together with $\varnothing,\R$. All nonempty closed sets intersect, so the space is normal. The point $1$ and the closed set $(-\infty,0]$ have no disjoint open neighborhoods, so it is not regular.
\end{proof}

\begin{theorem}\label{d10:07:heredity-products}
The properties $T_0,T_1$, Hausdorff, regular, and completely regular are hereditary to arbitrary subspaces. Each is preserved by arbitrary products. Normality is hereditary to closed subspaces and is preserved by closed continuous surjections. Consequently $T_3$ and Tychonoff are hereditary and productive, and $T_4$ is hereditary to closed subspaces and preserved by closed continuous surjections.
\end{theorem}
\begin{proof}
For the first three properties intersect witnessing open sets with the subspace. If $F$ is closed in $Y\subseteq X$, write $F=Y\cap C$ with $C$ closed in $X$. For $y\in Y\setminus F$ one has $y\notin C$; separate $y,C$ in $X$ by neighborhoods or a function and restrict to $Y$. This proves the remaining hereditary assertions.

For products the first three properties follow from a coordinate distinguishing two points. To prove regularity choose a basic neighborhood $\bigcap_{i\in E}\pi_i^{-1}(U_i)$ of $x$, with $E$ finite. In each relevant coordinate shrink $U_i$ to $V_i$ about $x_i$ with $\cl V_i\subseteq U_i$. The closure of the resulting basic neighborhood is contained in $\bigcap_{i\in E}\pi_i^{-1}(\cl V_i)$, so shrinking proves regularity. For complete regularity choose $f_i(x_i)=1$ and $f_i=0$ outside $U_i$. The function $\prod_{i\in E}f_i\circ\pi_i$ is $1$ at $x$ and zero outside the basic neighborhood. For $E=\varnothing$ use the constant $1$. This yields the required function for a point outside a closed set. Products that are empty have all the stated properties.

If $Y$ is closed in a normal $X$, disjoint closed sets in $Y$ are closed in $X$, and ambient separating neighborhoods restrict to $Y$. If $q:X\to Y$ is a closed continuous surjection, pull disjoint closed sets $A,B$ back to $X$ and choose disjoint open neighborhoods $U,V$ of their inverse images. The open sets $Y\setminus q(X\setminus U)$ and $Y\setminus q(X\setminus V)$ contain $A,B$ and are disjoint: a point in both would have its nonempty fiber contained in $U\cap V$. Finally a closed surjection from a $T_1$ space has closed singleton images, so also preserves $T_1$.
\end{proof}

\begin{theorem}\label{d10:07:compact-normal}
Every compact Hausdorff space is normal, and hence is $T_4$.
\end{theorem}
\begin{proof}
Let $A,B$ be disjoint closed subsets of compact Hausdorff $X$. They are compact. If one is empty the conclusion is immediate. For $a\in A,b\in B$ choose disjoint open neighborhoods $U_{a,b},V_{a,b}$. For fixed $a$, finitely many $V_{a,b}$ cover $B$; intersect the corresponding $U_{a,b}$ to get an open neighborhood $U_a$ of $a$, and take their $V_{a,b}$-union $V_a$, which contains $B$ and is disjoint from $U_a$. Finitely many $U_a$ cover $A$. Their union and the intersection of the corresponding $V_a$ are disjoint open neighborhoods of $A,B$. Hausdorffness supplies $T_1$.
\end{proof}

\begin{proposition}\label{d10:07:metric-normal}
Every metric space is $T_4$. In fact for disjoint nonempty closed sets $A,B$ the formula
\[
 f(x)=\frac{\dist(x,A)}{\dist(x,A)+\dist(x,B)}
\]
defines a continuous function equal to $0$ on $A$ and $1$ on $B$. Every closed subset of a metric space is a $G_\delta$, that is, a countable intersection of open sets.
\end{proposition}
\begin{proof}
For a nonempty subset $A$, the function $\dist(x,A)=\inf_{a\in A}d(x,a)$ satisfies $|\dist(x,A)-\dist(y,A)|\leq d(x,y)$, by the triangle inequality and then taking infima. Its zero set is $\cl A$. Thus for closed disjoint $A,B$ the denominator is everywhere positive, giving the formula and normality by inverse images of disjoint neighborhoods of $0,1$. Empty sets can be handled by constant functions. Metric balls of radius less than $d(x,y)/2$ separate distinct points. Finally $A=\bigcap_{n\geq1}\{x:\dist(x,A)<1/n\}$ for nonempty closed $A$; for $A=\varnothing$ use the intersection of the constant empty open set.
\end{proof}

\section{The Urysohn construction}
\begin{theorem}[Urysohn separation]\label{d10:07:urysohn}
A space $X$ is normal if and only if, for any disjoint closed $A,B\subseteq X$, a continuous $f:X\to[0,1]$ exists with $f|_A=0$ and $f|_B=1$. No $T_1$ hypothesis is needed.
\end{theorem}
\begin{proof}
Suppose $X$ is normal. Put $U_1=X\setminus B$ and use shrinking to choose open $U_0$ with $A\subseteq U_0\subseteq\cl U_0\subseteq U_1$. Insert open sets at the midpoint of each consecutive pair of dyadic indices already constructed: if $r<s$ and $\cl U_r\subseteq U_s$, shrinking gives $U_{(r+s)/2}$ with
\[
 \cl U_r\subseteq U_{(r+s)/2}\subseteq\cl U_{(r+s)/2}\subseteq U_s.
\]
Induction defines $U_r$ for every dyadic $r\in[0,1]$, and $\cl U_r\subseteq U_s$ whenever $r<s$. Define
\[
 f(x)=\inf\bigl(\{r\in[0,1]:r\text{ dyadic and }x\in U_r\}\cup\{1\}\bigr).
\]
Then $f=0$ on $A$ and $f=1$ on $B$. For $0<a\leq1$,
\[
 \{f<a\}=\bigcup_{r<a}U_r,
\]
where the union is over dyadics. For $0\leq a<1$,
\[
 \{f>a\}=\bigcup_{r>a}(X\setminus\cl U_r).
\]
To verify the latter identity, if $x\notin\cl U_r$, no $U_s$ with $s<r$ contains $x$, so $f(x)\geq r>a$. Conversely if $f(x)>a$, choose dyadics $a<r<s<f(x)$. Then $x\notin U_s$, and $\cl U_r\subseteq U_s$ gives $x\notin\cl U_r$. The other ranges of $a$ give $\varnothing$ or $X$. These identities prove continuity by open rays. The construction remains valid for empty $A$ or $B$.

Conversely, the disjoint open sets $\{f<1/3\},\{f>2/3\}$ separate the prescribed closed sets, proving normality.
\end{proof}

\begin{corollary}\label{d10:07:t4-tychonoff}
Every $T_4$ space is Tychonoff. In a normal space, if $F\subseteq U$ with $F$ closed and $U$ open, there is continuous $f:X\to[0,1]$ equal to $1$ on $F$ and zero on $X\setminus U$.
\end{corollary}
\begin{proof}
The function assertion is Urysohn separation for $F,X\setminus U$, with values reversed. If $X$ is also $T_1$, apply it to $F=\{x\}$; this gives complete regularity, and \cref{d10:07:hierarchy} gives Hausdorffness.
\end{proof}

\begin{proposition}[A strengthened exercise]\label{d10:07:zero-sets}
In a normal space a closed set $A$ is the zero set of a continuous function $g:X\to[0,1]$ if and only if it is a $G_\delta$. Thus in a normal space where every closed set is a $G_\delta$, disjoint closed $A,B$ can be separated with exact fibers $f^{-1}(0)=A$ and $f^{-1}(1)=B$.
\end{proposition}
\begin{proof}
A zero set is $\bigcap_{n\geq1}\{g<1/n\}$. Conversely write $A=\bigcap_{n\geq1}U_n$ with $U_n$ open. Choose continuous $g_n:X\to[0,1]$ zero on $A$ and $1$ on $X\setminus U_n$. The series $g=\sum_{n\geq1}2^{-n}g_n$ converges uniformly and has zero set exactly $A$. Its continuity follows directly: make the tails smaller than $\eps/3$ and use continuity of the finite partial sum at the point in question. For the final assertion choose such $g_A,g_B$ and set $f=g_A/(g_A+g_B)$; the denominator is positive because $A\cap B=\varnothing$. This also covers empty prescribed sets, using a positive constant for the empty zero set; on the empty space the assertion is vacuous. Exact fibers are not part of the unrestricted Urysohn lemma.
\end{proof}

\section{Tietze extension, including the unbounded case}
\begin{lemma}\label{d10:07:approximate-extension}
Let $A$ be closed in a normal space $X$, and let continuous $h:A\to[-a,a]$ be given, with $a>0$. There is continuous $g:X\to[-a/3,a/3]$ with $|h-g|\leq2a/3$ on $A$.
\end{lemma}
\begin{proof}
The sets $A_- =\{h\leq-a/3\}$ and $A_+=\{h\geq a/3\}$ are disjoint and closed in $X$. Urysohn separation, followed by an affine change of values, gives $g=-a/3$ on $A_-$ and $g=a/3$ on $A_+$. At points with $|h|<a/3$, the estimate follows from $|g|\leq a/3$; on the two remaining regions it follows by subtracting the indicated endpoint value. Empty level sets cause no change.
\end{proof}

\begin{theorem}[Bounded Tietze extension]\label{d10:07:tietze-bounded}
If $X$ is normal, $A\subseteq X$ is closed, and $f:A\to[a,b]$ is continuous, where $a\leq b$ are finite real numbers, then $f$ has a continuous extension $F:X\to[a,b]$.
\end{theorem}
\begin{proof}
The case $a=b$ uses a constant. An affine change reduces the other cases to $f:A\to[-1,1]$. Put $r_0=f$. Repeatedly apply the lemma with $a_n=(2/3)^n$, for $n=0,1,\ldots$, to obtain continuous $g_n:X\to[-a_n/3,a_n/3]$ and residuals $r_{n+1}=r_n-g_n|_A$ satisfying $|r_{n+1}|\leq(2/3)^{n+1}$. The series
\[
 F=\sum_{n=0}^{\infty}g_n
\]
converges uniformly: its tail beginning at $N$ is bounded by $\sum_{n\geq N}(2/3)^n/3=(2/3)^N$. Its sum has absolute value at most $1$, and $f-\sum_{n=0}^{N-1}g_n|_A=r_N\to0$, so $F|_A=f$. To check continuity, fix $x$ and $\eps>0$, make both tails at $x$ and nearby points less than $\eps/3$, and use continuity of the finite partial sum to control the remaining difference by $\eps/3$. This argument needs no compactness or separation of $X$. Reversing the affine change proves the theorem, including the empty-domain case.
\end{proof}

\begin{theorem}[Real Tietze extension]\label{d10:07:tietze-real}
A space is normal if and only if every continuous real-valued function on every closed subspace has a continuous real-valued extension to the space. The function on the closed subspace may be unbounded.
\end{theorem}
\begin{proof}
Suppose $X$ is normal and $f:A\to\R$ is continuous. The homeomorphism $h(t)=t/(1+|t|)$ sends $\R$ onto $(-1,1)$. Bounded extension gives $G:X\to[-1,1]$ extending $h\circ f$. The closed set $D=G^{-1}(\{-1,1\})$ is disjoint from $A$. Choose $u:X\to[0,1]$ with $u|_A=1$ and $u|_D=0$. Then $H=uG$ takes values in $(-1,1)$: on $D$ it is zero, and off $D$ one has $|H|\leq|G|<1$. Hence $F=h^{-1}\circ H$ is continuous and extends $f$. This step is essential; simply applying $h^{-1}$ to an arbitrary bounded extension could encounter its endpoint singularities.

Conversely let disjoint closed $A,B$ be given. The function on $A\cup B$ equal to $0$ on $A$ and $1$ on $B$ is continuous, since both pieces are clopen in $A\cup B$. Extend it to a real-valued function and take inverse images of $(-\infty,1/3)$ and $(2/3,\infty)$. These disjoint open neighborhoods prove normality. Empty pieces can be separated directly.
\end{proof}

\begin{corollary}\label{d10:07:finite-coordinate-extension}
If $A$ is closed in normal $X$, a continuous map $A\to\R^n$ extends to $X\to\R^n$ for every finite $n\geq0$. A continuous map into a product of nonempty closed real intervals extends into the same product; closed rays and $\R$ itself are allowed as interval factors.
\end{corollary}
\begin{proof}
Extend each coordinate by real Tietze and, where necessary, compose with the continuous clipping map onto its interval: $t\mapsto\max(a,\min(b,t))$ for $[a,b]$, $t\mapsto\max(a,t)$ for $[a,\infty)$, and $t\mapsto\min(b,t)$ for $(-\infty,b]$. These maps fix the prescribed interval, so the extensions retain the original values. Assemble the extensions by the product universal property. The zero-coordinate product is a singleton. Choice may be used to select extensions for an arbitrary index family. Nonempty factors are necessary for an unconditional extension assertion when $A=\varnothing$. No analogous assertion for every possible target space follows: coordinate extensions need not remain in an arbitrary subspace of the product.
\end{proof}

\begin{remark}\label{d10:07:normal-scope}
Normality does not have the arbitrary-subspace or arbitrary-product permanence proved above for regularity. The proved Sorgenfrey-plane example in \cref{d10:08:sorgenfrey-plane} demonstrates product failure; combined with the cube embedding in \cref{d10:10:tychonoff-cube} and compact Hausdorff normality it also demonstrates subspace failure. This chapter adopts the scalar real extension theorem and its coordinate consequences. Absolute retract theory, vector-valued extension into locally convex spaces, and complete-normality classification from WIL section 15 and later chapters are not assumed here.
\end{remark}

% END SOURCE chapters/07.tex


% BEGIN SOURCE chapters/08.tex
\chapter{Countability, density, and Lindelof spaces}\label{ch:d10-08}
\source{MUN section 30 and its exercises, section 32; WIL section 16; STE examples 20, 82, 84.}

\section{Four distinct countability conditions}
\begin{definition}\label{d10:08:definitions}
A space is first countable if each point has a countable neighborhood basis, and second countable if its topology has a countable basis. A set $D$ is dense if $\cl D=X$, and $X$ is separable if it has a countable dense subset. A space is Lindelof if every open cover has a countable subcover. Here countable includes finite and empty, so the empty space has all four properties. A countable neighborhood basis may be replaced by a countable open neighborhood basis: choose an open neighborhood inside each member. By taking successive finite intersections it may further be made decreasing.
\end{definition}

\begin{theorem}\label{d10:08:first-countable-sequences}
In a first countable space $X$, a point belongs to $\cl A$ if and only if it is the limit of a sequence in $A$. A map $f:X\to Y$ from a first countable space is continuous if and only if it preserves limits of sequences. No separation assumption on either space is needed.
\end{theorem}
\begin{proof}
If $x\in\cl A$, let $U_n$ be a decreasing open neighborhood basis and choose $a_n\in A\cap U_n$. For every neighborhood $U$ of $x$, some $U_N\subseteq U$, and $a_n\in U$ for $n\geq N$; thus $a_n\to x$. Conversely a limit of points of $A$ belongs to $\cl A$ by the neighborhood criterion. Continuous maps preserve convergence by taking inverse images of neighborhoods. If $f$ is discontinuous at $x$, there is an open neighborhood $V$ of $f(x)$ such that every $U_n$ meets $X\setminus f^{-1}(V)$. Choose $x_n$ there. Then $x_n\to x$ but $f(x_n)$ is never in $V$, contradicting preservation of convergence. The selection uses countable choice.
\end{proof}

\begin{theorem}\label{d10:08:second-implications}
Every second countable space is first countable, separable, and Lindelof. Its every subspace is second countable and hence is also separable and Lindelof. Every discrete subspace of a second countable space is countable. A separable space has at most countably many pairwise disjoint nonempty open subsets.
\end{theorem}
\begin{proof}
If $\mathcal B$ is a countable basis, its members containing a point form a neighborhood basis. Choose one point from each nonempty member of $\mathcal B$; the resulting countable set meets every nonempty open set and is dense. Given an open cover, consider the basis elements contained in at least one covering member, and for each choose one such member. These countably many choices cover $X$. Intersecting $\mathcal B$ with a subspace gives a countable basis there. If $E$ is discrete, choose $B_e\in\mathcal B$ with $B_e\cap E=\{e\}$; distinct points give distinct basis members. Finally, from each member of a disjoint open family choose a point in a fixed countable dense set. Distinct members receive distinct points.
\end{proof}

\begin{proposition}\label{d10:08:permanence}
First countability is hereditary. Closed subspaces of Lindelof spaces are Lindelof, and open subspaces of separable spaces are separable. Continuous surjective images preserve separability and the Lindelof property; open continuous surjections preserve first and second countability. Countable products preserve first countability, second countability, and separability. A product of a Lindelof space with a compact space is Lindelof.
\end{proposition}
\begin{proof}
Intersect neighborhood bases with the subspace. To cover a closed subspace, extend its relative open cover to ambient open sets and append its open complement; a countable ambient subcover restricts as desired. For open $U$, the intersection of a dense set with $U$ is dense in $U$. A continuous surjection carries a dense set to a dense set: every nonempty open set in the target has nonempty open inverse image. Pulling an open cover back proves preservation of Lindelofness. For an open surjection $q$, images of neighborhood bases at a point give neighborhood bases at its image, and images of a global basis give a global basis; this proves the two countability assertions.

For a countable product, finite restrictions in countable coordinate bases form a countable basis, globally or at each specified point. For separability, handle the empty-product and empty-factor cases first. Otherwise choose dense countable $D_n$ in each factor and a base point $a_n\in D_n$. The points of $\prod_nD_n$ differing from $(a_n)$ in finitely many coordinates form a countable dense subset. All choices here are permitted by our standing choice convention.

Finally let $X$ be Lindelof, $K$ compact, and let $\mathcal U$ cover $X\times K$. If $K$ is empty there is nothing to prove. For each $x$, compactness of $\{x\}\times K$ supplies finitely many members of $\mathcal U$ covering this slice. The tube lemma, \cref{d10:05:tube}, gives an open neighborhood $V_x$ such that $V_x\times K$ lies in their union. Countably many $V_x$ cover $X$, and the associated finite families together form a countable subcover of the product.
\end{proof}

\begin{proposition}[A useful basis exercise]\label{d10:08:basis-extraction}
If $X$ is second countable, every basis $\mathcal C$ for $X$ contains a countable subfamily which is still a basis. Also, in a first countable $T_1$ space every singleton is a $G_\delta$.
\end{proposition}
\begin{proof}
Let $(B_n)$ be a countable basis. For each pair $(n,m)$ for which possible choose $C_{n,m}\in\mathcal C$ with $B_n\subseteq C_{n,m}\subseteq B_m$. If $x\in U$ with $U$ open, choose $x\in B_m\subseteq U$, then $x\in C\subseteq B_m$ from $\mathcal C$, and then $x\in B_n\subseteq C$. The chosen $C_{n,m}$ exists and contains $x$ while lying in $U$, proving the claim. For the second assertion let $U_n$ be a neighborhood basis at $x$. If $y\ne x$, the open set $X\setminus\{y\}$ contains some $U_n$, so $y\notin\bigcap_nU_n$. Hence $\{x\}=\bigcap_nU_n$ after replacing the $U_n$ by open neighborhoods.
\end{proof}

\section{The metric equivalence and a separation consequence}
\begin{theorem}\label{d10:08:metric-equivalence}
For a metrizable space the following are equivalent: second countability, separability, and the Lindelof property. Every metrizable space is first countable. Every compact metrizable space is second countable.
\end{theorem}
\begin{proof}
Fix a compatible metric $d$. Balls of radii $1/n$ give first countability. The implications from second countability were proved without a metric. If $D$ is countable dense, the balls $B(p,r)$ with $p\in D$ and positive rational $r$ form a countable basis: for $x\in U$ choose $\eps>0$ with $B(x,\eps)\subseteq U$, then $p\in D$ with $d(x,p)<\eps/3$ and rational $r$ with $d(x,p)<r<\eps-d(x,p)$. This ball contains $x$ and lies in $U$.

If $X$ is Lindelof, for each $n\geq1$ choose a countable subcover of the cover by $1/n$-balls. The union $D$ of all their centers is countable and dense: every $B(x,\eps)$ contains such a center once $2/n<\eps$, because a chosen $1/n$-ball covering $x$ has its center within $1/n$ of $x$. If $X$ is empty, take $D=\varnothing$. Compactness implies Lindelofness directly and therefore gives the last assertion.
\end{proof}

\begin{theorem}\label{d10:08:regular-lindelof-normal}
Every regular Lindelof space is normal. Consequently every regular second countable $T_1$ space is $T_4$.
\end{theorem}
\begin{proof}
Let $A,B$ be disjoint closed subsets. Regularity gives, for each $a\in A$, an open $U_a$ containing $a$ with $\cl U_a\cap B=\varnothing$, and similarly open $V_b$ with $\cl V_b\cap A=\varnothing$. The closed subspaces $A,B$ are Lindelof, so choose countably many $U_n$ covering $A$ and countably many $V_n$ covering $B$. Finite nonempty families may be repeated; if either closed set is empty the separation is immediate. Define
\[
 U=\bigcup_{n\geq1}\left(U_n\setminus\bigcup_{i\leq n}\cl V_i\right),\qquad
 V=\bigcup_{n\geq1}\left(V_n\setminus\bigcup_{i\leq n}\cl U_i\right).
\]
These sets are open and contain $A,B$. If a point belonged to the $n$th piece of $U$ and the $m$th piece of $V$, then $m\leq n$ would contradict its exclusion from $\cl V_m$, while $n\leq m$ would contradict its exclusion from $\cl U_n$. Thus $U,V$ are disjoint. The argument uses regularity in the convention without $T_1$; the last assertion follows by adding the given $T_1$ condition.
\end{proof}

\section{Counterexamples with complete checks}
\begin{example}\label{d10:08:lindelof-subspace-failure}
The Lindelof property is not hereditary to arbitrary, or even open, subspaces. Let $E$ be an uncountable discrete set and add a point $\infty$. Declare all subsets of $E$ open and declare a set containing $\infty$ open precisely when its complement in $E$ is finite. This is a topology. The resulting space is compact: a cover member containing $\infty$ leaves only finitely many points to cover. It is Hausdorff, since two points of $E$ are isolated and $e,\infty$ have disjoint neighborhoods $\{e\}$ and $X\setminus\{e\}$. But the open subspace $E$ has the singleton cover with no countable subcover. This construction directly checks the claimed properties and does not presuppose a compactification theorem.
\end{example}

\begin{example}\label{d10:08:cocountable}
On an uncountable set $X$, let the nonempty open sets have countable complements. This space is $T_1$ and Lindelof, but neither separable nor first countable.

The complement of a singleton is open. From any cover choose one nonempty member $U$ and then at most one covering member for each of the countably many points outside $U$. Every countable set is closed and is proper, so it cannot be dense. If $(U_n)$ were a neighborhood basis at $x$, then $\bigcap_nU_n$ has countable complement and is uncountable. Choose $y\ne x$ in this intersection; no $U_n$ is contained in the neighborhood $X\setminus\{y\}$, a contradiction. The same intersection argument shows that no singleton is a $G_\delta$.
\end{example}

\begin{example}\label{d10:08:fan}
The singleton $G_\delta$ criterion does not characterize first countability. Take $F=\{p\}\cup(\N_{\geq1}\times\N_{\geq1})$, with all pairs isolated, and neighborhoods of $p$ required to contain a tail of each branch $\{n\}\times\N_{\geq1}$. This is $T_1$: removing one point preserves each tail. It is countable, so every singleton is the countable intersection of the open complements of all the other singletons. But if $(U_n)$ were a neighborhood basis at $p$, choose $(n,k_n)\in U_n$ on the $n$th branch. The set $V=F\setminus\{(n,k_n):n\geq1\}$ is a neighborhood of $p$ and contains no $U_n$, a contradiction. This supplies a fully specified converse failure for \cref{d10:08:basis-extraction}.
\end{example}

\begin{example}\label{d10:08:sorgenfrey-line}
The Sorgenfrey line $S$ is $\R$ with basis $[a,b)$, $a<b$. It is $T_3$, first countable, separable, and Lindelof, but is not second countable and therefore is not metrizable.

Each basic interval is clopen: its complement is $(-\infty,a)\cup[b,\infty)$, open in this topology. The topology is $T_1$, so a clopen basis proves regularity by shrinking. At $x$, the intervals $[x,x+1/n)$ form a neighborhood basis. Every basic interval meets $\Q$, proving separability. If $\mathcal B$ were a countable basis, choose for each $x$ a member $B_x$ with $x\in B_x\subseteq[x,x+1)$. For $x<y$, $B_x$ cannot equal $B_y$, since $x\notin B_y$. This injects $\R$ into $\mathcal B$, a contradiction.

To prove Lindelofness, refine a given cover by intervals $[x,b_x)$ about each point. The ordinary open union $O=\bigcup_x(x,b_x)$ is Lindelof in the Euclidean topology, since it is second countable. Thus countably many $(x,b_x)$ cover $O$, and the corresponding half-open intervals cover $O$. For distinct $x<y$ in $E=\R\setminus O$ one has $b_x\leq y$, since $y\notin(x,b_x)$. The nonempty ordinary intervals $(x,b_x)$ for $x\in E$ are therefore pairwise disjoint, so choosing a rational in each shows $E$ countable. Adding one interval for each point of $E$ gives a countable cover, and passing to the original covering members proves the assertion. By \cref{d10:08:regular-lindelof-normal}, $S$ is also normal.
\end{example}

\begin{lemma}[Jones's counting argument]\label{d10:08:jones}
If a space has a dense subset $D$ and a closed discrete subspace $E$ with $|E|\geq2^{|D|}$, then it is not normal.
\end{lemma}
\begin{proof}
If it were normal, each $T\subseteq E$ and its complement $E\setminus T$, being disjoint closed subsets of the space, would have disjoint open neighborhoods $U_T,V_T$. The map $T\mapsto U_T\cap D$ is injective. Indeed if $e\in T\setminus T'$, then the nonempty open set $U_T\cap V_{T'}$ meets $D$; any such point lies in $U_T\cap D$ and outside $U_{T'}\cap D$. Reverse the roles if necessary. Thus $2^{|E|}\leq2^{|D|}$, contradicting Cantor's theorem and $|E|\geq2^{|D|}$.
\end{proof}

\begin{example}\label{d10:08:sorgenfrey-plane}
The plane $S\times S$ is separable, first countable, and $T_3$, but is neither Lindelof nor normal. It has a closed subspace which is not separable. Hence separability need not pass even to closed subspaces, and products of two Lindelof normal spaces need be neither Lindelof nor normal.

The product permanence results give the positive properties, with dense set $\Q^2$. The anti-diagonal $E=\{(t,-t):t\in\R\}$ is closed: at a point off it choose a sufficiently small basic half-open rectangle on which the sum of coordinates has a fixed sign. It is discrete, because $[t,t+\eps)\times[-t,-t+\eps)$ meets $E$ only in $(t,-t)$. It is uncountable, so cannot be separable or Lindelof as a discrete space. Closed heredity of Lindelofness proves the plane is not Lindelof. Finally $|E|=2^{\aleph_0}$ while $\Q^2$ is countable, so the lemma proves nonnormality. The properties of each factor were established in the previous example.
\end{example}

\begin{remark}\label{d10:08:implication-scope}
The equivalence of the three global conditions requires metrizability. The Sorgenfrey line disproves its unrestricted extension even among $T_3$ spaces. First countability alone is much weaker: an uncountable discrete metric space is first countable but is neither separable, Lindelof, nor second countable. The source exercises involving larger cardinal invariants, hereditary Lindelof spaces, and general product cardinal bounds are not used as unproved implications here.
\end{remark}

% END SOURCE chapters/08.tex


% BEGIN SOURCE chapters/09.tex
\chapter{Metrics, completeness, and total boundedness}\label{ch:d10-09}
\source{MUN sections 20--21, 43 and exercise 43.9, section 45; WIL sections 2, 22, 24 and exercises 24A--24D.}

\section{The topology and the uniform information of a metric}
\begin{definition}\label{d10:09:metric}
A metric on $X$ is a function $d:X\times X\to[0,\infty)$ satisfying symmetry, $d(x,y)=0$ if and only if $x=y$, and $d(x,z)\leq d(x,y)+d(y,z)$. Its open balls are $B_d(x,r)=\{y:d(x,y)<r\}$ for $r>0$. A set is open if it contains a ball about each of its points. A space is metrizable if some metric induces its topology. A pseudometric satisfies the same axioms except that distinct points may have distance zero.
\end{definition}

\begin{proposition}\label{d10:09:metric-topology}
Metric balls form a basis. The restricted metric gives the subspace topology, and every metric space is Hausdorff and first countable. A map between metric spaces is continuous precisely when it satisfies the pointwise $\eps$--$\delta$ condition. Also $x\in\cl A$ precisely when $B(x,r)\cap A\ne\varnothing$ for every $r>0$, and a sequence converges to $x$ precisely when its distances to $x$ tend to zero.
\end{proposition}
\begin{proof}
If $y\in B(x,r)$, then $B(y,r-d(x,y))\subseteq B(x,r)$ by the triangle inequality. Intersections of two balls containing a point therefore contain a ball about it, proving the basis assertion. Intersecting balls with a subspace proves the restriction assertion. Distinct points have disjoint balls of radius $d(x,y)/3$, and the balls $B(x,1/n)$ form a countable neighborhood basis. Taking inverse images of balls is exactly the pointwise $\eps$--$\delta$ condition. The closure and convergence assertions are the corresponding neighborhood tests applied to this basis.
\end{proof}

\begin{definition}\label{d10:09:equivalences}
Metrics on the same set are topologically equivalent if they induce the same open sets. They are uniformly equivalent if both identity maps between the metric spaces are uniformly continuous. Uniform continuity of $f:(X,d)\to(Y,e)$ means that for every $\eps>0$ some $\delta>0$ satisfies $d(x,x')<\delta\Rightarrow e(f(x),f(x'))<\eps$, independently of the points. A map is $L$-Lipschitz if $e(f(x),f(x'))\leq Ld(x,x')$. An isometry preserves all distances; a surjective isometry is an isometric isomorphism. Metrics are bilipschitz equivalent when $cd\leq e\leq Cd$ for constants $c,C>0$.
\end{definition}

\begin{proposition}\label{d10:09:metric-comparison}
Two metrics $d,e$ give the same topology if and only if for every $x$ and every $\eps>0$ some $\delta>0$ satisfies $B_d(x,\delta)\subseteq B_e(x,\eps)$, and conversely with $d,e$ interchanged. The metrics $d$ and $\min(1,d)$ are uniformly equivalent. Bilipschitz equivalence implies uniform equivalence, which implies topological equivalence; neither converse holds in general.
\end{proposition}
\begin{proof}
The first criterion is continuity of both identity maps using the preceding proposition. For $\bar d=\min(1,d)$, the triangle inequality follows from $\min(1,s+t)\leq\min(1,s)+\min(1,t)$. Balls of radius less than $1$ agree for $d,\bar d$, and the same cutoff works independently of $x$, proving uniform equivalence. The implication from bilipschitz estimates is the $\eps$--$\delta$ estimate with constant $C$ and then $1/c$. On $\R$, the metrics $|x-y|$ and $\min(1,|x-y|)$ are uniformly equivalent but not bilipschitz equivalent, because the latter is bounded. On $(0,1)$, the usual metric and $e(x,y)=|h(x)-h(y)|$, with $h(x)=\tan(\pi(x-1/2))$, give the same topology. The sequence $1/(n+1)$, $n\geq1$, is Cauchy in the first metric and not in the second, so the metrics cannot be uniformly equivalent, as the next proposition also explains.
\end{proof}

\section{Cauchy sequences and complete spaces}
\begin{definition}\label{d10:09:complete-definition}
A sequence $(x_n)_{n\geq1}$ is Cauchy if for every $\eps>0$ some $N$ satisfies $d(x_m,x_n)<\eps$ for all $m,n\geq N$. A metric space is complete if every Cauchy sequence converges in it. A topological space is completely metrizable if it has a compatible complete metric. Thus completeness refers to a specified metric, while complete metrizability refers to a topology. The empty metric space is complete.
\end{definition}

\begin{proposition}\label{d10:09:cauchy-basics}
Every convergent sequence is Cauchy. A Cauchy sequence with a convergent subsequence converges to that same limit. Uniformly continuous maps carry Cauchy sequences to Cauchy sequences. Closed subspaces of complete metric spaces are complete, and complete subspaces of any metric space are closed. Completeness is invariant under uniform equivalence of metrics.
\end{proposition}
\begin{proof}
For a convergent sequence use two distances less than $\eps/2$ to the limit. For a Cauchy sequence with a subsequence tending to $x$, make all sufficiently late terms pairwise closer than $\eps/2$ and pick a sufficiently late subsequence term within $\eps/2$ of $x$. The triangle inequality then puts every late term within $\eps$ of $x$. Uniform continuity transfers the Cauchy test using its $\delta$ for the chosen $\eps$. A Cauchy sequence in a closed subspace has a limit in the complete ambient space, and closedness retains that limit. Conversely, if $A$ is complete and $x\in\cl A$, choose $a_n\in A$ with $d(a_n,x)<1/n$. It is Cauchy in $A$, hence converges to a point of $A$, which must equal $x$ by uniqueness of metric limits. For uniform equivalence use preservation of Cauchy sequences and the equality of the two topologies, which preserves limits.
\end{proof}

\begin{example}\label{d10:09:complete-examples}
Euclidean $\R^n$ is complete: a Cauchy sequence has Cauchy coordinate sequences, which converge by real completeness; the maximum coordinate distance tends to zero, and $\max_i|v_i|\leq\|v\|_2\leq\sqrt n\max_i|v_i|$ gives the Euclidean conclusion. This also handles $n=0$ as a singleton. The usual $(0,1)$ is incomplete, while its pullback metric through the homeomorphism $h$ in \cref{d10:09:metric-comparison} is complete, since it is isometric to $\R$. Thus a topology does not determine completeness of an arbitrarily selected compatible metric.
\end{example}

\begin{theorem}[Cantor intersection criterion]\label{d10:09:cantor-intersection}
A metric space is complete if and only if every decreasing sequence $(F_n)$ of nonempty closed subsets with $\diam F_n\to0$ has nonempty intersection. Such an intersection has exactly one point. Here $\diam A=\sup\{d(a,b):a,b\in A\}$, allowing $+\infty$, and $\diam\varnothing=0$.
\end{theorem}
\begin{proof}
If complete, choose $x_n\in F_n$. For $m,n\geq N$ the distance is at most $\diam F_N$, so $(x_n)$ is Cauchy. Its limit lies in every $F_N$, by closedness. Two intersection points have distance at most every $\diam F_n$, hence coincide. Conversely let $(x_n)$ be Cauchy and put $F_n=\cl\{x_k:k\geq n\}$. These sets are nonempty, closed, and decreasing. Their diameters tend to zero: any distance bound on the tail passes to its closure by continuity of the metric. An intersection point $x$ is approached by $(x_n)$, since $x_n,x\in F_n$ and the diameters tend to zero. For the empty space both conditions hold vacuously.
\end{proof}

\section{Completion by Cauchy sequences}
\begin{theorem}\label{d10:09:completion}
Every metric space $(X,d)$ has a completion: a complete metric space $(\widehat X,\widehat d)$ and an isometric embedding $j:X\to\widehat X$ with dense image.
\end{theorem}
\begin{proof}
Let $\mathcal C$ be the set of Cauchy sequences in $X$. Declare $(x_n)\sim(y_n)$ when $d(x_n,y_n)\to0$. Reflexivity, symmetry, and transitivity follow from the metric axioms. For two Cauchy sequences the real sequence $d(x_n,y_n)$ is Cauchy, because
\[
 |d(x_n,y_n)-d(x_m,y_m)|\leq d(x_n,x_m)+d(y_n,y_m).
\]
Its limit therefore exists and is finite. Define
\[
 \widehat X=\mathcal C/{\sim},\qquad
 \widehat d([(x_n)],[(y_n)])=\lim_n d(x_n,y_n).
\]
Replacing either representative changes the displayed distances by at most quantities tending to zero; thus the formula is well defined. Nonnegativity, symmetry, and the triangle inequality pass to the limit. Distance zero is exactly the equivalence relation, proving positive definiteness.

Send $x$ to the class of its constant sequence. This preserves distances and is injective. If $q=[(x_n)]$, then $j(x_n)\to q$: for $n$ large, the distances $d(x_n,x_m)$ are small for every sufficiently large $m$, hence so is their limit as $m\to\infty$. This proves density.

To prove completeness let $(q_n)$ be Cauchy in $\widehat X$. By density choose $x_n\in X$ with $\widehat d(j(x_n),q_n)<1/n$. The triangle inequality shows $(x_n)$ is Cauchy in $X$. Put $q=[(x_n)]$. The preceding paragraph gives $j(x_n)\to q$, and the same $1/n$ estimate gives $q_n\to q$. If $X=\varnothing$, there are no sequences in $\mathcal C$, so $\widehat X=\varnothing$; the conclusion holds without choosing any points.
\end{proof}

\begin{theorem}[Universal property and uniqueness]\label{d10:09:completion-universal}
If $Y$ is a complete metric space and $f:X\to Y$ is uniformly continuous, there is a unique uniformly continuous $\widehat f:\widehat X\to Y$ with $\widehat f\circ j=f$. If $f$ is $L$-Lipschitz, so is $\widehat f$. Completions are unique up to a unique surjective isometry which commutes with the embeddings of $X$.
\end{theorem}
\begin{proof}
For $q=[(x_n)]$ define $\widehat f(q)=\lim_n f(x_n)$; the limit exists because uniformly continuous maps preserve Cauchy sequences. If $(x_n)\sim(y_n)$, uniform continuity makes $e(f(x_n),f(y_n))\to0$, so the limits agree. The formula extends $f$. For uniform continuity, given $\eps>0$ choose $\delta>0$ making $d(x,y)<\delta$ imply $e(f(x),f(y))<\eps/2$. If $\widehat d(q,r)<\delta/2$, representatives satisfy $d(x_n,y_n)<\delta$ for all large $n$; passage to limits gives $e(\widehat f(q),\widehat f(r))\leq\eps/2<\eps$. Lipschitz estimates pass to limits in the same way. Any continuous extension must take $q=\lim j(x_n)$ to the displayed limit, proving uniqueness even among continuous extensions.

For any other completion $(Z,i)$, apply the formula to the isometric map $i:X\to Z$. The resulting map $T:\widehat X\to Z$ preserves distances by taking limits. Its image is complete, hence closed, and contains the dense set $i(X)$, so it is all of $Z$. Uniqueness follows from density. This argument applies as well when $X$ and both completions are empty. The universal property concerns uniformly continuous maps; an arbitrary continuous map need not extend to a completion, as $x\mapsto1/x$ on $(0,1)$ illustrates.
\end{proof}

\section{Finite nets and compactness}
\begin{definition}\label{d10:09:total-bounded}
A metric space is totally bounded if for each $\eps>0$ it is the union of finitely many open $\eps$-balls with centers in the space. The empty union covers the empty space. The finite centers are often called an $\eps$-net; this finite approximation terminology is unrelated to a Moore--Smith net.
\end{definition}

\begin{proposition}\label{d10:09:total-bounded-basics}
Total boundedness is hereditary to subspaces, passes to closures in a metric space, and is preserved by uniformly continuous maps. A totally bounded space is bounded and separable. Boundedness alone does not imply total boundedness. A metric space is totally bounded if and only if every sequence has a Cauchy subsequence.
\end{proposition}
\begin{proof}
Cover the ambient totally bounded space by finitely many $\eps/2$-balls. For a subspace choose, in each ball meeting it, a point of the subspace. The corresponding $\eps$-balls cover the subspace by the triangle inequality. A finite cover of a subset by $\eps/2$-balls puts its closure inside the union of the corresponding closed $\eps/2$-balls, hence inside open $\eps$-balls. For a uniformly continuous image use finitely many $\delta$-balls, with $\delta$ chosen for the desired image radius, and map their centers.

For a nonempty totally bounded space a finite cover by unit balls gives a finite upper bound on all distances by the triangle inequality and the finitely many center distances. Choose finite $1/n$-nets; their union is countable dense. The infinite discrete metric $d(x,y)=1$ for $x\ne y$ is bounded but requires infinitely many $1/2$-balls.

Given a sequence in a totally bounded space, one ball of radius $1$ contains infinitely many of its terms. Among these terms one ball of radius $1/2$ contains infinitely many, and inductively obtain decreasing infinite index sets with their terms in a single $1/k$-ball at stage $k$. Choose increasing indices from these sets. Late chosen terms are within $2/k$ of one another, giving a Cauchy subsequence. Conversely, if no finite collection of $\eps$-balls covers $X$, recursively choose $x_n$ outside the balls about previously chosen points. The pairwise distances are at least $\eps$, so no subsequence is Cauchy. Empty and finite cases obey all assertions directly.
\end{proof}

\begin{theorem}\label{d10:09:compact-equivalences}
For a metric space the following are equivalent: compactness; sequential compactness (every sequence has a convergent subsequence); completeness and total boundedness.
\end{theorem}
\begin{proof}
Compactness gives total boundedness by applying a finite subcover to the cover by all $\eps$-balls. A Cauchy sequence in a nonempty compact space has a cluster point: the closed tail closures have the finite intersection property, so intersect. At an intersection point choose a subsequence approaching it using $1/k$-balls and increasing indices; every neighborhood meets every tail. The Cauchy-subsequence principle then makes the whole sequence converge. Thus the space is complete. Completeness and total boundedness give sequential compactness by the preceding proposition.

Suppose now that $X$ is sequentially compact. It is totally bounded, since otherwise the separated sequence constructed above would have no convergent subsequence. For any open cover $\mathcal U$ there is a $\delta>0$ such that every ball $B(x,\delta)$ is contained in some member. If not, choose $x_n$ for which $B(x_n,1/n)$ is contained in no member. A subsequence converges to $x$. Choose $U\in\mathcal U$ and $r>0$ with $B(x,r)\subseteq U$. For sufficiently late subsequence indices, $d(x_n,x)<r/2$ and $1/n<r/2$, forcing $B(x_n,1/n)\subseteq U$, a contradiction. A finite cover by $\delta$-balls, supplied by total boundedness, now gives a finite subcover from $\mathcal U$. The empty space is compact by its empty subcover and the other two conditions are vacuous there.
\end{proof}

\begin{corollary}\label{d10:09:precompact-completion}
A metric space is totally bounded if and only if its completion is compact. For a subset $A$ of a complete metric space, $\cl A$ is compact if and only if $A$ is totally bounded. Uniformly equivalent metrics have the same totally bounded subsets.
\end{corollary}
\begin{proof}
Total boundedness passes from $j(X)$ to its closure, which is the complete space $\widehat X$, so the compactness criterion applies. Conversely compactness gives total boundedness and then its hereditary version gives the condition on $X$. Apply the same argument to $A$ and its closed, hence complete, closure in the second assertion. The last assertion follows by applying preservation under uniformly continuous maps to both identity maps.
\end{proof}

\begin{theorem}[Lebesgue number and uniform continuity]\label{d10:09:lebesgue}
For any open cover of a compact metric space there is $\lambda>0$ such that each nonempty subset of diameter less than $\lambda$ lies in a single covering member. A continuous map from a compact metric space into any metric space is uniformly continuous.
\end{theorem}
\begin{proof}
The ball-containment assertion in the preceding proof gives $\delta>0$. For a nonempty $A$ of diameter less than $\delta$, pick $a\in A$; then $A\subseteq B(a,\delta)$, proving the first assertion with $\lambda=\delta$. Empty $X$ has no nonempty subsets and satisfies the assertion for any $\lambda$. If uniform continuity of $f$ failed, some $\eps>0$ and pairs $x_n,y_n$ would satisfy $d(x_n,y_n)<1/n$ and $e(f(x_n),f(y_n))\geq\eps$. Compactness gives a subsequence $x_{n_k}\to x$, and the triangle inequality gives $y_{n_k}\to x$. Continuity of $f$ contradicts the lower bound on the image distances.
\end{proof}

\begin{remark}\label{d10:09:metric-scope}
Compactness is topological, whereas total boundedness and completeness separately retain uniform information. The usual $(0,1)$ is totally bounded; the compatible complete metric pulled back from $\R$ is not totally bounded, since otherwise $(0,1)$ would be compact. Bounded equivalent metrics therefore must not be confused with totally bounded ones. Function-space completeness and the Ascoli theorem are proved in Chapter~13; they are not foundations of the completion constructed here. Pseudometric and uniform completions require their own separation conventions and are addressed in Chapter~12.
\end{remark}

% END SOURCE chapters/09.tex


% BEGIN SOURCE chapters/10.tex
\chapter{Metrization and embeddings}\label{ch:d10-10}
\source{MUN sections 20, 33--34; WIL sections 14, 22--24, especially 23.1 and 24.12.}

An embedding must recover neighborhoods as well as distinguish points. A family of continuous real functions can do both jobs, but point separation by itself is insufficient. This distinction is the organizing principle of the chapter.

\section{Evaluation maps and cubes}
\begin{theorem}\label{d10:10:evaluation-embedding}
Let $(f_i:X\to Y_i)_{i\in I}$ be continuous maps. Their evaluation map $e:X\to\prod_iY_i$, $e(x)=(f_i(x))$, is an embedding if and only if it is injective and the topology of $X$ is generated by the sets $f_i^{-1}(V)$, with $V\subseteq Y_i$ open. In particular it suffices that the maps distinguish distinct points and that for every open $U$ and $x\in U$, some $i$ and open $V\subseteq Y_i$ satisfy
\[
 x\in f_i^{-1}(V)\subseteq U.
\]
For $[0,1]$-valued functions, a sufficient neighborhood condition is that some $f_i$ be positive at $x$ and zero outside $U$.
\end{theorem}
\begin{proof}
The map is continuous by the product universal property. If its coordinate inverse images generate the topology, every open subset of $X$ is a union of finite intersections of such inverse images, and hence has open image in $e(X)$ under the bijection $e:X\to e(X)$. Thus its inverse is continuous. Conversely an embedding pulls the product's subbasis back to a subbasis of $X$, giving the stated topology. Under the sufficient neighborhood condition every open $U$ is a union of the indicated coordinate inverse images, so these sets generate the topology. In the final condition use $V=(0,1]$, open in $[0,1]$. Injectivity must still be checked unless already implied by a separation axiom and the neighborhood condition.
\end{proof}

\begin{theorem}\label{d10:10:tychonoff-cube}
A space is Tychonoff if and only if it embeds in a cube $[0,1]^I$ for some index set $I$. For a Tychonoff space the evaluation map indexed by all continuous $[0,1]$-valued functions is one such embedding.
\end{theorem}
\begin{proof}
If $X$ is Tychonoff, for $x\in U$ with $U$ open, complete regularity supplies a continuous $f:X\to[0,1]$ with $f(x)=1$ and $f=0$ on $X\setminus U$. Thus all such functions satisfy the neighborhood criterion. For $x\ne y$, the singleton $\{y\}$ is closed, and a function separating $x$ from it proves injectivity. The evaluation theorem gives the embedding. If $X$ is empty its map to any cube is an embedding as well. Conversely $[0,1]$ is metrizable and hence $T_4$, therefore Tychonoff by \cref{d10:07:metric-normal,d10:07:t4-tychonoff}. Arbitrary products and subspaces of Tychonoff spaces are Tychonoff by \cref{d10:07:heredity-products}. This proves the converse without presupposing any compactification theorem. Compactness of the cube is the separate product compactness result \cref{d10:05:tychonoff}.
\end{proof}

\begin{proposition}\label{d10:10:compact-point-separation}
If $X$ is compact and $Y$ Hausdorff, every continuous injection $X\to Y$ is an embedding. Thus on a compact space a point-separating family of continuous maps into Hausdorff spaces already gives an embedding into their product. Without compactness this implication fails.
\end{proposition}
\begin{proof}
The first assertion is \cref{d10:05:compact-bijection} applied to the image, which is Hausdorff as a subspace. The product of Hausdorff spaces is Hausdorff, proving the second assertion. For the last take the identity from the discrete real line to the usual real line. Its ordinary coordinate function separates points and is continuous, but the inverse is not continuous. This also explains why injectivity in the evaluation theorem cannot replace the topology condition.
\end{proof}

\section{Countable products and Urysohn metrization}
\begin{theorem}\label{d10:10:product-metric}
If $(X_n,d_n)$, $n\geq1$, are metric spaces, the formula
\[
 D(x,y)=\sum_{n=1}^{\infty}2^{-n}\min\{1,d_n(x_n,y_n)\}
\]
is a metric on their product and induces the product topology. If each $d_n$ is complete, $D$ is complete. More generally a nonempty product of topological spaces is metrizable if and only if each factor is metrizable and all but countably many factors are singletons.
\end{theorem}
\begin{proof}
Every summand is a bounded pseudometric on the product. Summing its triangle inequality gives the triangle inequality for $D$, and zero sum forces equality in every coordinate. Given finitely many restrictions $d_i(x_i,y_i)<r_i$ defining a basic neighborhood, replace each $r_i$ by $\min(r_i,1/2)$. A sufficiently small $D$-ball, with radius at most the finitely many $2^{-i}r_i$, satisfies all restrictions. Conversely, given $\eps>0$, choose $N$ so $\sum_{n>N}2^{-n}<\eps/2$ and impose $d_i(x_i,y_i)<\eps/2$ for $i\leq N$. The resulting basic neighborhood lies in $B_D(x,\eps)$. This proves equality of topologies, also for an empty product set.

If $(x^{(k)})$ is $D$-Cauchy, each coordinate sequence is $d_n$-Cauchy: to control coordinate distances below $\eta<1$, use $D$-distances below $2^{-n}\eta$. Complete coordinate metrics give limits $x_n$. The finite-head, small-tail estimate just used shows $D(x^{(k)},x)\to0$. Empty products with an empty factor are complete vacuously.

For the general necessity, fix a point $a$ in the nonempty product. Each factor is homeomorphic to the slice obtained by fixing all other coordinates, hence is metrizable. At $a$ choose a countable neighborhood basis $(U_n)$ and basic neighborhoods $V_n$ with $a\in V_n\subseteq U_n$. Let $J$ be the countable union of the finite coordinate supports of the $V_n$. If a nonsingleton factor had index $i\notin J$, its metrizable topology would provide a proper neighborhood $W_i$ of $a_i$. Some $U_n$ would lie in $\pi_i^{-1}(W_i)$, but $V_n$ allows every value of the $i$th coordinate, a contradiction. Thus only countably many factors are nonsingletons. Conversely omit singleton coordinates and use the displayed metric, with the usual finite-product variant if only finitely many remain. The zero-factor product is a singleton. The nonempty qualification is essential: a product with an empty factor is metrizable whatever the other factors are.
\end{proof}

\begin{theorem}[Urysohn metrization]\label{d10:10:urysohn-metrization}
Every regular second countable $T_1$ space is metrizable and embeds in the countable cube $[0,1]^{\N_{\geq1}}$. Equivalently the separation assumption can be written regular plus $T_0$. No separation assumption is implicit in the word regular here.
\end{theorem}
\begin{proof}
Such a space is Lindelof by \cref{d10:08:second-implications}, hence normal by \cref{d10:08:regular-lindelof-normal}. Let $(B_n)$ be a countable basis, repeating members if it is finite. The empty space embeds directly, so assume it is nonempty. For every ordered pair $(i,j)$ with $\cl B_i\subseteq B_j$, Urysohn separation provides continuous $f_{i,j}:X\to[0,1]$ with
\[
 f_{i,j}|_{\cl B_i}=1,\qquad f_{i,j}|_{X\setminus B_j}=0.
\]
These functions form a countable family. For any $x\in U$ with $U$ open, choose $x\in B_j\subseteq U$. By regular shrinking choose an open $V$ with $x\in V\subseteq\cl V\subseteq B_j$, and then $x\in B_i\subseteq V$. Thus $\cl B_i\subseteq B_j$, and $f_{i,j}$ is positive at $x$ and zero outside $U$. Since $T_1$ makes $X\setminus\{y\}$ a neighborhood of $x$ for $x\ne y$, this property also distinguishes points. Enumerate the functions as $(f_n)_{n\geq1}$, padding a finite family with zero functions. The evaluation criterion gives an embedding into $[0,1]^{\N_{\geq1}}$.

The countable product metric restricts to a compatible metric on the image, and pulling it back gives the explicit compatible metric
\[
 d(x,y)=\sum_{n=1}^{\infty}2^{-n}|f_n(x)-f_n(y)|.
\]
Positive definiteness follows from point separation, and its topology has already been proved to be exactly the original topology. Finally regularity plus $T_0$ implies $T_1$ by \cref{d10:07:hierarchy}, proving the alternative formulation. The two-point indiscrete space is regular and second countable but not metrizable, so some point-separation qualification is necessary.
\end{proof}

\begin{corollary}\label{d10:10:hilbert-cube-universal}
A space embeds in the countable cube if and only if it is second countable and $T_3$, equivalently if it is separable and metrizable. Every compact Hausdorff second countable space is metrizable. A metrizable space has a compatible totally bounded metric if and only if it is separable.
\end{corollary}
\begin{proof}
The cube is second countable by countable-product permanence, and metrizable by the product theorem; its subspaces are second countable and metrizable, hence $T_3$. Conversely the preceding theorem gives the required embedding, and \cref{d10:08:metric-equivalence} gives the separable characterization. Compact Hausdorff spaces are normal and regular, so the next assertion follows. The countable cube is compact by \cref{d10:05:interval,d10:05:tychonoff}, and therefore totally bounded for its displayed metric. Every subspace inherits total boundedness; hence the embedding gives a compatible totally bounded metric to every separable metrizable space. Conversely any totally bounded compatible metric supplies a countable dense set by \cref{d10:09:total-bounded-basics}.
\end{proof}

\begin{remark}\label{d10:10:embedding-distances}
The embedding in a cube preserves topology; it need not preserve the original distances or the original Cauchy sequences. In particular embedding an incomplete space into a compact cube does not make its image closed. A compatible complete metric, a uniformly equivalent metric, a topologically equivalent metric, and an isometric embedding answer different questions. These distinctions are witnessed by the metrics in \cref{d10:09:metric-comparison,d10:09:complete-examples}.
\end{remark}

\section{Complete metrizability and countable intersections}
\begin{theorem}\label{d10:10:complete-gdelta}
Let $A$ be a subspace of a metric space $(X,d)$.
\begin{enumerate}[label=(\alph*)]
\item If $A$ is completely metrizable, it is a $G_\delta$ subset of $X$.
\item If $X$ is complete and $A$ is a $G_\delta$ subset, then $A$ is completely metrizable.
\end{enumerate}
Consequently a metrizable space is completely metrizable if and only if its image is a $G_\delta$ in a metric completion, and in that case it is a $G_\delta$ in every metric space in which it embeds.
\end{theorem}
\begin{proof}
For (a), let $\rho$ be a complete compatible metric on $A$. If $A$ is empty it is itself open, so assume it is nonempty. For every $a\in A$ and $n\geq1$, compatibility permits an ambient open neighborhood $U_{a,n}$ such that
\[
 a\in U_{a,n},\qquad \diam_d U_{a,n}<1/n,\qquad
 \diam_\rho(U_{a,n}\cap A)<1/n.
\]
To obtain it, first choose a relative neighborhood of $a$ inside a $\rho$-ball of radius $1/(3n)$, realize it by an ambient open set, and intersect with a sufficiently small $d$-ball about $a$. Set $G_n=\bigcup_{a\in A}U_{a,n}$. Certainly $A\subseteq\bigcap_nG_n$. For $x$ in the latter intersection choose $a_n\in A$ with $x\in U_{a_n,n}$. Then $d(a_n,x)<1/n$. For a fixed $k$, choose one $U_{a,k}$ containing $x$. Since it is ambient open and $a_n\to x$ in $d$, every sufficiently late $a_n$ lies in $U_{a,k}$. Thus sufficiently late pairs have $\rho$-distance less than $1/k$. The sequence is $\rho$-Cauchy and converges in $\rho$ to some $a\in A$. Compatibility gives convergence in $d$ as well, and uniqueness of $d$-limits forces $x=a$. Hence $A=\bigcap_nG_n$.

For (b), write $A=\bigcap_{n\geq1}G_n$ with $G_n$ open. For $a\in A$ define
\[
 u_n(a)=\frac{1}{\dist_d(a,X\setminus G_n)}
\]
when $X\setminus G_n$ is nonempty, and define $u_n=0$ when $G_n=X$. The distance is strictly positive at points of $G_n$, so these functions are continuous on $A$. Map $A$ to its graph
\[
 H=\{(a,(u_n(a))_{n\geq1}):a\in A\}\subseteq X\times\R^{\N_{\geq1}}.
\]
This graph map is an embedding: it is continuous by coordinates and its inverse is the first projection. The ambient countable product has a complete metric by \cref{d10:10:product-metric}. The graph is closed. To prove this, take a convergent sequence of graph points $(a_k,(u_n(a_k)))\to(x,(t_n))$. If $x\notin G_n$, then $X\setminus G_n$ is nonempty and $\dist(a_k,X\setminus G_n)\leq d(a_k,x)\to0$, so $u_n(a_k)\to+\infty$, contradicting its finite limit $t_n$. Therefore $x\in A$, and continuity of each $u_n$ at $x$ gives $t_n=u_n(x)$, putting the limit in $H$. In a metric space sequential closedness implies closedness by \cref{d10:08:first-countable-sequences}. Thus $H$ is a closed subspace of a complete metric space and is complete. Pull back its metric to $A$. The empty case has empty closed graph and is covered as well.

For the final equivalence, embed the given metrizable space in a completion of any compatible metric and apply (a) and (b). Part (a) then applies to every other metric embedding, because complete metrizability is preserved by homeomorphism.
\end{proof}

\begin{corollary}\label{d10:10:open-gdelta-complete}
Open subspaces and countable intersections of open subspaces of a completely metrizable space are completely metrizable. In particular the irrational real numbers are completely metrizable. Closed subspaces are also completely metrizable.
\end{corollary}
\begin{proof}
Fix a compatible complete metric on the ambient space and apply the theorem to its $G_\delta$ subsets. The irrationals are $\bigcap_{q\in\Q}(\R\setminus\{q\})$, a countable intersection of open sets in complete $\R$. Closed subsets inherit completeness directly by \cref{d10:09:cauchy-basics}.
\end{proof}

\begin{remark}\label{d10:10:metrization-scope}
The proved results here include Urysohn's metrization theorem, cube embeddings, countable-product metrization, and the metric $G_\delta$ characterization of complete metrizability. The broader Nagata--Smirnov theorem and the paracompact local metrization theorem require locally finite covers and are treated with their proofs in Chapter~15. WIL section 23's development and star-refinement criteria, general metrization of topological groups, and Lavrentieff's extension of homeomorphisms are not invoked here. The direct proof of \cref{d10:10:complete-gdelta} avoids assuming that extension theorem. Cantor models, Cantor classification, and universal compact-metric images are proved in Chapter~18. Infinite-dimensional Hilbert-space classification, cut-point classifications of arcs and circles, and the full Hahn--Mazurkiewicz theorem remain outside the adopted manuscript scope.
\end{remark}

% END SOURCE chapters/10.tex


% BEGIN SOURCE chapters/11.tex
\chapter[Local compactness and compactifications]{Local compactness, proper maps, and compactifications}\label{ch:d10-11}
\source{MUN \S\S29,38 and exercise 26.12, including exercises 29.1,29.5,38.10; WIL \S\S18--19; STE Examples 110--111.}

\section{Compact neighborhoods}

\begin{definition}\label{d10:11:local-compactness}
A space is \emph{locally compact} if each point has a compact neighborhood: a compact set containing an open set containing the point. Neither compactness nor local compactness includes a separation axiom. A Hausdorff locally compact space will be called LCH.
\end{definition}

\begin{lemma}[Shrinking in an LCH space]\label{d10:11:compact-shrinking}
If $X$ is LCH, $x\in U$, and $U$ is open, there is an open $V$ such that
\[
 x\in V\subseteq\cl_X V\subseteq U,\qquad \cl_X V\text{ is compact}.
\]
More generally, if $C\subseteq U$ is compact, there is such a $V$ with $C\subseteq V$. In particular $X$ is regular.
\end{lemma}
\begin{proof}
Choose a compact neighborhood $K$ of $x$ and an open $O$ with $x\in O\subseteq K$. A compact subset of a Hausdorff space is closed by \cref{d10:05:compact-hausdorff}, and the compact Hausdorff space $K$ is regular by \cref{d10:07:compact-normal,d10:07:hierarchy}. Shrink inside $K$ to a relatively open set $W$ with $x\in W$ and $\cl_K W\subseteq O\cap U$. Write $W=K\cap G$ with $G$ open in $X$, and put $V=G\cap O$. Then $V$ is open in $X$, contains $x$, and is contained in $W$. Since $K$ is closed, $\cl_X V\subseteq\cl_K W\subseteq U$; it is a closed subset of the compact set $K$. For compact $C$, choose finitely many such $V_x$ covering $C$ and take their union. Its closure is the finite union of their compact closures and remains inside $U$.
\end{proof}

\begin{proposition}[Subspaces and exhaustions]\label{d10:11:exhaustion}
Closed subspaces of a locally compact space are locally compact. Open subspaces of an LCH space are LCH. An LCH space $X$ is $\sigma$-compact if and only if it has compact sets $(K_n)_{n\geq1}$ with
\[
 K_n\subseteq\Int_X K_{n+1},\qquad X=\bigcup_{n\geq1}K_n.
\]
For such an exhaustion every compact subset of $X$ is contained in some $K_n$.
\end{proposition}
\begin{proof}
For closed $A\subseteq X$, intersect a compact neighborhood $K$ of $a$ with $A$: $K\cap A$ is closed in $K$, hence compact, and is a neighborhood in $A$. For open $A$ in an LCH space, apply \cref{d10:11:compact-shrinking} inside $A$; the resulting compact closure is a compact neighborhood in $A$.

An exhaustion plainly gives a countable compact cover. Conversely write $X=\bigcup_{n\geq1}C_n$ with $C_n$ compact. Set $K_0=\varnothing$. Given $K_{n-1}$, cover $K_{n-1}\cup C_n$ by finitely many open sets with compact closures and let $K_n$ be the union of those closures. This gives $K_{n-1}\cup C_n\subseteq\Int K_n$. The construction also works when the set being covered is empty, using the empty union. Finally the increasing open sets $\Int K_{n+1}$ cover $X$. A compact set has a finite subcover from them and consequently lies in one $K_n$.
\end{proof}

\begin{example}\label{d10:11:rationals}
The space $\Q$ is not locally compact. If a compact neighborhood $K$ of $q$ existed, it would contain $(q-r,q+r)\cap\Q$ for some $r>0$. Inclusion into $\R$ maps $K$ to a compact, hence closed subset of $\R$. Density of the rationals would force $K$ to contain the entire real interval $[q-r/2,q+r/2]$, contradicting $K\subseteq\Q$. Thus even a metrizable subspace of an LCH space need not be locally compact.
\end{example}

\section{Two uses of the word proper}

\begin{definition}\label{d10:11:proper}
A continuous map $f:X\to Y$ is \emph{proper} here if $f^{-1}(K)$ is compact for every compact $K\subseteq Y$. A continuous closed map with compact fibers is often called a \emph{perfect map}. The following result states exactly when these definitions coincide in this chapter.
\end{definition}

\begin{theorem}[Proper and perfect maps]\label{d10:11:proper-perfect}
Every continuous closed map with compact fibers is proper, without separation assumptions. Conversely, a continuous proper map into an LCH space is closed and has compact fibers. The source need not be Hausdorff in either assertion.
\end{theorem}
\begin{proof}
Suppose first that $f$ is closed and each $f^{-1}(y)$ is compact. Let $K\subseteq Y$ be compact and let $\mathcal U$ be an open cover of $f^{-1}(K)$ by sets open in $X$; any relative open cover can be represented this way. For $y\in K$, choose a finite subfamily whose union $U_y$ covers $f^{-1}(y)$, allowing an empty subfamily if the fiber is empty. Closedness gives the open neighborhood
\[
 V_y=Y\setminus f(X\setminus U_y)
\]
of $y$, and $f^{-1}(V_y)\subseteq U_y$. Finitely many $V_y$ cover $K$, and the associated finitely many members of $\mathcal U$ cover $f^{-1}(K)$.

Now suppose $f$ is proper and $Y$ is LCH. Singleton sets are compact, so fibers are compact. Let $F$ be closed in $X$ and $y\notin f(F)$. Choose a compact neighborhood $K$ of $y$. The set $F\cap f^{-1}(K)$ is closed in the compact space $f^{-1}(K)$, hence compact. Its image $D$ is compact and therefore closed in the Hausdorff space $Y$. The open set $\Int K\setminus D$ contains $y$ and misses $f(F)$: any point of $f(F)$ lying in $K$ belongs to $D$. Thus $f(F)$ is closed.
\end{proof}

\begin{proposition}\label{d10:11:proper-calculus}
Composites of proper maps are proper. A closed inclusion is proper. A continuous map from a compact space to a Hausdorff space is proper. A proper map to a compact space has compact domain.
\end{proposition}
\begin{proof}
For a composite use $(g\circ f)^{-1}(K)=f^{-1}(g^{-1}(K))$. For a closed inclusion, the relevant preimage is the intersection of a compact set with a closed set. For the third assertion, a compact subset of the Hausdorff target is closed, so its inverse image is closed in the compact domain. The last assertion applies the definition to the entire target.
\end{proof}

\begin{example}\label{d10:11:proper-boundary}
The identity from a two-point discrete space to the two-point indiscrete space is continuous and proper, because every subset of its domain is finite and compact. It is not closed. Thus properness alone does not imply closedness without hypotheses on the target. Also, compact domain alone is insufficient for properness: the map from $[0,1]$ to the two-point indiscrete space sending $0$ to $0$ and every positive point to $1$ is continuous, but the compact singleton $\{1\}$ has inverse image $(0,1]$, whose open cover $\{(1/n,2)\cap(0,1]:n\geq1\}$ has no finite subcover.
\end{example}

\section{Adjoining one point}

\begin{theorem}[The Alexandroff construction]\label{d10:11:one-point}
Let $X$ be LCH and let $X^+=X\sqcup\{\infty\}$. Declare its open sets to be the open subsets of $X$ and the sets
\[
 X^+\setminus K\quad(K\subseteq X\text{ compact}).
\]
This is a compact Hausdorff topology, and $X$ is an open subspace with its original topology. It is the unique compact Hausdorff topology on this set inducing the topology of $X$. The subspace $X$ is dense in $X^+$ exactly when $X$ is noncompact. If $X$ is compact, $\infty$ is isolated; if $X=\varnothing$, $X^+$ is a singleton.

Conversely, if $X$ is the complement of one point in a compact Hausdorff space, then $X$ is LCH and that space has this topology.
\end{theorem}
\begin{proof}
Compact subsets of $X$ are closed. Finite intersections of the displayed neighborhoods of infinity correspond to finite unions of compact sets. Intersecting such a neighborhood with an open subset of $X$ is open in $X$. For an arbitrary union containing infinity, select one constituent $X^+\setminus K_0$. The complement of the union is a closed subset of $K_0$, hence compact: relative to $X$ it is the intersection of closed compact complements and closed complements of ordinary open sets. These observations prove the topology axioms and the subspace assertion.

Any open cover has a member containing infinity, with compact complement $K$. Finitely many additional members cover $K$, so $X^+$ is compact. Distinct points of $X$ have disjoint open neighborhoods by Hausdorffness. To separate $x$ from infinity choose an open $U$ and compact $K$ with $x\in U\subseteq K$; then $U$ and $X^+\setminus K$ are disjoint. The point infinity lies outside $\cl X$ precisely when its singleton is open, which happens precisely when $X$ is compact.

For the converse, write $Z\setminus X=\{p\}$ with $Z$ compact Hausdorff. The singleton is closed, so $X$ is open. It is LCH by the open-subspace part of \cref{d10:11:exhaustion}, since $Z$ is compact and Hausdorff. The complement in $Z$ of an open neighborhood of $p$ is compact and contained in $X$. Conversely, each compact subset of $X$ is compact in $Z$ and thus closed in $Z$. These two facts identify exactly the neighborhoods of $p$ and prove uniqueness.
\end{proof}

\begin{definition}\label{d10:11:compactification}
A \emph{Hausdorff compactification} of $X$ is an embedding $i:X\to K$ into a compact Hausdorff space with dense image. We permit $K=i(X)$, including $X=K=\varnothing$. Thus $X^+$ is a compactification in this sense exactly for noncompact LCH $X$; the construction still exists when $X$ is compact, but then is not a dense compactification.
\end{definition}

\begin{proposition}[Extension at infinity]\label{d10:11:proper-infinity}
For a continuous map $f:X\to Y$ between LCH spaces, define $f^+:X^+\to Y^+$ by $f^+|_X=f$ and $f^+(\infty)=\infty$. Then $f^+$ is continuous exactly when $f$ is proper. In particular every homeomorphism extends to the one-point constructions.
\end{proposition}
\begin{proof}
If $f$ is proper, the inverse image of $Y^+\setminus K$ is $X^+\setminus f^{-1}(K)$ and is open; inverse images of open subsets of $Y$ are open in $X$ and hence in $X^+$. Conversely, if $f^+$ is continuous, for compact $K\subseteq Y$ its inverse image $f^{-1}(K)$ is closed in the compact space $X^+$ and therefore compact. Restricting the topology to $X$ preserves that compactness. A homeomorphism and its inverse preserve compact sets, proving the last assertion.
\end{proof}

\section{The Stone--\v Cech compactification}

\begin{definition}[Evaluation cubes]\label{d10:11:evaluation-embedding}
For a Tychonoff space $X$, put $I_X=C(X,[0,1])$ and
\[
 e_X:X\longrightarrow[0,1]^{I_X},\qquad e_X(x)_u=u(x).
\]
By \cref{d10:10:tychonoff-cube}, $e_X$ is an embedding; that earlier theorem proves the cube representation independently of compactification. A compact Hausdorff space is Tychonoff by \cref{d10:07:compact-normal,d10:07:urysohn}, and its evaluation image is compact and hence closed in its Hausdorff cube.
\end{definition}

\begin{theorem}[Construction and universal property]\label{d10:11:stone-cech}
For Tychonoff $X$, let
\[
 \beta X=\cl_{[0,1]^{I_X}} e_X(X).
\]
Then $(\beta X,e_X)$ is a Hausdorff compactification. For every compact Hausdorff $H$ and continuous $f:X\to H$, there is a unique continuous $\bar f:\beta X\to H$ with $\bar f e_X=f$. This property determines the compactification up to a unique homeomorphism fixing $X$. A space has a Hausdorff compactification if and only if it is Tychonoff.
\end{theorem}
\begin{proof}
The product theorem \cref{d10:05:tychonoff} makes the cube compact; it is Hausdorff by coordinate separation. Its closed subspace $\beta X$ is compact Hausdorff, and \cref{d10:10:tychonoff-cube} gives the dense embedding. Product compactness is the choice-dependent step of this construction.

For the extension, embed $H$ by $e_H$ in its cube. Define
\[
 T:[0,1]^{I_X}\longrightarrow[0,1]^{I_H},\qquad
 T(t)_v=t_{v\circ f}\quad(v\in I_H).
\]
Every coordinate of $T$ is a projection, so $T$ is continuous, and $Te_X=e_Hf$. A continuous map sends the closure of a set into the closure of its image: if an open neighborhood of its image missed that image set, its inverse image would contradict being in the original closure. Thus
\[
 T(\beta X)\subseteq\cl(e_Hf(X))\subseteq e_H(H),
\]
the last inclusion holding because $e_H(H)$ is closed. Set $\bar f=e_H^{-1}T|_{\beta X}$. If $H$ is empty the existence of $f$ forces $X$ and $\beta X$ empty, and the assertion holds with the unique empty map.

Two continuous maps from a space into a Hausdorff space that agree on a dense subset agree everywhere: their equalizer is closed, since any unequal pair of values has disjoint neighborhoods whose inverse images give an open set in the complement of the equalizer. This proves uniqueness. If another compactification has the same property, extend the two dense embeddings in both directions. Both composites restrict to the identity on the dense image of $X$ and hence equal the corresponding identity maps. They are inverse homeomorphisms.

Finally, a subspace of a Tychonoff space is Tychonoff. Indeed, for a point $a$ and a relatively closed set $F$ in a subspace $A$, write $F=A\cap D$ with $D$ closed in the ambient space and $a\notin D$, and restrict a function separating $a$ from $D$. Hausdorffness is also inherited. Thus a Hausdorff compactification forces $X$ Tychonoff; the converse is the construction just given.
\end{proof}

\begin{corollary}\label{d10:11:beta-consequences}
Every bounded continuous real-valued function on $X$ extends uniquely to $\beta X$. Every Hausdorff compactification of $X$ is a continuous image of $\beta X$ by a map fixing $X$. A continuous map $f:X\to Y$ between Tychonoff spaces extends uniquely to a map $\beta f:\beta X\to\beta Y$, and
\[
 \beta(\id_X)=\id_{\beta X},\qquad \beta(g\circ f)=\beta g\circ\beta f.
\]
If $X$ is compact Hausdorff, $e_X(X)=\beta X$; if $X$ is empty, $\beta X$ is empty.
\end{corollary}
\begin{proof}
A bounded real function takes its values in a compact interval; apply the universal property to that interval. For a compactification $i:X\to K$, extend $i$. The image of the extension is compact and therefore closed in $K$, and contains the dense set $i(X)$, so is all of $K$. To define $\beta f$, extend $e_Yf$. The two displayed identities follow from uniqueness after restriction to the dense copy of the domain. A compact cube image is already closed, giving the last assertions.
\end{proof}

\begin{remark}\label{d10:11:compactification-scope}
The construction concerns bounded continuous coordinates, not extensions of arbitrary unbounded functions. The remainder of $\beta\N$, cardinalities of remainders, and the lattice of all compactifications are not adopted here. Examples 110--111 of STE demonstrate how much more complicated this remainder can be than the single point in the Alexandroff construction.
\end{remark}

% END SOURCE chapters/11.tex


% BEGIN SOURCE chapters/12.tex
\chapter{Uniform spaces and uniform continuity}\label{ch:d10-12}
\source{WIL \S\S35--39, especially 35.2--35.8,37.1--37.5,39.1--39.13 and exercises 36D,39A; MUN \S\S21,43; STE Example 137.}

\section{Entourages and their topology}

\begin{definition}\label{d10:12:uniformity}
For relations $U,V\subseteq X\times X$, write
\begin{align*}
 U^{-1}&=\{(y,x):(x,y)\in U\},\\
 U\circ V&=\{(x,z):\exists y\ (x,y)\in V,(y,z)\in U\},\\
 U[x]&=\{y:(x,y)\in U\}.
\end{align*}
A \emph{uniformity} on a nonempty set $X$ is a proper filter $\mathcal U$ of subsets of $X\times X$, its \emph{entourages}, such that every entourage contains the diagonal $\Delta_X$, $U^{-1}\in\mathcal U$ whenever $U\in\mathcal U$, and for each $U\in\mathcal U$ some $V\in\mathcal U$ satisfies $V\circ V\subseteq U$. It is \emph{separated} if $\bigcap\mathcal U=\Delta_X$.

On the empty set we separately permit the unique degenerate entourage family $\{\varnothing\}$. It is not a proper filter; this is solely an empty-uniform-space convention. All convergence and Cauchy filters below remain proper. Thus there are no such filters on the empty set, and its completeness assertions are vacuous.
\end{definition}

\begin{lemma}\label{d10:12:entourage-calculus}
Every entourage contains a symmetric entourage, and for every entourage $U$ and positive integer $m$ there is a symmetric entourage $V$ with $V^m\subseteq U$. A nonempty family $\mathcal B$ generates a uniformity by taking supersets if its members contain the diagonal, its finite intersections contain members of $\mathcal B$, each $B^{-1}$ contains a member of $\mathcal B$, and each $B$ contains $C\circ C$ for some $C\in\mathcal B$. For nonempty $X$ these conditions ensure the generated filter is proper.
\end{lemma}
\begin{proof}
Replace $U$ by $U\cap U^{-1}$ to obtain symmetry. Repeatedly choose square roots and symmetrize; if $2^k\geq m$, the resulting $V$ satisfies $V^m\subseteq V^{2^k}\subseteq U$, since the diagonal belongs to $V$. For the base criterion, upward closure is imposed, and the finite-intersection condition gives closure under intersections and nonemptiness of the family. The inverse and square-root conditions pass to supersets. Since every base member contains a nonempty diagonal when $X\ne\varnothing$, the empty set cannot enter the generated filter.
\end{proof}

\begin{theorem}[The induced topology]\label{d10:12:uniform-topology}
Declare $O\subseteq X$ open when each $x\in O$ has some $U\in\mathcal U$ with $U[x]\subseteq O$. This is a topology, and the sets $U[x]$ form a neighborhood base at $x$. It is Hausdorff exactly when the uniformity is separated. Each uniform topology is regular, without an implicit $T_1$ assertion.
\end{theorem}
\begin{proof}
The empty set and $X$ satisfy the condition. Unions do because each point lies in a constituent; intersections of two open sets do by intersecting their entourages. To verify the neighborhood assertion, for an entourage $U$ put
\[
 O=\{y:\exists W\in\mathcal U\text{ with }W[y]\subseteq U[x]\}.
\]
It contains $x$. If $y\in O$ and $W[y]\subseteq U[x]$, choose symmetric $V$ with $V^2\subseteq W$. For $z\in V[y]$, one has $V[z]\subseteq W[y]$, so $z\in O$. Thus $O$ is open and $x\in O\subseteq U[x]$.

If $(x,y)$ belongs to every entourage, then $y$ belongs to every neighborhood of $x$, which is impossible for distinct points in a Hausdorff space. Conversely, if $(x,y)\notin U$, choose symmetric $V$ with $V^2\subseteq U$. The neighborhoods $V[x],V[y]$ are disjoint, since a common point would put $(x,y)$ in $V^2$; their interiors give disjoint open neighborhoods.

For regularity, given an open $O$ containing $x$, take symmetric $V$ and $U$ with $U[x]\subseteq O$ and $V^2\subseteq U$. If $y\in\cl(V[x])$, its neighborhood $V[y]$ meets $V[x]$, so $y\in V^2[x]\subseteq O$. Hence $x\in\Int(V[x])$ and $\cl(\Int(V[x]))\subseteq O$. This is the point-neighborhood shrinking form of regularity, which does not require points to be closed. All statements for empty $X$ hold directly.
\end{proof}

\begin{lemma}[Closed entourages]\label{d10:12:closed-entourages}
Every entourage contains a closed symmetric entourage, where closedness is in the product topology of $X\times X$. More precisely, if $V$ is a symmetric entourage, then $\cl_{X\times X}V\subseteq V^3$. In particular every entourage $U$ contains a closed symmetric entourage $C$ with $C^3\subseteq U$.
\end{lemma}
\begin{proof}
If $(x,y)\in\cl V$, the product neighborhood $V[x]\times V[y]$ meets $V$, say at $(a,b)$. Symmetry gives $(x,a),(a,b),(b,y)\in V$, so $(x,y)\in V^3$. Given an entourage $U$, choose symmetric $V$ with $V^3\subseteq U$. Its closure is symmetric, contains the entourage $V$, and is contained in $U$. For the last assertion first choose a symmetric entourage $T$ with $T^3\subseteq U$, and then a closed symmetric entourage $C\subseteq T$ as just constructed. On the empty space the same statement holds with the empty relation.
\end{proof}

\begin{definition}\label{d10:12:uniform-continuity}
A map $f:(X,\mathcal U)\to(Y,\mathcal V)$ is \emph{uniformly continuous} if $(f\times f)^{-1}(V)\in\mathcal U$ for every $V\in\mathcal V$. A bijection whose two directions are uniformly continuous is a \emph{uniform isomorphism}. An injective map inducing a uniform isomorphism onto a uniform subspace is a \emph{uniform embedding}. Maps with empty domain have their evident vacuous uniform continuity.
\end{definition}

\begin{proposition}\label{d10:12:continuity-calculus}
Uniformly continuous maps are continuous for the induced topologies. Identity maps and composites are uniformly continuous. Two uniformities can induce the same topology and different notions of uniform continuity and completeness.
\end{proposition}
\begin{proof}
If $O$ is a neighborhood of $f(x)$, choose $V[f(x)]\subseteq O$. For $U=(f\times f)^{-1}(V)$, one has $f(U[x])\subseteq O$, proving continuity at $x$. Identity and composition follow from the corresponding laws of inverse images. For the distinction, on $\R$ let
\[
 d(x,y)=|x-y|,\qquad \rho(x,y)=|\arctan x-\arctan y|.
\]
Both induce the usual topology, since arctangent is a homeomorphism onto $(-\pi/2,\pi/2)$. The sequence $n$ is $\rho$-Cauchy and has no $\rho$-limit in $\R$, whereas $(\R,d)$ is complete. The identity from $\rho$ to $d$ is not uniformly continuous: $\rho(n,n+1)\to0$ but $d(n,n+1)=1$. Thus completeness of a specified uniformity is not a topological invariant. The existence of some compatible complete metric, called complete metrizability, is a different property.
\end{proof}

\section{Subspaces, products, and pseudometrics}

\begin{proposition}\label{d10:12:subspace-product}
On $A\subseteq X$ the relations $U\cap(A\times A)$ form a base for the subspace uniformity. On $P=\prod_{i\in I}X_i$ the finite intersections
\[
 \bigcap_{i\in F}(\pi_i\times\pi_i)^{-1}(U_i),\qquad F\subseteq I\text{ finite},\quad U_i\in\mathcal U_i,
\]
form a base for the product uniformity. Their induced topologies are the subspace and product topologies. A map into a product is uniformly continuous exactly when all its coordinate maps are uniformly continuous. Products and subspaces of separated uniform spaces are separated.
\end{proposition}
\begin{proof}
Restrictions preserve the diagonal, inverse, and finite intersections; a restricted square of a relation is contained in the restriction of its square. This proves the base criterion for a subspace. For products, intersect the finitely many coordinate requirements and take square roots in each coordinate. The same criterion applies. The basic neighborhoods are respectively $U[a]\cap A$ and the products restricting finitely many coordinates to $U_i[x_i]$. The neighborhood-base result of \cref{d10:12:uniform-topology} proves the asserted topological identities. Inverse images of the displayed product base establish both directions of the mapping property. Separation follows by restricting the intersection of entourages in a subspace, or detecting a distinct product coordinate. Empty subspaces and products with an empty factor receive the degenerate uniformity. The empty product is the one-point space with its unique uniformity.
\end{proof}

\begin{definition}\label{d10:12:pseudometric}
A pseudometric is a function $p:X\times X\to[0,\infty)$ satisfying $p(x,x)=0$, symmetry, and the triangle inequality; $p(x,y)=0$ need not imply $x=y$. A family $(p_i)_{i\in I}$ generates a uniformity with base given by the finite intersections of
\[
 U_{i,\eps}=\{(x,y):p_i(x,y)<\eps\}\quad(\eps>0).
\]
For the empty family this is the indiscrete uniformity. The base criterion follows from $U_{i,\eps/2}^2\subseteq U_{i,\eps}$.
\end{definition}

\begin{proposition}[Countable pseudometric families]\label{d10:12:countable-pseudometrics}
A countable family $(p_n)_{n\geq1}$ generates the same uniformity as the single pseudometric
\[
 p(x,y)=\sum_{n=1}^\infty 2^{-n}\min\{1,p_n(x,y)\}.
\]
It is a metric exactly when the family separates pairs of distinct points. For one pseudometric $q$, the relation $x\sim y$ if $q(x,y)=0$ is an equivalence relation and
\[
 q_0([x],[y])=q(x,y)
\]
is a well-defined metric on $X_0=X/{\sim}$. The uniformity on $X$ is the inverse-image uniformity of this metric under the quotient map $r:X\to X_0$.
\end{proposition}
\begin{proof}
The inequality $\min(1,a+b)\leq\min(1,a)+\min(1,b)$ proves each summand's triangle inequality; summing proves that of $p$. Symmetry and vanishing on the diagonal follow termwise. The sum vanishes exactly when all the terms vanish.

Given $\eps>0$, choose $N$ with $\sum_{n>N}2^{-n}<\eps/2$ and impose $p_n(x,y)<\eps/2$ for $1\leq n\leq N$. Then the head of the sum is less than $\eps/2$, so $p(x,y)<\eps$. Conversely, to force finitely many conditions $p_{n_j}(x,y)<a_j$, it is enough to impose
\[
 p(x,y)<\min_j 2^{-n_j}\min\{1,a_j\}.
\]
These two refinements identify the uniformities. A finite family can be extended by zero pseudometrics to apply the same formula.

For $q$, transitivity of $\sim$ follows from the triangle inequality. If $x\sim x'$ and $y\sim y'$, the two triangle inequalities give $q(x,y)=q(x',y')$. Thus $q_0$ is well defined and is a metric. The equation $q=q_0\circ(r\times r)$ identifies the entourage bases.
\end{proof}

\begin{lemma}[Dyadic chains]\label{d10:12:dyadic-chains}
Let $V_0=X\times X$ and let $(V_n)_{n\geq1}$ be symmetric entourages satisfying $V_n^3\subseteq V_{n-1}$. There is a bounded pseudometric $p$ such that, for every $n\geq0$,
\[
 V_n\subseteq\{p\leq2^{-n}\},\qquad
 \{p<2^{-(n+1)}\}\subseteq V_n.
\]
The pseudometric $p:X\times X\to\R$ is uniformly continuous for the product uniformity of the original space.
\end{lemma}
\begin{proof}
A weighted chain from $x$ to $y$ consists of finitely many vertices $x=x_0,\ldots,x_m=y$ and integers $k_1,\ldots,k_m\geq0$ with $(x_{j-1},x_j)\in V_{k_j}$. Its cost is $\sum_{j=1}^m2^{-k_j}$. The empty chain from a point to itself has cost zero. Define $p(x,y)$ as the infimum of all these costs. A one-edge chain with index zero always exists, so $0\leq p\leq1$. Reversing chains proves symmetry. Concatenating chains with costs arbitrarily close to their infima proves the triangle inequality: for every $\eta>0$, choose the two costs within $\eta/2$ of the respective infima and then let $\eta\downarrow0$. Empty chains prove $p(x,x)=0$.

We prove simultaneously for all $r\geq1$, by induction on the number of edges, that a chain of cost less than $2^{-r}$ has its endpoint pair in $V_{r-1}$. The empty chain is covered by the diagonal. For a nonempty chain let its total cost be $s<2^{-r}$. Choose the edge whose interval of partial costs contains $s/2$. The chains strictly before and after this edge have costs at most $s/2<2^{-(r+1)}$ and have fewer edges. By induction their endpoint pairs lie in $V_r$. The central edge has cost at most $s<2^{-r}$, so its index exceeds $r$ and it too belongs to $V_r$. Here the relations are nested, since the diagonal and $V_n^3\subseteq V_{n-1}$ imply $V_n\subseteq V_{n-1}$. Composing the three pieces puts the endpoint pair in $V_r^3\subseteq V_{r-1}$. This argument also proves the one-edge case, with empty outer chains, so completes the induction.

The one-edge chain of index $n$ proves $V_n\subseteq\{p\leq2^{-n}\}$. If $p(x,y)<2^{-(n+1)}$, the definition of infimum supplies a chain with cost below that number, and the preceding result puts $(x,y)$ in $V_n$.

Finally
\[
 |p(x,y)-p(x',y')|\leq p(x,x')+p(y,y').
\]
Given $\eps>0$, choose $n$ with $2\cdot2^{-n}<\eps$. Requiring $(x,x'),(y,y')\in V_n$ makes the right side at most $2\cdot2^{-n}$. These are product entourage conditions, proving uniform continuity, and hence continuity, of $p$.
\end{proof}

\begin{theorem}[Uniform pseudometrization and representation]\label{d10:12:pseudometric-representation}
A uniformity has a countable entourage base if and only if one pseudometric generates it. In this countable-base case it is generated by a metric exactly when it is separated. Every uniformity is generated by a family of bounded uniformly continuous pseudometrics, each continuous on $X\times X$ in the induced product topology.
\end{theorem}
\begin{proof}
A pseudometric supplies the countable base $\{p<1/n\}$, $n\geq1$. Conversely let $(B_n)$ be a countable entourage base, repeating members if necessary. Inductively choose symmetric $V_n$ with
\[
 V_0=X\times X,\qquad V_n^3\subseteq V_{n-1}\cap B_n.
\]
Their existence follows from \cref{d10:12:entourage-calculus}. They form a base because $V_n\subseteq B_n$. Apply \cref{d10:12:dyadic-chains}. Every $\{p<\eps\}$ contains some $V_n$, by taking $2^{-n}<\eps$, so is an original entourage. Conversely every original entourage contains some $V_n$ and hence contains $\{p<2^{-(n+1)}\}$. Thus the uniformities agree. The equality $p(x,y)=0$ is equivalent to membership in every entourage, so $p$ is a metric exactly when the uniformity is separated.

For an arbitrary entourage $U$, choose symmetric $V_1\subseteq U$ and continue a sequence with $V_n^3\subseteq V_{n-1}$, $V_0=X\times X$. Its chain pseudometric $p_U$ is uniformly continuous and satisfies $\{p_U<1/4\}\subseteq V_1\subseteq U$. Each of its pseudometric balls is an original entourage, by the same $V_n$ argument. Taking one such pseudometric for each $U$, or taking the set of all pseudometrics obtained from such sequences, gives a family generating exactly the original uniformity. Continuity on $X\times X$ follows from the final assertion of the chain lemma and \cref{d10:12:subspace-product,d10:12:continuity-calculus}. The empty space has the unique zero pseudometric and the declared degenerate uniformity, so the conclusions hold there too.
\end{proof}

\begin{corollary}[Uniformizable topologies]\label{d10:12:uniformizable}
A topology is induced by a uniformity if and only if it is completely regular, in the convention without an implicit separation axiom. A separated uniform topology is Tychonoff.
\end{corollary}
\begin{proof}
Suppose a closed set $F$ misses $x$ in a uniform space. Choose $U[x]\subseteq X\setminus F$ and a pseudometric $p$ from \cref{d10:12:pseudometric-representation} with $\{p<1/4\}\subseteq U$. Then $g(y)=\max\{0,1-4p(x,y)\}$ is continuous, is $1$ at $x$, and is zero on $F$. This proves complete regularity. Separatedness gives Hausdorffness by \cref{d10:12:uniform-topology}.

Conversely, on a completely regular space take all continuous $f:X\to[0,1]$ and the pseudometrics $p_f(x,y)=|f(x)-f(y)|$. Their finite-intersection balls at a point are open in the original topology, so their uniform topology is no finer. Given $x\in O$ with $O$ originally open, choose $f(x)=1$ and $f=0$ outside $O$. The ball $\{y:p_f(x,y)<1/2\}$ lies in $O$, proving the reverse containment of topologies. Empty $X$ is immediate under the declared convention.
\end{proof}

\section{Cauchy filters and completeness}

\begin{definition}\label{d10:12:cauchy}
A proper filter $\mathcal F$ on $X$ is \emph{Cauchy} if for every entourage $U$ some $A\in\mathcal F$ satisfies $A\times A\subseteq U$. The uniform space is \emph{complete} if every Cauchy filter converges. It is \emph{totally bounded} if for every entourage $U$ finitely many sets $U[x]$ cover $X$; the empty space is totally bounded using the empty cover.
\end{definition}

\begin{lemma}\label{d10:12:cauchy-cluster}
A convergent filter is Cauchy. A Cauchy filter with a cluster point converges to that point. Uniformly continuous maps send Cauchy filters to Cauchy filters under the image-filter operation.
\end{lemma}
\begin{proof}
If $\mathcal F\to x$, choose symmetric $V$ with $V^2\subseteq U$. The neighborhood $V[x]$ belongs to $\mathcal F$, and $V[x]\times V[x]\subseteq U$. If $x$ is instead a cluster point of a Cauchy filter, choose $A\in\mathcal F$ with $A\times A\subseteq V$ where $V$ is symmetric and $V^2\subseteq U$. Clustering means $A\cap V[x]\ne\varnothing$; choose $a$ there. Every $b\in A$ satisfies $(x,b)\in V^2\subseteq U$. Thus $U[x]\in\mathcal F$, for every $U$, giving convergence. Finally, pull back an entourage of the target, choose a small filter member in the source, and take its image; that image belongs to the image filter and is small in the target.
\end{proof}

\begin{theorem}\label{d10:12:compact-complete-tb}
A uniform space is compact if and only if it is complete and totally bounded. No separatedness is needed.
\end{theorem}
\begin{proof}
In a compact space every proper filter has a cluster point, by the filter characterization of compactness in Chapter~5. A Cauchy filter therefore converges by \cref{d10:12:cauchy-cluster}. For an entourage $U$, the interiors of the neighborhoods $U[x]$ cover $X$; a finite subcover proves total boundedness.

Conversely, let $\mathcal F$ be an ultrafilter. Given $U$, choose symmetric $V$ with $V^2\subseteq U$. Total boundedness gives a finite cover by $V[x_i]$. An ultrafilter contains one member of every finite cover: if it contained none, it would contain all their complements and hence the empty intersection. Thus $V[x_i]\in\mathcal F$ for some $i$, and this member has square contained in $U$. Hence $\mathcal F$ is Cauchy, and completeness makes it converge. The ultrafilter characterization of compactness (including its use of choice) now gives compactness. For the empty space compactness, completeness and total boundedness all hold directly.
\end{proof}

\begin{proposition}\label{d10:12:complete-subspace-product}
A closed subspace of a complete uniform space is complete. A complete subspace of a separated uniform space is closed. A product of complete uniform spaces is complete; if the product is nonempty and complete, each factor is complete.
\end{proposition}
\begin{proof}
A Cauchy filter on a subspace gives a Cauchy image filter on the ambient space. If the subspace is closed, any ambient limit belongs to it, proving the first assertion. For the second, let $x\in\cl A$. The traces on $A$ of the neighborhoods of $x$ form a proper filter. It is Cauchy since, for symmetric $V$ with $V^2\subseteq U$, the trace of $V[x]$ has square in $U$. Completeness gives a limit $a\in A$. Its ambient image also converges to $x$, so Hausdorffness forces $x=a$.

Each coordinate image of a product Cauchy filter is Cauchy, hence converges if its factor is complete. Choose coordinate limits; a finite intersection of their inverse-image neighborhoods belongs to the original filter. Thus it converges in the product topology. For the converse, fix a point of the nonempty product, embed a factor by keeping other coordinates fixed, and push a Cauchy filter forward along this uniformly continuous section. A product limit projects to a limit of the original filter. A product with an empty factor is empty and complete, but makes no assertion about its other factors.
\end{proof}

\begin{proposition}[Metric and filter completeness]\label{d10:12:metric-completeness}
For a metric uniformity, completeness as defined above is equivalent to the convergence of every Cauchy sequence.
\end{proposition}
\begin{proof}
A Cauchy sequence has a proper tail filter which is Cauchy, so uniform completeness implies sequential completeness. Conversely let $\mathcal F$ be a Cauchy filter in a sequentially complete metric space. Choose nested nonempty $A_n\in\mathcal F$ of diameter less than $2^{-n}$, using finite intersections after choosing small members, and choose $a_n\in A_n$. Then $(a_n)$ is Cauchy and has a limit $a$. For $b\in A_n$ and all $k\geq n$ one has $d(b,a_k)<2^{-n}$; passing to the real limit gives $d(b,a)\leq2^{-n}$. Thus $A_n\subseteq B(a,\eps)$ for $2^{-n}<\eps$, so the filter converges to $a$.
\end{proof}

\section{Completion and extension}

\begin{theorem}[Separated completion for pseudometric uniformities]\label{d10:12:pseudometric-completion}
Let $X$ have the uniformity of one pseudometric $p$, or equivalently of a specified countable pseudometric family as in \cref{d10:12:countable-pseudometrics}. Form the metric quotient $(X_0,p_0)$ and its metric completion $(\widehat X,\widehat p)$. The map
\[
 j:X\longrightarrow X_0\longrightarrow\widehat X
\]
has dense image and induces precisely the given uniformity on $X$ by inverse images. It is an embedding exactly when $p$ separates points. The target is complete and separated. For every uniformly continuous $f:X\to Y$ into a complete metric space, there is a unique uniformly continuous $\widehat f:\widehat X\to Y$ with $\widehat f j=f$. If $X$ is empty take $\widehat X$ empty.
\end{theorem}
\begin{proof}
The metric quotient is justified in \cref{d10:12:countable-pseudometrics}; its isometric dense completion and completeness follow from \cref{d10:09:completion,d10:12:metric-completeness}. Uniform continuity forces $f(x)=f(y)$ when $p(x,y)=0$: the distance of their images is less than every positive number. Therefore $f$ factors through a uniformly continuous map $f_0:X_0\to Y$. The metric extension result \cref{d10:09:completion-universal} yields the desired extension. For clarity, its construction takes a sequence $a_n\in X_0$ approaching $z\in\widehat X$ and sets $\widehat f(z)=\lim f_0(a_n)$; uniform continuity makes this a Cauchy sequence, makes its limit independent of the sequence, and preserves the uniform modulus after choosing a smaller target tolerance. Uniqueness follows by density and Hausdorffness of $Y$. The inverse-image assertion follows from the isometry of $X_0$ into its completion. Its kernel on $X$ is exactly $p=0$, so injectivity is precisely separation; if injective the same entourage identity makes it a uniform embedding.
\end{proof}

\begin{lemma}[Extension from a dense uniform subspace]\label{d10:12:dense-extension}
Let $A$ be dense in a uniform space $Z$, with its induced uniformity, and let $Y$ be a complete separated uniform space. Every uniformly continuous map $f:A\to Y$ extends uniquely to a uniformly continuous map $\bar f:Z\to Y$.
\end{lemma}
\begin{proof}
For empty $Z$ the assertion is immediate. Otherwise density implies $A$ is nonempty, and existence of $f$ implies $Y$ is nonempty. For $z\in Z$, the traces $A\cap N$ of neighborhoods $N$ of $z$ generate a proper filter $\mathcal F_z$ on $A$: any finite intersection is the trace of a neighborhood and meets $A$ by density. It is Cauchy. Indeed for an ambient entourage $D$, choose symmetric $V$ with $V^2\subseteq D$; then $A\cap V[z]$ is a member whose square lies in $D\cap(A\times A)$. By \cref{d10:12:cauchy-cluster}, the image filter is Cauchy in $Y$. Completeness supplies a limit and separatedness makes it unique. Define $\bar f(z)$ to be this limit. When $z\in A$, $\mathcal F_z$ converges to $z$ in $A$, so continuity of $f$ makes the image converge to $f(z)$; therefore $\bar f(z)=f(z)$.

To prove uniform continuity, let $W$ be an entourage of $Y$ and choose a symmetric closed entourage $C$ with $C^3\subseteq W$ by \cref{d10:12:closed-entourages}. Uniform continuity of $f$ gives an ambient entourage $D$ of $Z$ such that
\[
 (a,b)\in D,\quad a,b\in A\quad\Longrightarrow\quad(f(a),f(b))\in C.
\]
Choose symmetric $U$ with $U^3\subseteq D$. Suppose $(z,z')\in U$. The trace $A\cap U[z]$ belongs to $\mathcal F_z$. Since the image filter converges to $\bar f(z)$, the inverse image of the neighborhood $C[\bar f(z)]$ also belongs to $\mathcal F_z$. Properness lets us choose $a$ in their intersection. Similarly choose $b\in A\cap U[z']$ with $f(b)\in C[\bar f(z')]$. Then $(a,b)\in U^3\subseteq D$, and symmetry gives
\[
 (\bar f(z),\bar f(z'))\in C\circ C\circ C\subseteq W.
\]
Thus $(\bar f\times\bar f)^{-1}(W)$ contains the entourage $U$, proving uniform continuity. Closedness of $C$ is available from the general closed-entourage lemma; the three-factor estimate itself only needs symmetry.

For uniqueness, two continuous extensions into the Hausdorff space $Y$ cannot disagree at a point: disjoint neighborhoods of two unequal values would pull back to a neighborhood whose intersection with the dense set $A$ would force unequal values there. Hence the extension is unique even among continuous extensions.
\end{proof}

\begin{theorem}[Completion of a separated uniform space]\label{d10:12:uniform-completion}
Every separated uniform space $X$ has a complete separated uniform space $\widehat X$ and a dense uniform embedding $j:X\to\widehat X$. Every uniformly continuous $f:X\to Y$ into a complete separated uniform space extends uniquely to a uniformly continuous $\widehat f:\widehat X\to Y$. The pair $(\widehat X,j)$ is unique up to a unique uniform isomorphism commuting with $j$.
\end{theorem}
\begin{proof}
For empty $X$ choose empty $\widehat X$. Otherwise let $\mathcal P$ be a family of bounded pseudometrics generating its uniformity, as constructed in \cref{d10:12:pseudometric-representation}. For each $p\in\mathcal P$ form the metric quotient $X_p$ by zero $p$-distance and its metric completion $\widehat X_p$ from \cref{d10:09:completion}. Let
\[
 e:X\longrightarrow P=\prod_{p\in\mathcal P}\widehat X_p
\]
be the map whose $p$-coordinate is the metric-quotient map followed by its isometric inclusion. Every coordinate is uniformly continuous, so $e$ is uniformly continuous by \cref{d10:12:subspace-product}. The induced uniformity on $X$ is exactly the given one: finite coordinate metric restrictions pull back to the finite pseudometric restrictions generating it. If $e(x)=e(y)$, then $p(x,y)=0$ for every $p$, so $(x,y)$ belongs to every entourage and separatedness gives $x=y$. Therefore $e$ is a uniform embedding.

Every factor is complete in the filter sense by \cref{d10:12:metric-completeness}, so $P$ is complete by \cref{d10:12:complete-subspace-product}. It is separated by \cref{d10:12:subspace-product}. Define $\widehat X=\cl_P e(X)$. A closed subspace of a complete uniform space is complete, and separatedness passes to subspaces, proving the required properties and density. Apply \cref{d10:12:dense-extension} to the map on $e(X)$ corresponding to $f$ to obtain the universal extension.

If $(Z,i)$ is any other complete separated uniform space with a dense uniform embedding of $X$, extend $i\circ j^{-1}$ from $j(X)$ to a uniformly continuous $I:\widehat X\to Z$, and extend $j\circ i^{-1}$ to $J:Z\to\widehat X$. Their composites agree with the identities on the respective dense copies of $X$. Uniqueness in the extension lemma gives $JI=\id_{\widehat X}$ and $IJ=\id_Z$. Thus $I$ is a uniform isomorphism, and density also gives its uniqueness. This proof uses the already proved metric completion and arbitrary-product completeness, not an assumed general completion theorem.
\end{proof}

\begin{proposition}[The separated quotient]\label{d10:12:separated-quotient}
For an arbitrary uniform space put $x\sim y$ when $(x,y)$ belongs to every entourage. This is an equivalence relation. The quotient $X_0=X/{\sim}$ has a separated uniformity with base
\[
 (q\times q)(U),\qquad U\in\mathcal U,
\]
where $q:X\to X_0$ sends a point to its equivalence class. The given uniformity on $X$ is its inverse-image uniformity. Every uniformly continuous map from $X$ into a separated uniform space factors uniquely and uniformly continuously through $q$. Consequently composing $q$ with the completion of $X_0$ gives a separated completion map with dense image and the same extension property into complete separated targets. It is injective exactly when $X$ is separated.
\end{proposition}
\begin{proof}
Reflexivity and symmetry follow from the entourage axioms. For transitivity choose $V^2\subseteq U$; two equivalences give two pairs in $V$ and hence their composite in $U$, for every $U$.

The displayed image relations contain the diagonal; their finite intersections contain images of finite intersections, and inversion commutes with images. For the square-root condition choose symmetric $V$ with $V^3\subseteq U$. If classes are related successively by two image-$V$ pairs, choose representatives $(x,y)\in V$ and $(y',z)\in V$ with $y\sim y'$. Then $(y,y')\in V$, so $(x,z)\in V^3\subseteq U$. This proves the base criterion.

Each inverse image $(q\times q)^{-1}((q\times q)(U))$ contains $U$, so is an original entourage. Conversely, for symmetric $V$ with $V^3\subseteq U$, changing either representative within its equivalence class costs a pair in $V$, and therefore
\[
 (q\times q)^{-1}((q\times q)(V))\subseteq V^3\subseteq U.
\]
This proves equality of inverse-image and original uniformities. If two distinct classes were related by every quotient entourage, this inclusion would put any representatives in every original entourage, contradicting their distinctness. The quotient is thus separated.

Uniform continuity into a separated target forces equivalent points to have equal images, since their image pair belongs to every target entourage. This defines the unique factor $f_0$. For a target entourage $W$, the source entourage $(f\times f)^{-1}(W)$ has quotient image exactly $(f_0\times f_0)^{-1}(W)$, so the factor is uniformly continuous. Apply \cref{d10:12:uniform-completion} to $X_0$ for the last assertion. When $X$ is empty, take $X_0$ and its completion empty, with all statements vacuous as declared earlier.
\end{proof}

\begin{remark}[The remaining scope boundary]\label{d10:12:completion-scope}
The completion theorem now applies to arbitrary separated uniformities, including those without a countable base. For nonseparated $X$ the separated completion map identifies exactly the indistinguishable pairs; it is not an injective inclusion unless $X$ was already separated. No uniqueness theorem for nonseparated completions by inclusions is asserted. Covering-uniformity conventions, proximities, and hyperspace completion remain outside this chapter's adopted results.
\end{remark}

% END SOURCE chapters/12.tex


% BEGIN SOURCE chapters/13.tex
\chapter{Function spaces and the Ascoli theorem}\label{ch:d10-13}
\source{MUN \S\S43,45--47, especially exercises 46.3--9 and 47.2; WIL \S\S42--43, including 43.7,43.11--43.15 and exercises 43D--43E.}

\section{Three topologies on a function set}

\begin{definition}\label{d10:13:function-topologies}
Write $Y^X$ for all functions $X\to Y$ and $C(X,Y)$ for the continuous ones. The \emph{pointwise topology} on either set is the product topology, or its restriction. Its subbasic sets prescribe $f(x)\in O$ for one $x\in X$ and one open $O\subseteq Y$.

The \emph{compact-open topology} has subbasic sets
\[
 [K,O]=\{f:f(K)\subseteq O\},\qquad K\subseteq X\text{ compact},\quad O\subseteq Y\text{ open}.
\]
If $(Y,d)$ is metric, put $d_b=\min\{1,d\}$ and
\[
 D(f,g)=\sup_{x\in X}d_b(f(x),g(x)),\qquad \sup\varnothing=0.
\]
The function $D$ is a metric: it is finite, separates distinct functions by evaluating where they differ, and inherits the triangle inequality from $d_b$ by taking suprema. It gives the \emph{uniform topology}. The \emph{topology of uniform convergence on compact sets} has neighborhood base at $f$ given by finite intersections of
\[
 B_K(f,\eps)=\{g:\sup_{x\in K}d_b(f(x),g(x))<\eps\}.
\]
Its uniformity is generated by the pseudometrics $D_K(f,g)=\sup_Kd_b(f,g)$. On an empty function set these constructions give the unique topology; if $X=\varnothing$, the function set has one element.
\end{definition}

\begin{proposition}\label{d10:13:topology-comparison}
For metric $Y$, on $Y^X$ the pointwise topology is contained in the compact convergence topology, which is contained in the uniform topology. On $C(X,Y)$ the compact convergence topology equals the compact-open topology, without a separation assumption on $X$. If $X$ is compact, compact convergence and compact-open on $C(X,Y)$ equal the uniform topology. In particular the uniform topology on $C(X,Y)$ for compact $X$ depends only on the topology of $Y$, although the associated uniformity on an unrestricted domain need not do so.
\end{proposition}
\begin{proof}
Singletons are compact, giving the first containment. Every $D_K$ is bounded by $D$, giving the second. If $X$ is compact, $K=X$ supplies the reverse containment with the uniform topology.

Let $f\in[K,O]$ with $f$ continuous. The compact set $f(K)$ has a positive uniform margin inside $O$: for each $a\in f(K)$ choose $r_a>0$ with $B_d(a,2r_a)\subseteq O$, and cover $f(K)$ by finitely many $B_d(a_i,r_{a_i})$. Taking $0<\delta<\min\{1,r_{a_1},\ldots,r_{a_m}\}$ ensures that every point at distance less than $\delta$ from $f(K)$ lies in $O$. Thus $B_K(f,\delta)\subseteq[K,O]$. Empty $K$ causes no restriction.

For the other containment fix $0<\eps<1$. Cover $f(K)$ by finitely many balls $B_d(f(a_i),\eps/4)$ with $a_i\in K$ and put
\[
 K_i=\{x\in K:d(f(x),f(a_i))\leq\eps/4\},\quad
 O_i=B_d(f(a_i),\eps/2).
\]
Each $K_i$ is closed in the compact space $K$, hence compact, and the $K_i$ cover $K$. The compact-open neighborhood $\bigcap_i[K_i,O_i]$ contains $f$. For $g$ in this neighborhood and $x\in K_i$, the triangle inequality gives $d(f(x),g(x))<3\eps/4$, so its $D_K$-distance from $f$ is at most $3\eps/4<\eps$. This neighborhood lies in $B_K(f,\eps)$. Finite intersections give equality of the topologies. The definition of the compact-open topology contains no metric, proving the final assertion.
\end{proof}

\begin{example}\label{d10:13:strict-comparisons}
On $[0,1]$, $f_n(x)=x^n$ converges pointwise to the function that is zero on $[0,1)$ and one at $1$. The limit is discontinuous, and convergence is not uniform: at points sufficiently close to $1$ the difference from zero is greater than $1/2$. This shows that pointwise limits need not retain continuity.

To witness strictness within $C([0,1],\R)$ itself, use
\[
 h_n(x)=\max\{0,1-|nx-1|\},\qquad n\geq1.
\]
For $x=0$ every value is zero, and for fixed $x>0$ the value is eventually zero, so $h_n\to0$ pointwise with a continuous limit. But $h_n(1/n)=1$, so convergence is not uniform, and hence not compact convergence on this compact domain. On $\R$, $g_n(x)=x/n$ converges uniformly on every compact set to zero, because a compact subset of $\R$ is bounded. It does not converge uniformly on $\R$, since $\sup_x\min(1,|x|/n)=1$. Thus the two containments can be strict on spaces of continuous functions.
\end{example}

\begin{proposition}[Uniform limits and completeness]\label{d10:13:uniform-complete}
A uniform limit of continuous maps into a metric space is continuous, for sequences and for nets. If $Y$ is complete metric, $(C(X,Y),D)$ is complete for every topological space $X$. If $X$ is LCH, $C(X,Y)$ is also complete in the uniformity of compact convergence.
\end{proposition}
\begin{proof}
Suppose $f_\lambda\to f$ uniformly. Given $x\in X$ and $\eps>0$, choose one index $\lambda$ with $\sup_zd_b(f_\lambda(z),f(z))<\min(1,\eps/3)$. Continuity of this fixed $f_\lambda$ supplies a neighborhood $U$ of $x$ with $d(f_\lambda(y),f_\lambda(x))<\eps/3$ for $y\in U$. The three-term triangle inequality gives $d(f(y),f(x))<\eps$.

For completeness first take a $D$-Cauchy sequence $(f_n)$. At each $x$ it is $d$-Cauchy, so define $f(x)=\lim_n f_n(x)$. Given $0<\eps<1$, choose $N$ with $D(f_n,f_m)<\eps/2$ for $m,n\geq N$. Taking the pointwise limit in $m$ gives $d_b(f_n(x),f(x))\leq\eps/2$ for every $x$, whence $D(f_n,f)\leq\eps/2<\eps$. The limit is continuous by the preceding paragraph. This proves metric completeness, and \cref{d10:12:metric-completeness} gives filter completeness. If $C(X,Y)$ is empty there are no Cauchy sequences or proper filters, so the assertion holds vacuously.

Now let $\mathcal F$ be a Cauchy filter for compact convergence. For each compact $K\subseteq X$, restriction sends it to a Cauchy filter on the complete metric space $C(K,Y)$ just considered; call its unique limit $f_K$. Restriction from $K$ to any compact $L\subseteq K$ is uniformly continuous, so uniqueness gives $f_K|_L=f_L$. In particular the singleton limits assemble a function $f:X\to Y$ with $f|_K=f_K$. Every point of an LCH space has a compact neighborhood, and on its open interior $f$ is continuous as the restriction of $f_K$. Hence $f$ is continuous on $X$. The filter converges to $f$ on every compact set, and finite intersections of these convergence neighborhoods belong to it. This is convergence in the compact convergence topology and proves completeness.
\end{proof}

\section{Evaluation, parameters, and composition}

\begin{proposition}\label{d10:13:functorial-maps}
In the compact-open topology, evaluation at a fixed point $x$ is continuous, and $C(X,Y)$ is Hausdorff when $Y$ is Hausdorff. For continuous $a:W\to X$ and $b:Y\to Z$, the maps
\[
 f\longmapsto f\circ a:C(X,Y)\to C(W,Y),\qquad
 f\longmapsto b\circ f:C(X,Y)\to C(X,Z)
\]
are continuous. No local compactness hypothesis is needed for these fixed precomposition and postcomposition maps.
\end{proposition}
\begin{proof}
The inverse image of $O$ under evaluation is $[\{x\},O]$. If two functions differ, evaluate at a point where they differ and separate their values by disjoint open sets in $Y$. For precomposition the inverse image of $[K,O]$ is $[a(K),O]$, since continuous images of compact sets are compact. For postcomposition it is $[K,b^{-1}(O)]$.
\end{proof}

\begin{theorem}[Joint evaluation and the parameter criterion]\label{d10:13:evaluation-exponential}
If $X$ is LCH, joint evaluation
\[
 \operatorname{ev}:C(X,Y)\times X\to Y,\qquad(f,x)\mapsto f(x)
\]
is continuous for arbitrary $Y$. For arbitrary $X,Z,Y$, if $h:Z\times X\to Y$ is continuous, the map $\widetilde h:Z\to C(X,Y)$ given by $\widetilde h(z)(x)=h(z,x)$ is continuous. Conversely, continuity of $\widetilde h$ implies continuity of $h$ when $X$ is LCH.
\end{theorem}
\begin{proof}
If $f(x)\in O$ with $O$ open, \cref{d10:11:compact-shrinking} gives an open $V$ with $x\in V$, $K=\cl V$ compact, and $K\subseteq f^{-1}(O)$. The product neighborhood $[K,O]\times V$ of $(f,x)$ maps into $O$.

For the parameter assertion, every slice $h(z,-)$ is continuous. Suppose $\widetilde h(z_0)\in[K,O]$. For each $x\in K$, continuity of $h$ gives neighborhoods $W_x$ of $z_0$ and $V_x$ of $x$ with $h(W_x\times V_x)\subseteq O$. Finitely many $V_{x_i}$ cover $K$. Their corresponding finite intersection $W=\bigcap_i W_{x_i}$ satisfies $h(W\times K)\subseteq O$, so $\widetilde h(W)\subseteq[K,O]$. This proves continuity on a subbase. For the converse, write $h$ as the composite of $(z,x)\mapsto(\widetilde h(z),x)$ with continuous joint evaluation.
\end{proof}

\begin{theorem}[Joint composition]\label{d10:13:composition}
If the middle space $Y$ is LCH, the composition map
\[
 C(X,Y)\times C(Y,Z)\longrightarrow C(X,Z),\quad(f,g)\mapsto g\circ f
\]
is continuous in the compact-open topologies for arbitrary $X,Z$.
\end{theorem}
\begin{proof}
Suppose $g(f(K))\subseteq O$ for compact $K\subseteq X$ and open $O\subseteq Z$. The compact set $f(K)$ is contained in the open set $g^{-1}(O)$. The compact-set version of \cref{d10:11:compact-shrinking} gives an open $V$ and compact $L=\cl V$ with $f(K)\subseteq V\subseteq L\subseteq g^{-1}(O)$. Then $[K,V]\times[L,O]$ is a neighborhood of $(f,g)$, and every pair in it has composite in $[K,O]$. For a finite intersection of subbasic neighborhoods, intersect the corresponding product neighborhoods. This proves continuity.
\end{proof}

\begin{corollary}[A quotient-product application]\label{d10:13:quotient-product}
If $q:A\to B$ is a quotient map and $X$ is LCH, then $\id_X\times q:X\times A\to X\times B$ is a quotient map.
\end{corollary}
\begin{proof}
Let $p:X\times A\to W$ be the quotient of $X\times A$ by the fibers of $\id_X\times q$. There is a bijection $r:W\to X\times B$ with $rp=\id_X\times q$, and it is continuous by the quotient property of $p$. To prove its inverse continuous, regard $p$ as a continuous map $A\times X\to W$. The parameter criterion makes $a\mapsto(x\mapsto p(x,a))$ continuous into $C(X,W)$. This function is constant on every fiber of $q$, so the quotient property gives a continuous map $B\to C(X,W)$. Its associated map $X\times B\to W$ is continuous by the converse parameter criterion; it is exactly $r^{-1}$. Thus $r$ is a homeomorphism, and $\id_X\times q=rp$ is quotient.
\end{proof}

\section{Equicontinuity and finite sampling}

\begin{definition}\label{d10:13:equicontinuity}
For metric $Y$, a family $\mathcal A\subseteq C(X,Y)$ is \emph{equicontinuous at $x$} if for every $\eps>0$ there is a neighborhood $U$ of $x$ such that
\[
 d(f(y),f(x))<\eps\quad(y\in U,\ f\in\mathcal A).
\]
It is equicontinuous if this holds at every point. It is \emph{pointwise relatively compact} if $\cl_Y\{f(x):f\in\mathcal A\}$ is compact for every $x$. Boundedness at each point is a weaker condition for a general metric target.
\end{definition}

\begin{lemma}[Pointwise closure preserves equicontinuity]\label{d10:13:equicontinuous-closure}
If $\mathcal A$ is equicontinuous, its closure $H$ in $Y^X$ for the pointwise topology consists of continuous maps and is equicontinuous. On any equicontinuous family of continuous maps the pointwise and compact convergence topologies agree.
\end{lemma}
\begin{proof}
Given $x$ and $\eps>0$, choose $U$ with $d(f(y),f(x))<\eps/2$ for all $f\in\mathcal A,y\in U$. For fixed $y\in U$ the set
\[
 \{g\in Y^X:d(g(y),g(x))\leq\eps/2\}
\]
is pointwise closed, since the two evaluation maps and the distance function are continuous. It contains $\mathcal A$ and therefore $H$. Thus all $g\in H$ satisfy $d(g(y),g(x))\leq\eps/2<\eps$ on the same $U$. This is equicontinuity and gives continuity of each $g$.

For the topological equality, let $\mathcal E$ be equicontinuous, $f\in\mathcal E$, and $K$ compact. For each $x\in K$ choose an open neighborhood $U_x$ on which $d(g(y),g(x))<\eps/4$ for every $g\in\mathcal E$. Finitely many $U_{x_i}$ cover $K$. The pointwise neighborhood given by $d(g(x_i),f(x_i))<\eps/4$ for all $i$ then forces, for $y\in U_{x_i}\cap K$,
\[
 d(g(y),f(y))\leq d(g(y),g(x_i))+d(g(x_i),f(x_i))+d(f(x_i),f(y))<3\eps/4.
\]
Taking the supremum gives $D_K(g,f)\leq3\eps/4<\eps$. Thus every compact convergence neighborhood contains a pointwise neighborhood relative to $\mathcal E$. The reverse containment always holds by \cref{d10:13:topology-comparison}.
\end{proof}

\begin{lemma}[A finite-sampling estimate]\label{d10:13:finite-sampling}
If $X$ is compact and $\mathcal A\subseteq C(X,Y)$ is equicontinuous with each set $\mathcal A(x)$ totally bounded, then $\mathcal A$ is totally bounded in the uniform metric $D$.
\end{lemma}
\begin{proof}
For $0<\eps<1$ choose a finite cover $U_1,\ldots,U_m$ of $X$ with centers $x_i$ such that $d(f(x),f(x_i))<\eps/4$ whenever $x\in U_i$ and $f\in\mathcal A$. At each $x_i$, cover $\mathcal A(x_i)$ by finitely many sets of diameter less than $\eps/4$; balls of radius $\eps/8$ suffice. Assign each $f$ one containing set at each of the finitely many centers, using a fixed ordering when choices overlap. There are finitely many assignment patterns. From each nonempty pattern choose a representative $f_j$. If $f$ has representative $f_j$, then for $x\in U_i$,
\[
 d(f(x),f_j(x))<\eps/4+\eps/4+\eps/4=3\eps/4.
\]
Consequently $D(f,f_j)\leq3\eps/4<\eps$, and finitely many $D$-balls cover $\mathcal A$. If $X$ or $\mathcal A$ is empty the assertion is immediate.
\end{proof}

\section{Ascoli with precise hypotheses}

\begin{theorem}[Ascoli, metric target]\label{d10:13:ascoli}
Let $X$ be any topological space, $Y$ a metric space, and $\mathcal A\subseteq C(X,Y)$.
\begin{enumerate}[label=(\alph*)]
\item If $\mathcal A$ is equicontinuous and pointwise relatively compact, its closure in the compact-open topology is compact.
\item If $X$ is LCH, the converse holds.
\end{enumerate}
In particular, for compact Hausdorff $X$ (hence for compact metric $X$), relative compactness in the uniform topology is equivalent to equicontinuity and pointwise relative compactness. No completeness assumption on $Y$ is required.
\end{theorem}
\begin{proof}
For sufficiency, let $K_x=\cl_Y\mathcal A(x)$ and $P=\prod_{x\in X}K_x\subseteq Y^X$. The factors are compact Hausdorff; by \cref{d10:05:tychonoff}, $P$ is compact. It is closed in $Y^X$ because each $K_x$ is closed and $P$ is the intersection of their coordinate inverse images. The pointwise closure $H$ of $\mathcal A$ is a closed subset of $P$, hence compact. By \cref{d10:13:equicontinuous-closure}, $H\subseteq C(X,Y)$ and its pointwise and compact convergence topologies coincide. They also coincide with its compact-open topology by \cref{d10:13:topology-comparison}. Thus $H$ is compact in $C(X,Y)$ for the compact-open topology. That function space is Hausdorff by \cref{d10:13:functorial-maps}, so $H$ is closed there. The compact-open closure of $\mathcal A$ is contained in $H$, and is a closed subset of a compact space; hence it is compact. In fact it equals $H$, because compact-open closure is contained in pointwise closure and the two relative topologies agree on $H$.

For necessity let $H$ be the compact-open closure of $\mathcal A$ and assume it compact. Evaluation at any fixed $x$ is continuous, so $H(x)$ is compact and closed in metric $Y$. It contains $\mathcal A(x)$, and hence its closure, proving pointwise relative compactness.

Fix $x\in X$ and $\eps>0$. Joint evaluation is continuous by \cref{d10:13:evaluation-exponential}. For each $f\in H$ choose an open neighborhood $N_f$ of $f$ in $C(X,Y)$ and an open neighborhood $U_f$ of $x$ such that
\[
 d(g(y),f(x))<\eps/3\quad(g\in N_f,\ y\in U_f).
\]
Finitely many $N_{f_1},\ldots,N_{f_m}$ cover $H$. Set $U=\bigcap_iU_{f_i}$. Given $g\in H$, choose $i$ with $g\in N_{f_i}$. Both $y\in U$ and $x$ belong to $U_{f_i}$, so
\[
 d(g(y),g(x))\leq d(g(y),f_i(x))+d(f_i(x),g(x))<2\eps/3<\eps.
\]
Thus $H$, and hence $\mathcal A$, is equicontinuous. Compact Hausdorff spaces are LCH, and uniform equals compact-open on compact domains, giving the final assertion. Empty families and empty domains satisfy the result directly; the product proof also gives $H=\varnothing$ when appropriate.
\end{proof}

\begin{corollary}[The classical real-vector form]\label{d10:13:classical-ascoli}
For compact Hausdorff $X$, a family in $C(X,\R^m)$ has compact uniform closure exactly when it is equicontinuous and pointwise bounded. Such a family is compact exactly when it is also uniformly closed. In particular a uniformly bounded, uniformly Lipschitz family of functions $[a,b]\to\R^m$ has compact closure.
\end{corollary}
\begin{proof}
Closed bounded subsets of finite-dimensional Euclidean space are compact: bounded boxes are finite products of compact intervals, and closed subsets of them are compact. Thus pointwise boundedness is equivalent to pointwise relative compactness. Apply \cref{d10:13:ascoli}; a compact subset of the Hausdorff function space is closed, and a closed subfamily of the compact closure is compact. A common Lipschitz constant gives equicontinuity by $\delta=\eps/(L+1)$, while a common bound gives pointwise boundedness.
\end{proof}

\begin{remark}\label{d10:13:ascoli-boundaries}
Boundedness alone does not replace relative compactness in an arbitrary metric target: the constant functions with values in an infinite discrete bounded metric space form an equicontinuous pointwise bounded family with no compact closure. For $f_n(x)=\sin(nx)$ on $[0,2\pi]$, pointwise boundedness holds but equicontinuity at $0$ fails, since $x_n=\pi/(2n)\to0$ and $f_n(x_n)=1$. Ascoli therefore excludes compact uniform closure. The Hausdorff-uniform-target version on general compactly generated spaces, and the topology of continuous convergence, are beyond the adopted formulation; WIL \S43 describes those broader interfaces.
\end{remark}

\begin{example}[Necessity can fail outside the LCH setting]\label{d10:13:ascoli-nonlch}
Give $X=\R$ the topology with basis $O\setminus C$, where $O$ is ordinary open and $C$ is countable. It is Hausdorff because it refines the ordinary topology. Its compact subsets are precisely the finite subsets. Indeed every countable subset $B$ is closed, since its complement is open, and is discrete, since $X\setminus(B\setminus\{b\})$ isolates $b$ within $B$. If an infinite compact set existed, choose a countably infinite subset $B$ of it. This would be a closed compact discrete infinite subspace, contradicting its singleton open cover. Conversely every finite subset of any space is compact.

Consequently the compact-open topology on $C(X,\R)$ is the pointwise topology. The functions $h_n(x)=\max\{0,1-|nx-1|\}$ are continuous on $X$, as they are continuous for the coarser ordinary topology, and converge pointwise to zero. Thus $H=\{0,h_1,h_2,\ldots\}$ is compact in the compact-open topology: an open cover member containing the limit contains all sufficiently late terms, leaving finitely many points to cover.

However $H$ is not equicontinuous at zero. Any neighborhood of zero contains $(-\delta,\delta)\setminus C$ for some $\delta>0$ and countable $C$. Choose $n$ so large that $3/(2n)<\delta$. The ordinary interval $(1/(2n),3/(2n))$ is uncountable, so has a point $x$ outside $C$; there $h_n(x)>1/2$ while $h_n(0)=0$. No common neighborhood gives the equicontinuity tolerance $1/2$. Every neighborhood in $X$ is infinite, so no finite compact set can be a neighborhood; $X$ is not locally compact. This example shows directly that the necessity assertion of \cref{d10:13:ascoli} cannot be extended to all Hausdorff domains.
\end{example}

% END SOURCE chapters/13.tex


% BEGIN SOURCE chapters/14.tex
\chapter[Baire category and contractions]{Baire category and contraction fixed points}\label{ch:d10-14}
\source{MUN \S48, including exercises 48.3--8,48.10, and exercise 43.5; WIL \S25 and 24.14--24.16.}

\section{Small sets and residual sets}

\begin{definition}\label{d10:14:category-definitions}
A subset $A$ of $X$ is \emph{nowhere dense in $X$} if $\Int_X\cl_XA=\varnothing$. It is \emph{meager in $X$}, or of first category in $X$, if it is a countable union of nowhere dense subsets of $X$. It is \emph{residual in $X$}, or comeager, if its complement is meager. A set is of second category in $X$ when it is not meager in $X$. All these statements refer to the named ambient space.

A \emph{Baire space} is a space in which every countable intersection of dense open subsets is dense. The empty space is Baire.
\end{definition}

\begin{proposition}[Equivalent Baire conditions]\label{d10:14:category-equivalences}
For any space $X$, the following are equivalent:
\begin{enumerate}[label=(\alph*)]
\item $X$ is Baire;
\item no nonempty open subset of $X$ is meager in $X$;
\item every countable union of closed sets with empty interiors has empty interior;
\item every residual subset of $X$ is dense.
\end{enumerate}
Every residual set contains a $G_\delta$ set obtained as an intersection of open dense sets. That contained $G_\delta$ is dense when $X$ is Baire. Conversely, any set containing a dense $G_\delta$ subset of $X$ is residual.
\end{proposition}
\begin{proof}
Subsets of nowhere dense sets are nowhere dense, since their closures are smaller. Every meager set is contained in a countable union of closed nowhere dense sets by taking closures of its defining pieces; intersecting these pieces with the given set also shows that subsets of meager sets are meager. Thus (b) and (c) are equivalent. Taking complements converts the union in (c) into an intersection of open dense sets, proving equivalence with (a).

If (b) holds and $R$ is residual, a nonempty open set cannot be contained in $X\setminus R$, so it meets $R$; this proves (d). If (d) holds and a nonempty open $O$ were meager, $X\setminus O$ would be residual and not dense, a contradiction.

Write $X\setminus R=\bigcup_n A_n$ with $A_n$ nowhere dense. Then $\bigcap_n(X\setminus\cl A_n)\subseteq R$ is the promised intersection of open dense sets; the Baire property makes it dense. Conversely, if $G=\bigcap_nU_n$ is dense with $U_n$ open, every $U_n$ is dense, so $X\setminus U_n$ is closed nowhere dense. The complement of a set containing $G$ is a subset of their countable union and hence meager.
\end{proof}

\begin{example}[Ambient and logical distinctions]\label{d10:14:category-boundaries}
In $\R$, the rationals are dense and meager, while the irrationals are residual. Thus density alone is weaker than residuality in a Baire space. In $\Q$ every singleton is closed with empty interior, so $\Q$ is meager in itself. Consequently even the empty set is residual in $\Q$ and is not dense there. A singleton in $\R$ is nowhere dense in $\R$ but, as a one-point space, is not meager in itself.

The nonempty space $\R\sqcup\Q$ is of second category in itself but is not Baire. It is not Baire because its open subspace $\Q$ is meager in the ambient sum (the rational singleton sets remain nowhere dense there). If the whole sum were meager, intersecting its nowhere dense pieces with its open subspace $\R$ would make $\R$ meager in itself, contradicting the Baire theorem below. This last observation distinguishes ``the whole space is of second category'' from the stronger Baire condition on every nonempty open subset.
\end{example}

\section{Two nested-set proofs}

\begin{lemma}[The shrinking-diameter intersection principle]\label{d10:14:nested-closed}
In a complete metric space, if $F_1\supseteq F_2\supseteq\cdots$ are nonempty closed sets and $\diam F_n\to0$, then their intersection is a singleton.
\end{lemma}
\begin{proof}
Choose $x_n\in F_n$. For $m,n\geq N$, both points belong to $F_N$, so $d(x_m,x_n)\leq\diam F_N$. The sequence is Cauchy, and completeness gives a limit $x$. Each $F_N$ contains the entire tail and is closed, so it contains $x$. If $y$ were another point of the intersection, $d(x,y)\leq\diam F_N$ for all $N$, forcing $y=x$.
\end{proof}

\begin{theorem}[Baire for complete metric spaces]\label{d10:14:baire-metric}
Every complete metric space is Baire. Hence every completely metrizable space is Baire, regardless of whether a particular compatible metric is complete.
\end{theorem}
\begin{proof}
Let $G_n$ be dense open sets and let $O$ be any nonempty open set. Choose $x_1\in O\cap G_1$ and $0<r_1<1/2$ so that the closed ball $F_1=\{y:d(y,x_1)\leq r_1\}$ lies in $O\cap G_1$; openness supplies a larger open ball first, so such a closed ball exists. Given $x_n,r_n$, the nonempty open ball $B(x_n,r_n)$ meets $G_{n+1}$. Choose $x_{n+1}$ in that intersection and a radius $0<r_{n+1}<2^{-(n+1)}$ whose closed ball $F_{n+1}$ lies in $B(x_n,r_n)\cap G_{n+1}$. Then $F_{n+1}\subseteq F_n$ and $\diam F_n\leq2r_n\to0$. By \cref{d10:14:nested-closed}, some point lies in every $F_n$, and hence in $O\cap\bigcap_nG_n$. This proves density. If the space is empty there is no nonempty $O$ to check. Complete metrizability means exactly that such a complete compatible metric exists, and the Baire property refers only to the topology.
\end{proof}

\begin{theorem}[Baire for locally compact Hausdorff spaces]\label{d10:14:baire-lch}
Every LCH space is Baire. In particular every compact Hausdorff space is Baire.
\end{theorem}
\begin{proof}
Given dense open $G_n$ and nonempty open $O$, use \cref{d10:11:compact-shrinking} to choose a nonempty open $V_1$ with compact closure $F_1\subseteq O\cap G_1$. Inductively, $V_n\cap G_{n+1}$ is nonempty open; choose a nonempty open $V_{n+1}$ with compact closure $F_{n+1}\subseteq V_n\cap G_{n+1}$. The closed sets $F_n$ are nested nonempty subsets of the compact space $F_1$. They have the finite intersection property, so \cref{d10:05:fip} gives $\bigcap_nF_n\ne\varnothing$. Any point of this intersection lies in $O$ and in all $G_n$. No bound on diameters is used in this proof. Hausdorffness enters the compact-neighborhood shrinking lemma; no Hausdorff assumption is built into the definition of a Baire space.
\end{proof}

\begin{proposition}[Open and dense $G_\delta$ subspaces]\label{d10:14:baire-subspaces}
Open subspaces of a Baire space are Baire. Dense $G_\delta$ subspaces of a Baire space are Baire. Every $G_\delta$ subspace of a complete metric space or of an LCH space is Baire, even when it is not dense in the original space.
\end{proposition}
\begin{proof}
First let $O\subseteq X$ be open. If $A$ is nowhere dense in $O$, then it is nowhere dense in $X$. Indeed $\cl_XA\subseteq\cl_XO$. If a nonempty open $W$ were contained in $\cl_XA$, it would meet $O$, since a nonempty open subset of $\cl_XO$ meets $O$. The nonempty open set $W\cap O$ would then lie in $\cl_OA$, a contradiction. Thus meager subsets of $O$ lift to meager subsets of $X$. A nonempty relatively open subset of $O$ is open in $X$, so cannot be meager. The criterion of \cref{d10:14:category-equivalences} proves $O$ Baire.

Suppose next that $D=\bigcap_nH_n$ is dense in Baire $X$, with $H_n$ open. Let $G_n$ be open dense subsets of $D$, and write $G_n=D\cap U_n$ with $U_n$ open in $X$. Each $H_n$ is dense because it contains $D$. Each $U_n$ is dense as well: if a nonempty open $W\subseteq X$ missed $U_n$, then $W\cap D$ would be a nonempty relatively open set missing $G_n$. The Baire property makes $\bigcap_nH_n\cap\bigcap_nU_n=\bigcap_nG_n$ dense in $X$, and hence dense in $D$.

Finally let $D$ be an arbitrary $G_\delta$ in an indicated ambient space and put $F=\cl D$. The closed subspace $F$ is complete metric in the first case and LCH in the second, so it is Baire by the preceding theorems. The set $D$ is a dense $G_\delta$ in $F$, so the second part applies. This is a direct category argument; the stronger complete-metrizability theorem for metric $G_\delta$ subspaces is proved in \cref{d10:10:complete-gdelta}.
\end{proof}

\begin{corollary}\label{d10:14:countable-dense}
The irrationals are Baire. No countable dense subset of $\R$ is a $G_\delta$ subset of $\R$.
\end{corollary}
\begin{proof}
The irrationals are the intersection of the open sets $\R\setminus\{q\}$ for $q\in\Q$, so are Baire by \cref{d10:14:baire-subspaces}. If a countable dense $D\subseteq\R$ were $G_\delta$, it would be Baire by the same proposition. It has no isolated points: every real neighborhood contains a smaller nonempty open interval avoiding its center, which meets $D$. Thus each singleton is nowhere dense in $D$, and countability would make the nonempty space $D$ meager in itself, contradicting its Baire property.
\end{proof}

\section{Consequences for pointwise limits}

\begin{theorem}[Continuity on a residual set]\label{d10:14:baire-one}
Let $X$ be Baire, $Y$ metric, and $f_n:X\to Y$ continuous with $f_n(x)\to f(x)$ for every $x$. Then the set of continuity points of $f$ is a dense $G_\delta$ and is residual in $X$.
\end{theorem}
\begin{proof}
For $\eps>0$ and $N\geq1$ put
\[
 A_N(\eps)=\bigcap_{n,m\geq N}\{x:d(f_n(x),f_m(x))\leq\eps\}.
\]
Each $A_N(\eps)$ is closed by continuity, and their union is $X$ because every convergent sequence is Cauchy. The open set $U_\eps=\bigcup_N\Int A_N(\eps)$ is dense. To see this, let $V$ be nonempty open. It is Baire by \cref{d10:14:baire-subspaces}. The closed sets $V\cap A_N(\eps)$ cover $V$, so at least one has nonempty interior in $V$; otherwise $V$ would be meager in itself. Since $V$ is open in $X$, that relative interior is an open subset of $X$ lying in $A_N(\eps)$, and meets $U_\eps$.

Let $x\in\bigcap_{k\geq1}U_{1/k}$. Given $\eta>0$, choose $k$ with $1/k<\eta/3$ and $N$ with $x\in\Int A_N(1/k)$. On a neighborhood $W$ of $x$ contained in that interior, shrink further so that $d(f_N(y),f_N(x))<\eta/3$. Passing to the limit in $m$ in the defining inequality for $A_N(1/k)$ gives $d(f_N(y),f(y))\leq1/k$ for $y\in W$, including $x$. The triangle inequality gives $d(f(y),f(x))<\eta$. Thus $f$ is continuous at $x$.

For any function into a metric space, its continuity set is $G_\delta$: let $W_k$ be the union of all open subsets $O$ of $X$ with $\diam f(O)<1/k$. A continuity point belongs to every $W_k$ by choosing a sufficiently small image ball, and membership in every $W_k$ implies continuity by comparing $f(y)$ with $f(x)$ in the same small-diameter image. Hence the continuity set is $\bigcap_kW_k$. We already showed that it contains the dense intersection of the open dense sets $U_{1/k}$. It is therefore dense and residual by \cref{d10:14:category-equivalences}.
\end{proof}

\begin{proposition}[A local uniform bound]\label{d10:14:local-uniform-boundedness}
If $X$ is a nonempty Baire space and $\mathcal A\subseteq C(X,\R)$ is pointwise bounded, there are a nonempty open $U\subseteq X$ and an integer $M\geq1$ with $|f(x)|\leq M$ for all $f\in\mathcal A,x\in U$.
\end{proposition}
\begin{proof}
For $M\geq1$ the set $F_M=\bigcap_{f\in\mathcal A}\{x:|f(x)|\leq M\}$ is closed, and pointwise boundedness gives $X=\bigcup_MF_M$. The Baire property and nonemptiness imply that some $F_M$ has nonempty interior. Take that interior as $U$. This conclusion is local; it does not assert a common bound on all of $X$.
\end{proof}

\section{Contraction iteration}

\begin{theorem}[Banach contraction principle with estimates]\label{d10:14:banach}
Let $(X,d)$ be a nonempty complete metric space and $T:X\to X$ satisfy
\[
 d(Tx,Ty)\leq qd(x,y)\qquad(x,y\in X),\qquad 0\leq q<1.
\]
There is a unique fixed point $p$. For every $x_0\in X$, define $x_{n+1}=Tx_n$ for $n\geq0$. Then $x_n\to p$, and
\begin{align*}
 d(x_n,p)&\leq\frac{q^n}{1-q}d(x_1,x_0) &&(n\geq0),\\
 d(x_n,p)&\leq\frac{1}{1-q}d(x_{n+1},x_n) &&(n\geq0),\\
 d(x_n,p)&\leq\frac{q}{1-q}d(x_n,x_{n-1}) &&(n\geq1).
\end{align*}
Here $q^0=1$, including $q=0$. More generally $d(z,p)\leq d(z,Tz)/(1-q)$ for every $z\in X$.
\end{theorem}
\begin{proof}
Induction gives $d(x_{n+1},x_n)\leq q^nd(x_1,x_0)$. For $m>n$, summing successive distances gives
\[
 d(x_m,x_n)\leq\sum_{k=n}^{m-1}q^kd(x_1,x_0)
 \leq\frac{q^n}{1-q}d(x_1,x_0).
\]
For $0<q<1$ the right side tends to zero, so the sequence is Cauchy; if $q=0$, it is constant from $x_1$ onward and is also Cauchy. Completeness gives a limit $p$. The contraction inequality makes $T$ continuous, so $Tx_n\to Tp$, whereas $Tx_n=x_{n+1}\to p$. Uniqueness of metric limits gives $Tp=p$. If $p'$ is another fixed point, $d(p,p')\leq qd(p,p')$, forcing equality of the points.

Letting $m\to\infty$ in the displayed sum gives the first estimate. Starting the sum at the current step instead gives $d(x_n,p)\leq d(x_{n+1},x_n)/(1-q)$. Apply $d(x_{n+1},x_n)\leq qd(x_n,x_{n-1})$ for the third. Finally
\[
 d(z,p)\leq d(z,Tz)+d(Tz,Tp)\leq d(z,Tz)+qd(z,p),
\]
and rearrange. None of this uses Baire category; completeness supports the direct iteration proof.
\end{proof}

\begin{corollary}[An iterate suffices]\label{d10:14:iterate-contraction}
If $X$ is nonempty complete metric and $T^m$ is a contraction for some integer $m\geq1$, then $T$ has a unique fixed point. No separate continuity assumption on $T$ is needed for this conclusion.
\end{corollary}
\begin{proof}
Let $p$ be the unique fixed point of $T^m$. Since $T^m(Tp)=T(T^mp)=Tp$, the point $Tp$ is another fixed point of $T^m$, so $Tp=p$. Every fixed point of $T$ is fixed by $T^m$ and must equal $p$.
\end{proof}

\begin{proposition}[Stability under perturbation]\label{d10:14:fixed-point-stability}
If contractions $S,T:X\to X$ on a nonempty complete metric space both have constant at most $q<1$ and $\sup_xd(Sx,Tx)\leq\delta$, their fixed points $p_S,p_T$ satisfy $d(p_S,p_T)\leq\delta/(1-q)$.
\end{proposition}
\begin{proof}
Using their fixed-point identities gives
\[
 d(p_S,p_T)=d(Sp_S,Tp_T)
 \leq d(Sp_S,Tp_S)+d(Tp_S,Tp_T)
 \leq\delta+qd(p_S,p_T).
\]
Rearrangement proves the estimate.
\end{proof}

\begin{example}[Why each hypothesis matters]\label{d10:14:contraction-boundaries}
On the incomplete space $(0,1)$, $T(x)=x/2$ is a contraction without a fixed point. On the complete space $\R$, translation $T(x)=x+1$ is nonexpansive, with constant $1$, and has no fixed point. The map $T(x)=x+1/x$ on the complete space $[1,\infty)$ strictly decreases the distance of every distinct pair but has no fixed point: for $x>y\geq1$,
\[
 T(x)-T(y)=(x-y)\left(1-\frac1{xy}\right),
\]
which lies strictly between $0$ and $x-y$. Thus a separate uniform constant below one is stronger than strict distance decrease. Finally the unique map on the empty complete metric space has no fixed point, explaining the explicit nonemptiness in the contraction principle.
\end{example}

% END SOURCE chapters/14.tex


% BEGIN SOURCE chapters/15.tex
\chapter{Local finiteness, paracompactness and partitions of unity}
\label{ch:d10-15}
\source{MUN sections 39--42; WIL section 20; KEL chapter 5, paracompactness and partition-of-unity problems.}

\section{Locally finite constructions}
\begin{definition}\label{d10:15:locally-finite}
An indexed family $(A_i)_{i\in I}$ is locally finite if every point has a
neighborhood meeting $A_i$ for only finitely many indices. A refinement of a
cover consists of sets each contained in some original member and still
covering the space. A family is $\sigma$-locally finite if it is a countable
union of locally finite families. A space is paracompact if every open cover
has a locally finite open refinement; Hausdorffness is a separate condition.
\end{definition}
Repeated nonempty sets count repeatedly in an indexed family. In particular,
local finiteness is stronger than requiring only finitely many members to
contain each point. Empty members cause no difficulty.

\begin{lemma}\label{d10:15:closure-sums}
If $(A_i)$ is locally finite, so is $(\cl A_i)$ and
\[
 \cl\Bigl(\bigcup_i A_i\Bigr)=\bigcup_i\cl A_i.
\]
Consequently a locally finite union of closed sets is closed. A family of
continuous real functions with locally finite supports has a continuous
pointwise sum, defined by the finitely many nonzero terms near each point.
\end{lemma}
\begin{proof}
An open set meeting $\cl A_i$ meets $A_i$. Thus the same open neighborhood
witnesses local finiteness for the closures. If $x$ belongs to none of the
closures, take a neighborhood meeting only $A_{i_1},\ldots,A_{i_k}$ and remove
their closures. The remaining neighborhood misses every $A_i$, so $x$ is
outside the closure of their union. The reverse inclusion always holds.
For functions, on a neighborhood meeting only finitely many supports, the
sum equals a fixed finite sum of continuous functions. Continuity is local.
\end{proof}

\begin{theorem}\label{d10:15:normal}
A paracompact Hausdorff space is $T_4$. A closed subspace of any
paracompact space is paracompact.
\end{theorem}
\begin{proof}
First separate a point $a$ from a closed set $B$ not containing it.
For each $b\in B$, Hausdorffness gives an open $V_b$ containing $b$ whose
closure excludes $a$: take disjoint neighborhoods of $a$ and $b$.
Refine $\{V_b:b\in B\}\cup\{X\setminus B\}$ by a locally finite open cover.
The subfamily meeting $B$ refines the $V_b$, covers $B$, and has union $V$
with $a\notin\cl V$ by \cref{d10:15:closure-sums}. This proves regularity.
For disjoint closed $A,B$, regularity gives each $b\in B$ a neighborhood
whose closure misses $A$. Repeat the same refinement argument to obtain
open $V\supset B$ with $\cl V\cap A=\varnothing$. Then $V$ and
$X\setminus\cl V$ separate $B$ and $A$. Hausdorffness implies $T_1$.

For a closed $Y\subset X$, write each relative member of a cover of $Y$
as $Y\cap U$ with $U$ open in $X$. Add $X\setminus Y$, refine in $X$,
and intersect the refinement with $Y$, discarding empty sets. This is the
required locally finite relative open refinement.
\end{proof}

\section{Why all metric spaces are paracompact}
\begin{lemma}[A metric refinement]\label{d10:15:metric-sigma}
Every open cover of a metric space has a $\sigma$-locally discrete open
refinement. Here locally discrete means that each point has a neighborhood
meeting at most one member of the indicated family.
\end{lemma}
\begin{proof}
The empty space is immediate. Well-order the open cover as $(U_\alpha)$;
this uses choice. For $n\geq1$ set
\[
 F_{n,\alpha}=\{x:B(x,1/n)\subset U_\alpha\}
                 \setminus\bigcup_{\beta<\alpha}U_\beta,\qquad
 V_{n,\alpha}=\bigcup_{x\in F_{n,\alpha}} B(x,1/(3n)).
\]
Each $V_{n,\alpha}$ is open and contained in $U_\alpha$. If $\alpha<\beta$,
$x\in F_{n,\alpha}$ and $y\in F_{n,\beta}$, then $y\notin U_\alpha$
whereas $B(x,1/n)\subset U_\alpha$, so $d(x,y)\geq1/n$.
The distance between a point of $V_{n,\alpha}$ and one of $V_{n,\beta}$
is therefore greater than $1/(3n)$. A ball of radius $1/(6n)$ meets
at most one such set. For coverage, given $x$, let $\alpha$ be the least
index with $x\in U_\alpha$. Choose $n$ with $B(x,1/n)\subset U_\alpha$.
Then $x\in F_{n,\alpha}\subset V_{n,\alpha}$. Taking all $n$ proves the claim.
\end{proof}

\begin{lemma}[Refinement upgrade]\label{d10:15:michael}
Suppose $X$ is regular and every open cover has a $\sigma$-locally finite
open refinement. Then $X$ is paracompact.
\end{lemma}
\begin{proof}
We establish three upgrades, each valid for every open cover.
First let $\bigcup_{n\geq1}\mathcal V_n$ be such an open refinement, and
write $W_n=\bigcup\mathcal V_n$. For $V\in\mathcal V_n$ put
\[
 S_{n,V}=V\setminus\bigcup_{i<n}W_i.
\]
These sets, although not necessarily open, cover $X$: use the least $n$
with $x\in W_n$. They are locally finite. Indeed, near $x$ choose
$N$ and $V_0\in\mathcal V_N$ containing $x$. The open set $V_0$
misses every $S_{n,V}$ with $n>N$. Intersect it with neighborhoods meeting
only finitely many members of each of $\mathcal V_1,\ldots,\mathcal V_N$.
The intersection meets only finitely many of the $S_{n,V}$.
Thus every open cover has a locally finite refinement by arbitrary sets.

Second, for a cover $\mathcal U$, regularity gives a covering by open
sets $G$ such that $\cl G$ is contained in some member of $\mathcal U$.
Apply the first upgrade, obtaining a locally finite arbitrary refinement
$(B_i)$ of these $G$. Then $(\cl B_i)$ is a locally finite closed
refinement of $\mathcal U$.

Third take a locally finite closed refinement $\mathcal B$ of $\mathcal U$.
The open sets meeting only finitely many members of $\mathcal B$ cover $X$.
By the second upgrade this cover has a locally finite closed refinement
$\mathcal C$. Every $C\in\mathcal C$ meets only finitely many $B\in\mathcal B$.
Define
\[
 E(B)=X\setminus\bigcup\{C\in\mathcal C:C\cap B=\varnothing\}.
\]
The union removed is closed, so $E(B)$ is open and contains $B$.
Choose $U(B)\in\mathcal U$ containing $B$ and replace $E(B)$ by
$E(B)\cap U(B)$. These sets cover and refine $\mathcal U$.
To check local finiteness, a neighborhood $W$ of $x$ meets only finitely
many $C_1,\ldots,C_k$. Because $\mathcal C$ covers $X$, $W$ is contained
in their union. If a $C_j$ meets $E(B)$, it cannot be disjoint from $B$,
by the definition of $E(B)$. Each $C_j$ meets only finitely many $B$.
It follows that $W$ meets only finitely many of the expanded sets, as required.
\end{proof}

\begin{theorem}\label{d10:15:metric-paracompact}
Every metric space is paracompact. Every regular Lindelof space is
paracompact. Compact spaces are paracompact without a separation hypothesis.
\end{theorem}
\begin{proof}
Metric spaces are regular and \cref{d10:15:metric-sigma,d10:15:michael}
apply. In a Lindelof space every open cover has a countable open
subcover; this is $\sigma$-locally finite, since each of its singleton
families is locally finite. Apply the upgrade. In a compact space a
finite subcover itself is a locally finite open refinement.
\end{proof}

\section{Shrinking and partitions of unity}
\begin{theorem}[Finite partitions on a normal space]\label{d10:15:finite-partition}
Every finite indexed open cover of a normal space admits a subordinate
finite partition of unity, without a $T_1$ or paracompactness assumption.
\end{theorem}
\begin{proof}
First obtain a finite shrinking. If $U_1,\ldots,U_k$ cover $X$, replace
$U_i$ in turn by an open $V_i$ as follows. At the $i$th step put
\[
 A_i=X\setminus(V_1\cup\cdots\cup V_{i-1}\cup U_{i+1}\cup\cdots\cup U_k).
\]
The previously maintained cover implies $A_i\subset U_i$.
It is closed. Normal separation of $A_i$ and $X\setminus U_i$ gives
$A_i\subset V_i\subset\cl V_i\subset U_i$. The cover is still a
cover after this replacement. Induction gives a finite open cover
$(V_i)$ with $\cl V_i\subset U_i$. Apply the same argument to $(V_i)$
to get a cover $(W_i)$ with $\cl W_i\subset V_i$.
By \cref{d10:07:urysohn} choose $h_i:X\to[0,1]$ equal to one on
$\cl W_i$ and zero on $X\setminus V_i$. Then
$\supp h_i\subset\cl V_i\subset U_i$, and $h=\sum_{i=1}^k h_i$
is continuous and positive since the $W_i$ cover.
The functions $h_i/h$ give the partition. If $k=0$, the cover
condition forces $X=\varnothing$, where the empty family suffices.
\end{proof}

\begin{lemma}[Indexed shrinking]\label{d10:15:shrinking}
Let $(U_i)_{i\in I}$ be an indexed open cover of a paracompact Hausdorff space. There is a
locally finite indexed open cover $(V_i)_{i\in I}$ with $\cl V_i\subset U_i$.
Some $V_i$ may be empty.
\end{lemma}
\begin{proof}
Regularity and paracompactness give a locally finite open cover
$(B_j)_{j\in J}$ such that for each $j$ there is $i(j)$ with
$\cl B_j\subset U_{i(j)}$. Set $V_i=\bigcup_{i(j)=i}B_j$.
Local finiteness gives
$\cl V_i=\bigcup_{i(j)=i}\cl B_j\subset U_i$.
A neighborhood meeting only $B_{j_1},\ldots,B_{j_k}$ meets $V_i$
only for $i\in\{i(j_1),\ldots,i(j_k)\}$. Thus the indexed family is locally finite.
\end{proof}

\begin{definition}\label{d10:15:partition-definition}
A partition of unity subordinate to $(U_i)$ consists of continuous
$\phi_i:X\to[0,1]$ with $\supp\phi_i\subset U_i$, locally finite indexed
supports, and $\sum_i\phi_i(x)=1$ at every point. The support is the closure
of the nonzero set, not just that set.
\end{definition}

\begin{theorem}\label{d10:15:partition}
Every indexed open cover of a paracompact Hausdorff space admits a
subordinate partition of unity. Conversely, if every indexed open cover
admits such a partition, the space is paracompact.
\end{theorem}
\begin{proof}
Shrink twice to get locally finite open covers $(V_i)$ and $(W_i)$ with
$\cl W_i\subset V_i$ and $\cl V_i\subset U_i$.
Normality and \cref{d10:07:urysohn} give a continuous $h_i:X\to[0,1]$
equal to $1$ on $\cl W_i$ and $0$ on $X\setminus V_i$.
Then $\supp h_i\subset\cl V_i\subset U_i$. The supports are locally finite.
Their sum $h$ is continuous by \cref{d10:15:closure-sums}, and $h(x)\geq1$
because the $W_i$ cover. Set $\phi_i=h_i/h$. This has all the asserted
properties; division by $h$ does not enlarge a support.
For the converse, the open sets $\{\phi_i>0\}$ cover, refine the original
cover, and are locally finite because the supports are. On the empty space
all conditions are vacuous.
\end{proof}

\begin{corollary}\label{d10:15:patching}
Suppose local continuous sections $s_i:U_i\to E$ of an affine equation
take values in a real topological vector space, and finite affine
combinations of solutions remain solutions. A subordinate partition of
unity patches them to a global continuous solution
$s(x)=\sum_i\phi_i(x)s_i(x)$.
\end{corollary}
\begin{proof}
Extend each $\phi_i s_i$ by zero outside $U_i$.
If $x\notin U_i$, the closed support of $\phi_i$ misses $x$, so some
neighborhood of $x$ has $\phi_i=0$. This proves continuity across the
boundary without asserting that $s_i$ itself extends. Near any point the
sum is finite and continuous. At that point the coefficients are
nonnegative and sum to one, so the affine equation is preserved.
\end{proof}

\section{General metrization}
\begin{corollary}[Embedding a compact finite atlas]\label{d10:15:compact-atlas-embedding}
Suppose a compact Hausdorff space $X$ has a finite open cover
$U_1,\ldots,U_k$ and homeomorphisms $\chi_i:U_i\to O_i\subset\R^n$.
Then $X$ embeds in $\R^{k(n+1)}$. In particular every compact
Hausdorff topological $n$-manifold has an embedding into some
finite-dimensional Euclidean space.
\end{corollary}
\begin{proof}
Compact Hausdorff spaces are normal. A finite subordinate partition
$(\phi_i)$ exists by \cref{d10:15:finite-partition}. Extend the
$n$ coordinate functions of $\phi_i\chi_i$ by zero outside $U_i$.
They are continuous there: any point outside $U_i$ is outside
the closed support of $\phi_i$, and so has a neighborhood on
which this product is identically zero. The map
\[
 F(x)=\bigl(\phi_i(x),\,\phi_i(x)\chi_i(x)\bigr)_{i=1}^k
\]
is therefore continuous. If $F(x)=F(y)$, choose $i$ with
$\phi_i(x)>0$. Then $\phi_i(y)$ has the same positive value,
both points belong to $U_i$, and division of the coordinate
block gives $\chi_i(x)=\chi_i(y)$. Thus $x=y$.
A continuous injection from compact to Hausdorff is an embedding.
A compact manifold has a finite chart subcover, giving the last
assertion. The empty atlas case forces $X=\varnothing$ and gives
its unique embedding into $\R^0$.
\end{proof}

\begin{lemma}\label{d10:15:perfect-normal}
A regular $T_1$ space with a $\sigma$-locally finite open base is normal,
and every closed set is a $G_\delta$. Each open set is a cozero set of
a continuous $[0,1]$-valued function.
\end{lemma}
\begin{proof}
Write the base as $\bigcup_n\mathcal B_n$, with each layer locally finite.
For closed $F$, let $O_n$ be the union of those $B\in\mathcal B_n$
with $\cl B\subset X\setminus F$. Regularity and the base property show
$X\setminus F=\bigcup_n O_n=\bigcup_n\cl O_n$, where the second equality
uses \cref{d10:15:closure-sums} to keep $\cl O_n$ disjoint from $F$.
Thus $F=\bigcap_n(X\setminus\cl O_n)$.

For disjoint closed $A,B$, construct open $G_n$ covering $A$, with
$\cl G_n\cap B=\varnothing$, by taking the $n$th-layer union of base
members whose closures miss $B$. Similarly construct $H_n$ covering $B$
whose closures miss $A$. The open sets
\[
 G=\bigcup_n\Bigl(G_n\setminus\bigcup_{j\leq n}\cl H_j\Bigr),\qquad
 H=\bigcup_n\Bigl(H_n\setminus\bigcup_{j\leq n}\cl G_j\Bigr)
\]
contain $A,B$ respectively and are disjoint: compare the two indices
of hypothetical memberships. Hence $X$ is normal.

If $F=\bigcap_n U_n$ is closed, Urysohn gives $f_n:X\to[0,1]$
equal to zero on $F$ and one on $X\setminus U_n$.
The uniformly convergent series $f=\sum_{n\geq1}2^{-n}f_n$ is continuous,
vanishes on $F$, and is positive at every point outside $F$.
Applying this to $F=X\setminus O$ proves the cozero assertion.
\end{proof}

\begin{theorem}[Nagata--Smirnov]\label{d10:15:nagata-smirnov}
A space is metrizable if and only if it is regular, $T_1$, and has a
$\sigma$-locally finite open base.
\end{theorem}
\begin{proof}
For a metric space, apply \cref{d10:15:metric-sigma} to the covers by
balls of radius $1/m$. The union over $m$ of the resulting countably
many locally finite layers is an open base: its elements have arbitrarily
small diameter at each point. It is again $\sigma$-locally finite.
Metric spaces are regular and $T_1$.

Conversely write the base as $\bigcup_{n\geq1}\mathcal B_n$.
By \cref{d10:15:perfect-normal}, for each $B\in\mathcal B_n$ there is
a continuous $f_{n,B}:X\to[0,2^{-n}]$ positive exactly on $B$.
Map $x$ to the vector $F(x)=(f_{n,B}(x))$ and put the supremum metric
on its bounded cube. This map is injective: if $x\ne y$, choose
a base neighborhood of $x$ excluding $y$.
It is continuous into the supremum metric. Given $\eps>0$, choose
$N$ with $2^{-N}<\eps/2$. Near $x$, each layer $n\leq N$ has only
finitely many supports of nonzero coordinate functions to consider,
because the family $(B)_{B\in\mathcal B_n}$ is locally finite.
Shrink a common neighborhood to make all these finitely many coordinates
vary by less than $\eps/2$; the other coordinates in these layers vanish
there. For $n>N$, their total coordinate variation is at most $2^{-n}$
and hence less than $\eps/2$. The supremum variation is at most $\eps/2$.

The inverse on $F(X)$ is continuous. If $x\in B\subset U$ with
$B\in\mathcal B_n$, then $f_{n,B}(x)>0$, and a sufficiently small
supremum ball about $F(x)$ has inverse image contained in
$\{f_{n,B}>0\}=B\subset U$. Pull back this metric along $F$.
It induces exactly the original topology.
\end{proof}

\begin{theorem}[Smirnov]\label{d10:15:smirnov}
A Hausdorff space is metrizable if and only if it is paracompact and locally
metrizable, where locally metrizable means that each point has an open
metrizable neighborhood.
\end{theorem}
\begin{proof}
The forward direction follows from \cref{d10:15:metric-paracompact}.
For the converse choose a cover by open metrizable $U_i$ and shrink it
to a locally finite cover $V_i$ with $\cl V_i\subset U_i$.
Each $U_i$ has a base $\bigcup_n\mathcal B_{i,n}$ with locally finite
layers, by the metric direction of \cref{d10:15:nagata-smirnov}.
For fixed $i,n$, the sets $B\cap V_i$, $B\in\mathcal B_{i,n}$,
are open in $X$ and locally finite in $X$: inside $U_i$ use its local
finiteness, and outside $\cl V_i$ use the disjoint neighborhood.
For fixed $n$ the union over $i$ is locally finite, since $(V_i)$ is
locally finite and only finitely many of these relative local checks are
needed near each point. The union over all $n,i$ is a base of $X$.
The space is regular and $T_1$ by \cref{d10:15:normal}, so
\cref{d10:15:nagata-smirnov} supplies a metric.
\end{proof}

None of these arguments assumes that an arbitrary product of paracompact
spaces is paracompact. The Sorgenfrey-plane example in
\cref{ch:d10-18} explains why that assertion would be false.

% END SOURCE chapters/15.tex


% BEGIN SOURCE chapters/16.tex
\chapter{Stone--Weierstrass approximation}
\label{ch:d10-16}
\source{WIL section 44, including the compact-open and locally compact approximation exercises; KEL chapter 7, approximation problems.}

\section{The lattice argument and its scalar ingredient}
For compact $X$ let $C(X,\R)$ denote the continuous real functions with
$\norm{f}_\infty=\sup_{x\in X}|f(x)|$; the norm on the unique function on
the empty space is zero. Compactness of continuous images makes this
supremum finite. A real subalgebra is closed under sums, products and real
scalar multiples. It is unital when it contains the constant function $1$.
For nonempty $X$ this includes every constant; for empty $X$ all functions
coincide.

\begin{lemma}[Dini's compactness argument]\label{d10:16:dini}
Suppose $X$ is compact and continuous $u_n:X\to\R$ increase pointwise to
a continuous $u$. Then $u_n\to u$ uniformly. The same assertion holds for
a decreasing sequence.
\end{lemma}
\begin{proof}
For $\eps>0$ the closed sets
$E_n=\{x:u(x)-u_n(x)\geq\eps\}$ decrease and have empty intersection.
If every $E_n$ were nonempty they would have the finite intersection
property and compactness would give a point in their intersection.
Thus some $E_n$ is empty, and then all subsequent errors are less than
$\eps$. Apply the result to $-u_n$ for the decreasing case.
\end{proof}

\begin{lemma}[Polynomial approximation of absolute value]\label{d10:16:absolute-polynomial}
There are real polynomials $p_n$ converging uniformly to $|t|$ on $[-1,1]$.
\end{lemma}
\begin{proof}
Set $p_0(t)=0$ and
\[
 p_{n+1}(t)=p_n(t)+\tfrac12\bigl(t^2-p_n(t)^2\bigr).
\]
For $s=|t|\leq1$, if $0\leq p_n(t)\leq s$, then
$0\leq\frac12(s-p_n(t))(s+p_n(t))\leq s-p_n(t)$.
Induction gives $0\leq p_n(t)\leq p_{n+1}(t)\leq s$.
Its scalar limit $\ell$ satisfies $\ell=\ell+(s^2-\ell^2)/2$,
so $\ell=s$. Apply \cref{d10:16:dini} on the compact real interval
to the continuous polynomials and their continuous limit.
This proves the only scalar approximation input we need; the general
polynomial approximation theorem is not being assumed.
\end{proof}

\begin{lemma}\label{d10:16:closed-lattice}
The uniform closure $\overline A$ of a unital real subalgebra
$A\subset C(X,\R)$ is a unital subalgebra and a lattice:
if $f,g\in\overline A$, then $|f|,\min(f,g),\max(f,g)\in\overline A$.
\end{lemma}
\begin{proof}
Uniform limits preserve continuity. For sums and scalar multiples take
the corresponding approximating sequences. For products, if
$a_n\to f$ and $b_n\to g$ uniformly, then $(\norm{a_n}_\infty)$ is
bounded and
\[
 \norm{a_nb_n-fg}_\infty
 \leq\norm{a_n}_\infty\norm{b_n-g}_\infty
           +\norm{g}_\infty\norm{a_n-f}_\infty\longrightarrow0.
\]
This proves the algebra assertion. If $\norm{f}_\infty\leq M$ with
$M>0$, the polynomials $M p_n(f/M)$ belong to $\overline A$ and
converge uniformly to $|f|$ by \cref{d10:16:absolute-polynomial}.
Closedness puts $|f|$ in $\overline A$. Finally use
\[
 \max(f,g)=\tfrac12(f+g+|f-g|),\qquad
 \min(f,g)=\tfrac12(f+g-|f-g|).
\]
\end{proof}

\begin{theorem}[Real Stone--Weierstrass]\label{d10:16:stone-weierstrass-real}
Let $X$ be compact Hausdorff. A unital real subalgebra $A\subset C(X,\R)$
that separates points is uniformly dense in $C(X,\R)$.
\end{theorem}
\begin{proof}
The empty case is immediate. Fix $f\in C(X,\R)$ and $\eps>0$.
For $x,y\in X$ choose $h_{x,y}\in A$ with
$h_{x,y}(x)=f(x)$ and $h_{x,y}(y)=f(y)$.
When $x=y$ use the constant $f(x)$. Otherwise take a separating $a\in A$
and its affine combination
\[
 h_{x,y}(z)=f(x)
       +\frac{f(y)-f(x)}{a(y)-a(x)}\bigl(a(z)-a(x)\bigr).
\]
For fixed $x$, the open sets
$\{z:h_{x,y}(z)>f(z)-\eps\}$, as $y$ varies, cover $X$.
Choose a finite subcover and let $g_x$ be the maximum of the corresponding
$h_{x,y}$. Then $g_x\in\overline A$, $g_x(x)=f(x)$, and
$g_x>f-\eps$ everywhere. The open sets
$\{z:g_x(z)<f(z)+\eps\}$ cover $X$ as $x$ varies.
Choose finitely many $x$ and take the minimum of their $g_x$.
The resulting $g\in\overline A$ satisfies
$f-\eps<g<f+\eps$ on $X$.
Thus the closed set $\overline A$ contains approximants arbitrarily
close to $f$ and must contain $f$. Equivalently, first approximate $g$
by a member of $A$ and let both errors tend to zero.
\end{proof}

\begin{corollary}\label{d10:16:polynomial}
On every compact $K\subset\R^n$, restrictions of real polynomials in
the coordinate functions are uniformly dense in $C(K,\R)$.
In particular this proves Weierstrass approximation on a compact interval.
\end{corollary}
\begin{proof}
Coordinates separate distinct points of $K$, and the polynomial algebra
contains constants. The compact subspace $K$ is Hausdorff. Apply the theorem.
For $n=0$, $K$ is empty or a singleton, so constants suffice.
\end{proof}

\section{Complex functions and the role of conjugation}
\begin{theorem}[Complex Stone--Weierstrass]\label{d10:16:stone-weierstrass-complex}
Let $X$ be compact Hausdorff. A unital complex subalgebra of $C(X,\C)$
that separates points and is closed under complex conjugation is uniformly dense.
\end{theorem}
\begin{proof}
Its real-valued members form a unital real subalgebra.
For $a\in A$, both $\operatorname{Re}a=(a+\bar a)/2$ and
$\operatorname{Im}a=(a-\bar a)/(2i)$ belong to this real algebra.
If $a$ distinguishes two points, one of these real functions does.
The real theorem approximates the real and imaginary parts of any
continuous complex function. Add the real approximant to $i$ times
the imaginary approximant; the complex uniform error is at most the
sum of their real uniform errors.
\end{proof}

\begin{example}\label{d10:16:conjugation-needed}
On the unit circle, ordinary complex polynomials in $z$ alone separate
points and contain constants, but are not uniformly dense.
\end{example}
\begin{proof}
The bounded linear functional
\[
 L(f)=\frac1{2\pi}\int_0^{2\pi}e^{it}f(e^{it})\,dt
\]
has $|L(f)|\leq\norm{f}_\infty$. For a polynomial in $z$, every
integrand term is $e^{ikt}$ with $k\geq1$ and integrates to zero by
one-variable calculus. But $L(\bar z)=1$. A uniform approximation
of $\bar z$ by these polynomials would contradict continuity of $L$.
Conjugation is therefore a real condition, not a cosmetic change of notation.
\end{proof}

\section{Vanishing at infinity}
\begin{definition}\label{d10:16:c0}
For locally compact Hausdorff $X$, $C_0(X,\R)$ consists of continuous
$f:X\to\R$ such that $\{x:|f(x)|\geq\eps\}$ is compact for every $\eps>0$.
Define $C_0(X,\C)$ by the same condition using the complex absolute value.
\end{definition}

\begin{lemma}\label{d10:16:c0-extension}
Let $X^*=X\cup\{\infty\}$ have the one-point compactification topology
of \cref{d10:11:one-point}; allow $\infty$ to be isolated if $X$ is
already compact. A function $f\in C(X,\R)$ belongs to $C_0(X,\R)$
if and only if its extension $f^*(\infty)=0$ is continuous on $X^*$.
Such functions are bounded, form a real algebra and are closed under uniform limits.
\end{lemma}
\begin{proof}
If the level sets are compact, the complement in $X^*$ of such a set
is a neighborhood of infinity on which $|f^*|<\eps$.
Conversely, continuity at infinity gives a compact $K\subset X$
outside which $|f|<\eps$. The closed level set
$\{|f|\geq\eps\}$ is a closed subset of $K$, hence compact.
The extension is continuous on the compact space $X^*$, so is bounded.
Sums, scalar multiples and products of the extensions remain continuous
and vanish at infinity, proving the algebra assertion.
A uniform limit of their zero-at-infinity extensions is again
continuous and zero there. This proves closedness.
\end{proof}

\begin{theorem}\label{d10:16:stone-weierstrass-c0}
Let $X$ be locally compact Hausdorff. A real subalgebra
$A\subset C_0(X,\R)$ that separates points and vanishes nowhere
(for each $x$ some $a\in A$ has $a(x)\ne0$) is uniformly dense in $C_0(X,\R)$.
The complex version holds with conjugation invariance as an additional hypothesis.
\end{theorem}
\begin{proof}
Extend every $a\in A$ by zero at infinity.
The unitized algebra $\{a^*+c:a\in A,c\in\R\}$ separates points
of $X^*$: the given separation works within $X$, and the nowhere-vanishing
condition separates each $x$ from infinity.
Apply \cref{d10:16:stone-weierstrass-real}.
If $\norm{f^*-(a^*+c)}_\infty<\eps$, evaluation at infinity gives
$|c|<\eps$, and hence $\norm{f-a}_\infty<2\eps$.
This proves density. The same reasoning with the complex theorem
handles conjugation-invariant complex algebras.
The compact case and the empty case follow from exactly the same
construction, including the isolated extra point.
\end{proof}

This is an approximation statement about a specified function algebra and
topology. It is not an assertion that arbitrary point-separating families
are uniformly dense, or that pointwise density suffices for uniform density.

% END SOURCE chapters/16.tex


% BEGIN SOURCE chapters/17.tex
\chapter{Finite-dimensional topological fixed points}
\label{ch:d10-17}
\source{MUN section 55 for the fixed-point/no-retraction connection; WIL section 34 proves the disk case using fundamental groups. The independent combinatorial route here proves all finite dimensions without assuming fundamental-group or homology results.}

The contraction theorem has a quantitative hypothesis and gives uniqueness.
Brouwer's theorem assumes only continuity on a nonempty compact convex
finite-dimensional set and gives existence. Neither theorem replaces the other.
We supply the triangulation and combinatorial steps so that no unproved
no-retraction or degree theorem is hidden in the argument.

\section{Simplexes and arbitrarily fine subdivisions}
\begin{definition}\label{d10:17:simplex}
The standard $n$-simplex is
\[
 \Delta^n=\{(t_0,\ldots,t_n)\in\R^{n+1}:t_i\geq0,\
                                    \textstyle\sum_i t_i=1\}.
\]
An affine simplex is the convex hull of affinely independent vertices.
Its faces are convex hulls of nonempty subsets of its vertices, with the
empty face allowed when discussing intersections. A finite geometric
triangulation is a finite collection of such simplexes covering the
given simplex, closed under faces, with two cells intersecting in a common
face or not at all. Its maximal cells have dimension $n$.
Its mesh is the largest diameter of its cells.
\end{definition}

\begin{lemma}[Barycentric subdivision]\label{d10:17:subdivision}
Every affine $n$-simplex has finite triangulations with mesh tending to zero
and compatible induced triangulations on all faces.
\end{lemma}
\begin{proof}
For each nonempty face $F$ let $b(F)$ be the mean of its vertices.
Take as new cells the convex hulls of the barycentres in strictly increasing
chains of faces, together with their faces. The vertices in a maximal chain
are affinely independent: in the order in which original vertices enter the
chain, the barycentric-coordinate matrix is triangular with nonzero diagonal.
Shorter chains are subfamilies of these independent vertex sets.
To see coverage and compatibility,
write a point in barycentric coordinates $t_i$ and order them
$t_{\sigma(0)}\geq\cdots\geq t_{\sigma(n)}\geq0$. Set
$t_{\sigma(n+1)}=0$. Then
\[
 x=\sum_{k=0}^n
 (k+1)\bigl(t_{\sigma(k)}-t_{\sigma(k+1)}\bigr)
             b(\{v_{\sigma(0)},\ldots,v_{\sigma(k)}\}).
\]
The coefficients are nonnegative and their sum is
$\sum_i t_i=1$. A coefficient is positive exactly at a strict drop between
successive ordered coordinates. The prefix set above each strict drop is
intrinsic to the point, independent of how tied coordinates are ordered.
Thus in two ordering cells containing the same point its positive
barycentre coefficients occur at common chain vertices. The intersection
is exactly the convex hull of those vertices common to the two chains,
which is a face of each. Equalities and zero coordinates merely set
some coefficients to zero.
Every face of the original simplex is treated by the same construction
on its vertex set. This proves that the cells form a finite triangulation,
not merely a covering by small convex sets. Applying it cell by cell to
a finite triangulation agrees on shared faces.

If $F\subset G$ are nonempty faces with $|F|=a$, $|G|=b$, then
\[
 b(G)=\frac ab b(F)+\frac{b-a}{b}b(G\setminus F)
\]
when $a<b$, so
$\norm{b(G)-b(F)}\leq(1-a/b)D\leq nD/(n+1)$,
where $D$ is the original diameter. A convex hull has diameter no
larger than the largest distance between its vertices: express the
difference of two convex combinations as a convex combination of
vertex differences and use the triangle inequality.
Consequently one subdivision reduces every cell diameter by at least
the factor $n/(n+1)$ when $n\geq1$. Iterate. For $n=0$ the single
point already has mesh zero.
\end{proof}

\begin{lemma}[Facet incidence]\label{d10:17:incidence}
In a finite triangulation of $\Delta^n$, $n\geq1$, an
$(n-1)$-dimensional face of the triangulation belongs to exactly two
$n$-cells if it is internal, and exactly one if it lies on the boundary.
The boundary faces induce triangulations of the original facets.
\end{lemma}
\begin{proof}
All balls and half-balls in this proof are relative to
$\operatorname{aff}(\Delta^n)$. Choose a relative interior point $p$ of
such a face $F$. A cell not
containing $p$ misses a sufficiently small ball about $p$, since the
complex is finite and its cells are closed. A cell containing $p$
intersects $F$ in a face containing a relative interior point of $F$,
so it contains $F$. For an $n$-cell containing $F$, near $p$ all its
other facet inequalities are strict. It therefore occupies precisely
one closed half-ball determined by the affine hyperplane of $F$.
Two distinct cells cannot occupy the same half-ball, since their
interiors would overlap, contradicting the common-face condition.
If $p$ is internal to the original simplex, both half-balls must be
covered, giving two cells; on its boundary exactly the inward half-ball
must be covered, giving one. A boundary $(n-1)$-cell lies in one
original facet, since an intersection of two original facets has
dimension at most $n-2$. Intersecting a cell with an original facet
is a face: the facet's affine coordinate is nonnegative on every
cell vertex, and its zero set in the cell is the convex hull of exactly
the vertices with zero coordinate. Restricting the finite common-face
covering to an original facet therefore proves the last assertion.
\end{proof}

\section{The parity proof of Sperner's lemma}
\begin{theorem}[Sperner]\label{d10:17:sperner}
Label every vertex of a finite triangulation of $\Delta^n$ by one of
$0,\ldots,n$, requiring that a vertex in an original face receive
one of that face's vertex labels. Then the number of $n$-cells
whose labels are all $0,\ldots,n$ is odd, in particular nonzero.
\end{theorem}
\begin{proof}
For $n=0$ there is one vertex and the assertion is immediate.
Induct on $n$. Call a triangulation facet special if its $n$ vertex
labels are exactly $0,\ldots,n-1$. Count incidences of special
facets with $n$-cells, modulo two. Internal special facets contribute
two and boundary special facets contribute one, by
\cref{d10:17:incidence}.
A special boundary facet cannot lie in the original face opposite
vertex $j<n$, since the face condition forbids label $j$ there.
Thus all special boundary facets lie in the face opposite $n$.
On that face the restricted labeling satisfies the induction
hypothesis. The number of its fully labeled $(n-1)$-cells is odd.
The incidence count is therefore odd.

An $n$-cell with all $n+1$ labels has exactly one special facet:
remove its vertex labeled $n$. A cell without label $n$ but with
all labels $0,\ldots,n-1$ has one repeated label and exactly two
special facets, obtained by deleting either copy. Every other
cell has zero special facets. Hence the same incidence parity
equals the number of fully labeled $n$-cells. That number is odd.
\end{proof}

\begin{theorem}[Brouwer on a simplex]\label{d10:17:brouwer-simplex}
Every continuous $f:\Delta^n\to\Delta^n$ has a fixed point.
\end{theorem}
\begin{proof}
For $n=0$ there is only one point. Take triangulations with mesh
tending to zero from \cref{d10:17:subdivision}. If any vertex is
already fixed, there is nothing to prove. Otherwise label a vertex
$v$ by an index $i$ for which $v_i>f_i(v)$. Such an index exists
because the coordinates of $v-f(v)$ sum to zero and are not all zero.
It satisfies $v_i>0$, since $f_i(v)\geq0$, so the label is
permitted on the face containing $v$.

For each triangulation choose a fully labeled cell by Sperner's lemma
and write $v_{r,i}$ for its vertex carrying label $i$.
The simplex is a closed bounded subset of Euclidean space and is
compact by \cref{d10:05:heine-borel}. By metric sequential
compactness, \cref{d10:09:compact-equivalences}, a subsequence of
$v_{r,0}$ tends to some $x\in\Delta^n$.
The mesh bound shows that along this same subsequence every
$v_{r,i}$ tends to $x$. Passing to the limit in
$(v_{r,i})_i>f_i(v_{r,i})$ gives $x_i\geq f_i(x)$ for all $i$.
Since both coordinate sums are one, every inequality is equality.
Thus $x=f(x)$. This is a metric compactness argument; no false
sequence criterion in an arbitrary topological space is being used.
\end{proof}

\section{General compact convex sets}
\begin{lemma}[Metric projection]\label{d10:17:projection}
For a nonempty compact convex $K\subset\R^m$, every $x\in\R^m$
has a unique nearest point $P(x)\in K$. This projection is a retraction
and satisfies $\norm{P(x)-P(y)}\leq\norm{x-y}$.
\end{lemma}
\begin{proof}
The continuous squared distance to $x$ attains its minimum on $K$.
If two different minimizers $p,q$ existed, convexity and
\[
 \norm{x-\tfrac12(p+q)}^2
 =\tfrac12\norm{x-p}^2+\tfrac12\norm{x-q}^2-\tfrac14\norm{p-q}^2
\]
would give a smaller value at their midpoint. Thus the minimizer is unique.
For $z\in K$ and $0<t\leq1$, minimality at $p=P(x)$ gives
\[
 0\leq \norm{x-p-t(z-p)}^2-\norm{x-p}^2
       =-2t\ip{x-p}{z-p}+t^2\norm{z-p}^2 .
\]
Divide by $t$ and let $t$ decrease to zero to obtain
$\ip{x-p}{z-p}\leq0$. Put $q=P(y)$ and use these inequalities with
$z=q$ and $z=p$ respectively. They give
\[
 \norm{p-q}^2\leq\ip{x-y}{p-q}
                 \leq\norm{x-y}\norm{p-q}.
\]
Divide when $p\ne q$, and treat $p=q$ directly. This proves the
Lipschitz estimate and continuity. If $x\in K$, its distance to
itself is zero, so $P(x)=x$.
\end{proof}

\begin{theorem}[Finite-dimensional Brouwer]\label{d10:17:brouwer-convex}
Let $K$ be a nonempty compact convex subset of $\R^m$.
Every continuous self-map $f:K\to K$ has a fixed point.
The same holds in any finite-dimensional real normed vector space.
\end{theorem}
\begin{proof}
For $m=0$ a nonempty subset has one point. Otherwise choose $M>0$
with $K\subset[-M,M]^m$. The affine simplex
\[
 S=\{x:x_i\geq-M\text{ for all }i,\ \textstyle\sum_i x_i\leq mM\}
\]
contains this cube. Its vertices are $v_0=(-M,\ldots,-M)$ and
$v_0+2mM e_i$, $1\leq i\leq m$; they are affinely independent.
An affine coordinate homeomorphism identifies $S$ with $\Delta^m$.
The continuous map $f\circ P:S\to K\subset S$, using
\cref{d10:17:projection}, has a fixed point $x$ by the simplex
theorem. That equality puts $x$ in $K$, so $P(x)=x$ and $f(x)=x$.
No nonempty-interior or full-affine-dimension assumption on $K$ was used.

For a real normed space with basis $e_1,\ldots,e_m$, the coordinate
map $T:\R^m\to V$ satisfies
$\norm{Tt}\leq(\sum_i\norm{e_i}^2)^{1/2}\norm{t}_2$.
The norm $\norm{Tt}$ is therefore continuous in $t$.
On the Euclidean unit sphere it has a strictly positive minimum $c$,
by compactness and injectivity; consequently
$c\norm{t}_2\leq\norm{Tt}$ for every $t$ by homogeneity.
These two inequalities make $T$ a homeomorphism with continuous inverse.
Conjugate $f$ and transport $K$ through this map. The zero-dimensional
case was already handled.
\end{proof}

\begin{corollary}[No retraction]\label{d10:17:no-retraction}
For $n\geq1$ there is no continuous retraction from the closed unit ball
$B^n\subset\R^n$ onto its boundary sphere $S^{n-1}$.
\end{corollary}
\begin{proof}
If $r$ were such a retraction, $x\mapsto-r(x)$ would be a continuous
self-map of the compact convex ball and would have a fixed point $x$.
Then $x=-r(x)$ lies on the sphere, whence $r(x)=x$.
It follows that $x=-x$ and $x=0$, contradicting $\norm{x}=1$.
\end{proof}

\begin{remark}\label{d10:17:fixed-point-boundaries}
Nonemptiness is essential: the empty self-map has no fixed point.
Convexity is essential even in the plane: a nontrivial rotation of the
unit circle has none. Compactness cannot be replaced just by boundedness:
$f(x)=(x+1)/2$ maps $(0,1)$ into itself without a fixed point.
The identity map on a nontrivial compact convex set has many fixed points,
so Brouwer alone asserts neither uniqueness nor a contraction estimate.
\end{remark}

% END SOURCE chapters/17.tex


% BEGIN SOURCE chapters/18.tex
\chapter{Canonical compact models and a counterexample atlas}
\label{ch:d10-18}
\source{WIL sections 28--30 for compact zero-dimensional models; STE examples 4,11,19--20,25,42--43,64 and the reference charts; MUN sections 13,23--25,29--32. The doubled-origin line is defined separately below and is not identified with the different planar construction in STE example 74.}

\section{The Cantor space as a reusable model}
\begin{theorem}\label{d10:18:cantor-model}
The space $C=\{0,1\}^{\N_{\geq1}}$ with its product topology is a nonempty
compact metrizable space with a clopen base and no isolated points.
It is homeomorphic to the middle-thirds Cantor subset of $[0,1]$ by
\[
 (a_n)\longmapsto \sum_{n\geq1}\frac{2a_n}{3^n}.
\]
\end{theorem}
\begin{proof}
The finite discrete factor is compact. Product compactness is
\cref{d10:05:tychonoff}, and the product is metrizable by
\cref{d10:10:product-metric}. Cylinder sets imposing finitely many
coordinates are clopen and form a base. No cylinder is a singleton,
since an unspecified later coordinate may be changed.

The displayed series converges uniformly in $(a_n)$ by its geometric
tail bound, and each partial sum is continuous. If two sequences first
differ at $k$, the contribution $2/3^k$ there is larger than the maximum
tail difference $\sum_{n>k}2/3^n=1/3^k$. Thus the map is injective.
At level $r$ of the middle-thirds construction, the retained intervals
are exactly
\[
 \left[\sum_{n=1}^r\frac{2a_n}{3^n},
       \sum_{n=1}^r\frac{2a_n}{3^n}+\frac1{3^r}\right].
\]
The image lies in all these finite unions. Conversely a point in their
intersection selects a unique nested retained interval at each level;
the intervals are separated by positive gaps and have diameter $3^{-r}$.
Its selected digits give precisely the displayed sum. This proves
surjectivity onto the Cantor set. A continuous bijection from compact
to Hausdorff is a homeomorphism by \cref{d10:05:compact-bijection}.
\end{proof}

\begin{lemma}\label{d10:18:clopen-partitions}
A compact zero-dimensional metric space has finite clopen partitions
of arbitrarily small mesh, and such a partition may refine any given
finite clopen partition. If it has no isolated points, every nonempty
clopen part can be split into any prescribed positive finite number of nonempty
clopen parts.
\end{lemma}
\begin{proof}
Choose for each point a clopen neighborhood contained in a ball of
radius $\eps/3$, and use compactness to select finitely many.
The nonempty atoms of the Boolean algebra generated by these clopen
sets form a finite clopen partition. Each atom is contained in one
cover member and has diameter at most $2\eps/3<\eps$. Intersect with the
parts of a prescribed partition to obtain a refinement.

For the second assertion, a nonempty clopen part cannot be a singleton,
since that would be an isolated point. Choose two points and use a
clopen neighborhood of one missing the other. Intersect it with the
part and take its relative complement to split the part into two
nonempty clopen sets. Each resulting set again has no isolated points,
because it is open in the original space. Repeat one split at a time.
\end{proof}

\begin{theorem}[Cantor characterization]\label{d10:18:cantor-characterization}
Every nonempty compact metrizable zero-dimensional space with no
isolated points is homeomorphic to $C$.
Equivalently in this statement, zero-dimensional may be replaced by
totally disconnected.
\end{theorem}
\begin{proof}
We prove that any two such spaces $X,Y$ are homeomorphic. Begin by
matching their single whole-space parts. At odd stage $r$, refine
each currently matched clopen part of $X$ into finitely many nonempty
clopen parts of diameter less than $2^{-r}$, using
\cref{d10:18:clopen-partitions}. Refine its partner in $Y$ into the
same number of nonempty clopen parts and match them arbitrarily.
At even stages do the small-mesh refinement in $Y$ and match by
splitting the corresponding parts in $X$. If useful one can require
at least two children by one extra split; the lemma permits this.
Partitions remain finite, matched and nested, and both meshes tend
to zero, because later refinement never increases a diameter.

A point $x$ selects a nested sequence of parts of $X$. Their partners
are nonempty nested closed subsets of compact $Y$, so intersect,
and their diameters tend to zero, so have exactly one point $h(x)$.
Reversing the construction gives an inverse: a point in either
space belongs to exactly one part of each partition.
For continuity, given a neighborhood of $h(x)$, choose a sufficiently
small partner part containing it and contained in that neighborhood.
The corresponding clopen part of $X$ is a neighborhood of $x$
mapped into it. Thus $h$ and its inverse are continuous.
Apply this to $Y=C$ from \cref{d10:18:cantor-model}.
Finally compact Hausdorff totally disconnected spaces have clopen
bases by \cref{d10:06:zero-dimensional}; a clopen base conversely
separates distinct points into disjoint clopen parts, so forces
components to be singletons.
\end{proof}

\begin{theorem}[Cantor surjection]\label{d10:18:cantor-surjection}
Every nonempty compact metric space is a continuous image of $C$.
\end{theorem}
\begin{proof}
Let $K_\varnothing=K$. Cover each already constructed nonempty compact
$K_s$ by finitely many relative closed balls of radius $2^{-r-2}$
with centers in $K_s$, where its children are to be at level $r+1$.
Use compactness first for the corresponding open balls; the closed
balls still cover. Intersect each with $K_s$. The children are
nonempty compact sets, cover their parent, and have diameter at
most $2^{-r-1}$. There are finitely many nodes at each level.
Let $m_{r+1}\geq1$ be the largest number of children of a node at
that level. Duplicate a child where necessary so every node has
exactly $m_{r+1}$ indexed children.

A branch in $\prod_{r\geq1}\{1,\ldots,m_r\}$ gives nested nonempty
compact sets with diameters tending to zero. Their intersection is
a unique point. This branch map is continuous: branches with a
common prefix of length $r$ give points at distance at most $2^{-r}$.
It is onto, since for any $x\in K$ one recursively chooses the
least child containing $x$.
For each $r$ choose a positive integer $\ell_r$ with
$2^{\ell_r}\geq m_r$ and a surjection from binary blocks of length
$\ell_r$ onto $\{1,\ldots,m_r\}$. Divide a binary sequence into these
successive finite blocks. The resulting map from $C$ to the branch
space is continuous and onto; each prescribed child sequence has
the concatenation of its least preimage blocks as a preimage.
Compose the two maps. Nonemptiness is necessary, since the
nonempty space $C$ cannot map onto the empty set.
\end{proof}

\begin{corollary}[A space-filling curve]\label{d10:18:peano}
There is a continuous surjection $[0,1]\to[0,1]^2$.
More generally, for every finite integer $m\geq1$, the cube $[0,1]^m$
is a continuous image of the interval.
\end{corollary}
\begin{proof}
By \cref{d10:18:cantor-model,d10:18:cantor-surjection}, the closed
middle-thirds Cantor subset $K\subset[0,1]$ admits a continuous
surjection onto the cube. Each coordinate is a continuous
$[0,1]$-valued function on $K$. The bounded Tietze theorem
\cref{d10:07:tietze-bounded}, followed if necessary by the continuous
clipping map $t\mapsto\min(1,\max(0,t))$, extends it to a
$[0,1]$-valued continuous function on the whole interval.
The coordinate extensions give a continuous map into the cube;
it is onto because its restriction to $K$ already is.
This proves the existence statement of MUN section 44 without
requiring its particular iterative curve construction.
\end{proof}

\section{Separation, countability and metric hypotheses}
\begin{example}[The Sierpinski space]\label{d10:18:sierpinski}
On $\{0,1\}$ take $\{\varnothing,\{1\},\{0,1\}\}$ as topology.
The space is $T_0$, compact and connected, but not $T_1$ or Hausdorff.
\end{example}
\begin{proof}
The topology axioms are immediate. The open singleton distinguishes
the two points, while $\{1\}$ is not closed. Every open cover of a
finite set has a finite subcover. The only nonempty proper open
set has nonopen complement, so there is no separation.
The constant sequence $1,1,\ldots$ converges to both points:
every neighborhood of either contains $1$.
\end{proof}

\begin{example}[Bounded, complete and not compact]\label{d10:18:discrete}
An infinite set with distance $d(x,y)=1$ for $x\ne y$ is bounded
and complete but not compact or totally bounded. An uncountable
such space is also locally compact and first countable, but not
separable, second countable or Lindelof.
\end{example}
\begin{proof}
A Cauchy sequence is eventually constant, using tolerance $1/2$.
Balls of radius less than one are singletons, so a finite ball cover
cannot cover an infinite set. The singleton open cover also disproves
compactness; for an uncountable set it disproves the Lindelof property.
Each singleton is a compact open neighborhood and a local base.
Every dense subset must contain each isolated point, so separability
fails in the uncountable case. A base must likewise contain every
singleton, so cannot be countable.
\end{proof}

\begin{example}[The $K$-topology]\label{d10:18:k-topology}
Let $K=\{1/n:n\geq1\}$. On $\R$ use the basis consisting of usual
open intervals and their differences with $K$.
This is a Hausdorff second-countable space that is not regular,
and hence is not metrizable.
\end{example}
\begin{proof}
Intersections of basis members have the required local basis
refinements. The topology refines the usual Hausdorff topology,
and rational-endpoint intervals of the two indicated forms give
a countable base. The set $K$ is closed, since its complement is
open, and it does not contain zero.
Every neighborhood $U$ of zero contains
$(-\eps,\eps)\setminus K$ for some $\eps>0$.
Choose $1/n\in(0,\eps)$. A basis neighborhood of $1/n$ must contain
an ordinary interval about it, since a deleted-$K$ interval does
not contain that point. Every such interval meets
$(-\eps,\eps)\setminus K$. Thus $1/n\in\cl U$.
No neighborhood of zero has closure disjoint from $K$, so
regularity fails. This explains why the regularity hypothesis in
Urysohn metrization cannot be reduced just to Hausdorffness.
\end{proof}

\begin{example}[A line with two origins]\label{d10:18:two-origins}
Take two disjoint copies of $\R$ and identify their points with
equal nonzero coordinate, keeping two different origins.
Give the result the quotient topology. It is $T_1$, second countable,
locally metrizable and paracompact, but not Hausdorff or metrizable.
\end{example}
\begin{proof}
Each copy maps homeomorphically onto the open subset missing the
other origin. Indeed its open intervals have open saturated inverse
images: the same interval in one copy and the interval with zero
removed in the other. These two open charts prove local
metrizability and second countability. Singleton inverse images
are closed in the disjoint sum, so points of the quotient are closed.
Neighborhoods of the two origins contain punctured intervals about
zero and therefore intersect; Hausdorffness fails.

For paracompactness, start with an open cover and choose one member
about each origin. Their union contains every nonzero point of
sufficiently small absolute value. Choose $\delta>0$ so each origin
has its own chart interval $(-\delta,\delta)$ inside its selected
member. On the ordinary line form an open cover consisting of
$(-\delta,\delta)$ and the coordinate portions of the original
cover away from $[-\delta/2,\delta/2]$. These portions are ordinary
open, since the quotient away from zero is the ordinary punctured
line. Refine this cover locally finitely by the metric
paracompactness theorem. Retain the refining members assigned to
the outer portions; discard those assigned to the central interval.
The retained members cover all points with $|x|\geq\delta$, lie
away from zero, and are locally finite also in the doubled line.
Add the two chosen chart intervals; they cover the remaining
points and add only two sets. This is a locally finite open
refinement of the original cover.
It shows explicitly the need for the Hausdorff condition in the
Smirnov theorem, not merely a failure of global second countability.
\end{proof}

\begin{example}[Fortissimo]\label{d10:18:fortissimo}
Let $D$ be uncountable and $X=D\cup\{p\}$. Points of $D$ are isolated;
a set containing $p$ is open when its complement is a countable
subset of $D$. This space is normal, Hausdorff, zero-dimensional
and Lindelof, but is not compact, first countable or separable.
\end{example}
\begin{proof}
Topology axioms follow because finite unions of countable sets
are countable. Separate two isolated points by their singletons,
and separate $p$ from $d$ by $X\setminus\{d\}$ and $\{d\}$.
If a closed set omits $p$, it is countable and also open.
For disjoint closed sets at least one omits $p$; that open set
and its open complement separate the pair. Thus the space is normal.
Singletons in $D$ and neighborhoods of $p$ are clopen, giving a
clopen base.

Any cover member containing $p$ leaves at most countably many
points, which can be covered by countably many more members.
For a countably infinite $A\subset D$, the cover
$\{X\setminus A\}\cup\{\{a\}:a\in A\}$ has no finite subcover.
If $(U_n)$ were a countable neighborhood base at $p$, choose
distinct $d_n\in U_n\setminus\{p\}$; each $U_n$ contains
uncountably many choices. The neighborhood $X\setminus\{d_n:n\geq1\}$
contains none of the $U_n$, a contradiction.
Finally every dense set contains every isolated point of $D$.
\end{proof}

\section{Ordinal compactness and convergence}
\begin{lemma}\label{d10:18:ordinal-compact}
For every ordinal $\beta$, the interval $[0,\beta]$ is compact
Hausdorff in its order topology.
\end{lemma}
\begin{proof}
Ordered spaces are Hausdorff by \cref{d10:01:ordered-spaces}.
For an open cover, suppose some initial interval has no finite
subcover and let $\alpha\leq\beta$ be the least such ordinal.
It is not zero. If it is a successor, add a cover member containing
its last point to a finite cover of the preceding interval.
If it is a limit, a cover member containing $\alpha$ contains a
tail $(\gamma,\alpha]$. The interval $[0,\gamma]$ already has a
finite subcover by minimality, giving the same contradiction.
For a point $\alpha<\beta$, intersect its cover neighborhood with
the subspace $[0,\alpha]$ in this argument.
Thus no such $\alpha$ exists.
\end{proof}

\begin{example}\label{d10:18:omega-one}
Let $\omega_1$ be the first uncountable ordinal, and let
$Y=[0,\omega_1)$ and $Z=[0,\omega_1]$, in their order topologies.
$Y$ is first countable, locally compact, countably compact and
sequentially compact, but not compact or Lindelof and not
$\sigma$-compact. $Z$ is compact Hausdorff but not first countable;
no sequence of points of $Y$ converges to $\omega_1$, although
$\omega_1\in\cl_Z Y$.
\end{example}
\begin{proof}
Every countable limit ordinal has a countable cofinal sequence:
enumerate its predecessors $\xi_1,\xi_2,\ldots$ and put
$\gamma_n=\max\{\xi_1,\ldots,\xi_n\}+1<\alpha$.
The intervals $(\gamma_n,\alpha]$ give a countable local
base at a countable limit $\alpha$; zero and successor points
are isolated. A compact interval $[0,\alpha]$ is a neighborhood
of its last point in $Y$, so $Y$ is locally compact.

Suppose a countable open cover $(U_n)$ of $Y$ has no finite
subcover. Choose $x_n\notin U_1\cup\cdots\cup U_n$.
Their supremum is a countable ordinal $\beta<\omega_1$.
Compactness of $[0,\beta]$ gives a finite subcover from the $U_n$,
contradicting the choice for all sufficiently large $n$.
This proves countable compactness; first countability then gives
sequential compactness by \cref{d10:05:countable-compactness}.
The cover by initial open intervals $[0,\alpha)$ has no finite
subcover and no countable subcover, because a countable set of
countable ordinals has countable supremum. Hence $Y$ is neither
compact nor Lindelof.
Every compact subset of $Y$ is bounded by some countable ordinal,
using that same cover and a finite subcover. Countably many such
subsets are still bounded, so cannot cover $Y$.

Compactness of $Z$ is \cref{d10:18:ordinal-compact}.
For a countable list of neighborhoods $V_n$ of $\omega_1$, choose
tails $(\alpha_n,\omega_1]\subset V_n$. Choose a countable ordinal
$\gamma$ strictly larger than every $\alpha_n$.
The smaller neighborhood $(\gamma,\omega_1]$ contains none of the
$V_n$, since every $V_n$ contains $\alpha_n+1\leq\gamma$.
Thus no countable list is a local base.
Any sequence in $Y$ is bounded by a countable ordinal and misses
a tail neighborhood of $\omega_1$. Every such neighborhood
nevertheless meets $Y$, proving the closure assertion.
\end{proof}

\section{An atlas that reuses proofs}
\begin{example}\label{d10:18:rational-incomplete}
The rationals in their usual metric are not complete.
\end{example}
\begin{proof}
Choose rational $q_n$ with $|q_n-\sqrt2|<2^{-n}$, using the density
of $\Q$ in $\R$. This sequence is Cauchy by the triangle inequality.
If it had a rational limit in the same metric, it would have two
different limits in $\R$, contradicting uniqueness of metric limits.
Its real limit $\sqrt2$ is not rational.
\end{proof}

The following entries point to the definitions and complete proofs already
given; they are not duplicate requirements to reprove the same mathematics.
\begin{longtable}{p{.25\textwidth}p{.42\textwidth}p{.23\textwidth}}
\toprule
Construction & Distinction witnessed & Proof interface\\
\midrule
\endhead
Topologist's sine closure &
Connected does not imply path connected; the exact added vertical set matters. &
\cref{d10:06:sine-curve}\\
Sorgenfrey line &
Separable and first countable need not imply second countable. &
\cref{d10:08:sorgenfrey-line}\\
Sorgenfrey plane &
Products of normal or paracompact spaces need not be normal or paracompact. &
\cref{d10:08:sorgenfrey-plane,d10:15:normal,d10:15:metric-paracompact}\\
Cocountable topology &
Closure and continuity need not be detected by sequences. &
\cref{d10:08:cocountable,d10:04:closure-net}\\
Compact binary cube &
A compact Hausdorff space need not be sequentially compact. &
\cref{d10:05:variant-counterexamples}\\
Rationals with their metric &
Local compactness and completeness do not follow from metrizability. &
\cref{d10:11:rationals,d10:18:rational-incomplete}\\
Uniformity of $\arctan$ &
The same topology can carry a complete and an incomplete compatible metric. &
\cref{d10:12:continuity-calculus}\\
Non-LCH Ascoli domain &
Compact-open compactness need not force equicontinuity on arbitrary domains. &
\cref{d10:13:ascoli-nonlch}\\
Quotient times $\Q$ &
A quotient map need not remain quotient after an arbitrary product. &
\cref{d10:03:quotient-product-failure}\\
\bottomrule
\end{longtable}

For the Sorgenfrey-plane row, each line is regular Lindelof and hence
paracompact by \cref{d10:15:metric-paracompact}; their product is Hausdorff
but not normal by the cited construction. By \cref{d10:15:normal}
that product cannot be paracompact either. No closure property of
arbitrary products is inferred.

\begin{proposition}[Evaluation can fail]\label{d10:18:evaluation-failure}
The map $C(\Q,\R)_{\mathrm{co}}\times\Q\to\R$, $(f,x)\mapsto f(x)$,
is not continuous at $(0,0)$, where $\mathrm{co}$ denotes compact-open.
\end{proposition}
\begin{proof}
A basic neighborhood of zero controls its values on finitely many
compact subsets of $\Q$; let $K$ be their union. Such a $K$ has empty
relative interior in $\Q$: if it contained a rational interval, its
image in $\R$ would be compact and therefore closed, forcing it
to contain an irrational limit of that interval.
Given a neighborhood $U$ of zero in $\Q$, choose $y\in U\setminus K$.
If $K$ is nonempty, the continuous function
\[
 h(z)=\frac{2\dist(z,K)}{\dist(z,K)+|z-y|}
\]
is zero on $K$ and equals $2$ at $y$; the denominator never vanishes
since $y\notin K$ and $K$ is closed. If $K$ is empty use $h=2$.
Thus $h$ belongs to the original basic neighborhood, while its
evaluation at $y$ lies outside $(-1,1)$. No product neighborhood
of $(0,0)$ is mapped into that interval.
\end{proof}

The reference charts in STE are useful for discovering examples, but their
entries must be translated into the conventions fixed here and proved for
the actual underlying set and topology. A name alone, or a checkmark in a
chart, is not an equivalence proof and is not an extra axiom.

% END SOURCE chapters/18.tex

\backmatter

% BEGIN SOURCE source-register.tex
\chapter{Source register and scope decisions}
\label{app:d10-sources}
The exact supplied PDFs, their SHA-256 values, physical page counts and
provenance are recorded in the adjacent textbook register and
\texttt{knowledge/SOURCES.json}. The source-to-result and chapter-disposition
files accompany this manuscript. Printed textbook pages and physical PDF
pages are different coordinates; Willard's reflowed electronic edition in
particular must not be cited by an invented constant offset.

\begin{description}[style=nextline,leftmargin=12mm]
\item[MUN] James R. Munkres, \emph{Topology}, second edition, Prentice Hall,
2000. The general-topology part supplies the core sequence. Sections
39--42 are expanded in \cref{ch:d10-15}; the existence of a space-filling
curve in section 44 follows by a different proof in
\cref{d10:18:peano}. Compact finite-atlas embedding is included in
\cref{d10:15:compact-atlas-embedding}. Part II's fundamental groups,
coverings, surface classification and group applications are routed to D23.
\item[WIL] Stephen Willard, \emph{General Topology}, original 1970 work,
unabridged Dover reprint, 2004. Convergence, uniformities, function spaces,
compactification and approximation supply complementary interfaces.
Its Cantor representation results are included in
\cref{ch:d10-18}. Its section 34 disk fixed-point result uses fundamental
groups; \cref{ch:d10-17} instead proves all finite dimensions
combinatorially. Proximity theory and its specialized compactification
correspondence are not admitted foundations.
\item[KEL] John L. Kelley, \emph{General Topology}, D. Van Nostrand,
copyright 1955, supplied as a later scanned printing. Direct page-image
checks compare subnet conventions, the chain pseudometric construction,
uniform completion, compact-open topology and Ascoli.
The image scan has no searchable text layer; a complete-book OCR or
page-by-page proof audit is not claimed. The actual checked page pairs
are listed in \texttt{audit/kelley-source-check.json}.
\item[STE] Lynn Arthur Steen and J. Arthur Seebach, Jr.,
\emph{Counterexamples in Topology}, second edition, Springer-Verlag, 1978.
The examples and charts guide distinctions of hypotheses. Each adopted
example is defined and justified in the notes, possibly through an earlier
proved interface. Not all 143 examples are adopted. In particular the
doubled-origin line in \cref{d10:18:two-origins} is not silently identified
with the different construction carrying a similar name in STE.
\end{description}

\section*{What is outside the present volume}
General dimension theory, Menger--Nobeling bounds, classification of curves
and manifolds, fundamental-group and covering-space theory, and the full
Hahn--Mazurkiewicz characterization belong to later geometry/topology
directions. Special normal-product criteria, proximity/Samuel theory,
continuous-convergence structures and topological Ascoli for general
nonmetric targets are likewise not asserted as proved results of this volume.
The exact per-source dispositions identify these choices; no omitted
advanced theorem may be used as an undeclared premise of an adopted result.

General uniform pseudometric representation and separated uniform completion
are \emph{included} in \cref{ch:d10-12}. Their earlier tentative routing during
drafting was superseded by complete constructions and cross-review.
The scalar nowhere-differentiable-function topic from MUN section 49 is
routed to analysis rather than introduced as a topology axiom.

\section*{How different presentations are merged}
One canonical definition is fixed for each object. Source results with
different conventions are compared by hypotheses and mathematical
conclusions, not only by the same theorem name. A metric completion is
related to its Cauchy-filter uniform completion; a compact Hausdorff
embedding is related to its cube realization; a compact-open family is
related to its uniform convergence structure only in the proved setting.
No theorem is generalized merely to inflate scope, and no source route
is treated as the only permissible proof.

% END SOURCE source-register.tex


% BEGIN SOURCE verification-register.tex
\chapter{Proof checking and formalization boundary}
\label{app:d10-verification}
The accompanying author and cross-review files record separate mathematical
readings, concrete findings and repairs. The review checked actual
statements and arguments: missing hypotheses, empty sets, non-Hausdorff
cases, metric versus topological completeness, sequence versus filter
tests, extension ranges and compactification universality.

Examples of repaired draft issues include the missing fixed connected set
in a union assertion, an omitted open-cover condition in a shrinking
statement, and ambiguous finite-cube language in a space-filling
corollary. Source-note repairs distinguish the doubled-origin quotient
line from a different named planar example, and the fundamental-group
disk proof from the all-dimensional Sperner argument.

The mathematical review is distinct from structural validation.
The latter checks stable labels and references, proof-environment coverage,
declared dependency links, source identities and a successful LaTeX build.
PDF rendering is inspected for clipping, missing mathematical glyphs,
unreadable tables and unresolved references. Neither this structural
validation nor the mathematical cross-review is a Lean kernel certificate.

A future Lean project should take the exact LaTeX source as its mathematical
input, together with the explicit foundations and scope. Main statements,
definitions and constructions, substantive adopted examples, conditions,
quantifiers and quantitative conclusions must be preserved. Equivalent
valid arguments may replace these proofs. Necessary representation
bridges must be established, but a source's proof-route-only step does
not become a new independent acceptance condition automatically.

Before formalization, the fixed Mathlib/MathGenome environment and reusable
interfaces should be inspected. This manuscript does not create a new
formalization run, change an agent or overwrite the existing library.
The exact current proof-audit and build status is in the external
\texttt{VERIFICATION\_REPORT.txt}, bound to the final source inventory.

% END SOURCE verification-register.tex

\end{document}
